% Generated by scripts/build_docs.py from the wiki. Edit the wiki or this template.
% Compile with Tectonic, LuaLaTeX, or XeLaTeX. No external images are required.
\documentclass[10pt,a4paper,oneside]{report}
\usepackage{fontspec}
\setmainfont{DejaVu Sans}
\setsansfont{DejaVu Sans}
\setmonofont[Scale=0.85]{DejaVu Sans Mono}
\usepackage[a4paper,left=20mm,right=19mm,top=25mm,bottom=23mm,
  headheight=28pt,headsep=7mm,footskip=12mm]{geometry}
\usepackage[table]{xcolor}
\usepackage{microtype}
\usepackage{amsmath,amssymb}
\usepackage{array,booktabs,longtable,calc}
\usepackage{graphicx}
\usepackage{fvextra}
\usepackage{enumitem}
\usepackage{titlesec}
\usepackage{fancyhdr,lastpage}
\usepackage{etoolbox}
\usepackage{xurl}
\usepackage{hyperref}
\usepackage{bookmark}
\definecolor{MGSRed}{HTML}{B71930}
\definecolor{MGSInk}{HTML}{1C2630}
\definecolor{MGSMuted}{HTML}{52616D}
\definecolor{MGSRule}{HTML}{C9D0D5}
\definecolor{MGSGray}{HTML}{F0F3F5}
\color{MGSInk}
\newcommand{\ProductName}{Magnecko Sense}
\newcommand{\BoardVersion}{MGS-1.0-}
\newcommand{\DocumentID}{MGS-UG-001}
\newcommand{\DocumentRevision}{A}
\newcommand{\DocumentDate}{October 2026}
\hypersetup{colorlinks=true,linkcolor=MGSRed,urlcolor=MGSRed,
  pdftitle={Magnecko Sense -- User Guide and Reference Manual},
  pdfauthor={Magnecko},bookmarksnumbered=true,bookmarksopen=false}
\urlstyle{tt}
\setlength{\parindent}{0pt}
\setlength{\parskip}{5pt}
\setlength{\emergencystretch}{2em}
\setlength{\tabcolsep}{4pt}
\renewcommand{\arraystretch}{1.18}
\setlength{\LTpre}{8pt}
\setlength{\LTpost}{8pt}
\setcounter{tocdepth}{1}
\setcounter{secnumdepth}{2}
\makeatletter
\renewcommand*\l@section{\@dottedtocline{1}{1.5em}{3.6em}}
\patchcmd{\l@chapter}{1.5em}{2.4em}{}{\PackageError{mgsmanual}{Could not size chapter numbers}{Check the table of contents layout.}}
\makeatother
\setlist{nosep,leftmargin=1.5em,topsep=4pt,parsep=2pt}
\setlist[itemize]{label=\textcolor{MGSRed}{\textbullet}}
\providecommand{\tightlist}{\setlength{\itemsep}{1pt}\setlength{\parskip}{1pt}}
\providecommand{\passthrough}[1]{#1}
\providecommand{\pandocbounded}[1]{#1}
\DefineVerbatimEnvironment{verbatim}{Verbatim}{breaklines=true,breakanywhere=true,
  fontsize=\small,baselinestretch=1,frame=single,rulecolor=\color{MGSRule},
  framesep=6pt,formatcom=\color{MGSInk}}
\AtBeginEnvironment{longtable}{\small}
\titleformat{\chapter}[display]{\bfseries\color{MGSRed}}
  {\large\chaptertitlename\ \thechapter}{5pt}{\LARGE}
\titlespacing*{\chapter}{0pt}{0pt}{17pt}
\titleformat{\section}{\large\bfseries\color{MGSRed}}{\thesection}{0.6em}{}
\titleformat{\subsection}{\normalsize\bfseries}{\thesubsection}{0.6em}{}
\titlespacing*{\section}{0pt}{12pt}{5pt}
\titlespacing*{\subsection}{0pt}{9pt}{4pt}
\renewcommand{\chaptermark}[1]{\markboth{\thechapter\quad #1}{}}
\pagestyle{fancy}
\fancyhf{}
\fancyhead[L]{\textcolor{MGSRed}{\large\bfseries MGS}\quad\small\ProductName}
\fancyhead[R]{\scriptsize\DocumentID\ Rev.\ \DocumentRevision\\\BoardVersion}
\fancyfoot[L]{\scriptsize\color{MGSMuted}\nouppercase{\leftmark}}
\fancyfoot[R]{\small\bfseries\thepage}
\renewcommand{\headrulewidth}{0.5pt}
\renewcommand{\footrulewidth}{0.3pt}
\fancypagestyle{plain}{\fancyhf{}
  \fancyhead[L]{\textcolor{MGSRed}{\large\bfseries MGS}\quad\small\ProductName}
  \fancyhead[R]{\scriptsize\DocumentID\ Rev.\ \DocumentRevision\\\BoardVersion}
  \fancyfoot[L]{\scriptsize\color{MGSMuted}\DocumentDate\quad |\quad PRELIMINARY}
  \fancyfoot[R]{\small\bfseries\thepage}}
\begin{document}
\raggedright
\hypersetup{pageanchor=false}
\begin{titlepage}
\thispagestyle{empty}
{\color{MGSRed}\Large\bfseries magnecko}\par
\vspace{22mm}
{\color{MGSRed}\small\bfseries DEPTH + INERTIAL SENSING}\par
\vspace{4mm}
{\fontsize{32}{38}\selectfont\bfseries\ProductName}\par
\vspace{6mm}
{\fontsize{22}{29}\selectfont User Guide and\\Reference Manual\par}
\vspace{10mm}
{\large Development, integration, diagnostics, and validation}\par
\vspace{9mm}
\textcolor{MGSRed}{\rule{\linewidth}{1.2pt}}\par
\vspace{8mm}
\begin{tabular}{@{}ll}
Document & \DocumentID\quad Revision \DocumentRevision \\
Hardware baseline & \BoardVersion \\
Publication & \DocumentDate \\
Status & Preliminary engineering reference
\end{tabular}
\vfill
{\small This manual is generated from the maintained project wiki. It covers
build and programming workflows, boot diagnostics, device configuration, host
protocols, desktop tools, reconstruction, and the board acceptance test plan.
The historical memory-layout proposal is retained separately as an appendix.\par}
\vspace{6mm}
{\small\color{MGSMuted}Firmware repository: \texttt{magnecko-sense-fw}\\
Source paths are relative to the repository root. Chapter links are internal;
web references remain clickable. Test-plan checkboxes do not record completed tests.\par}
\end{titlepage}
\hypersetup{pageanchor=true}
\pagenumbering{roman}
\tableofcontents
\clearpage
\pagenumbering{arabic}
\chapter{Hardware identification and revisions}\label{readme--hardware-revision-standard}

{\small\color{MGSMuted}Source: \nolinkurl{README.md}\par}\medskip

Board identifiers use \textbf{\texttt{MGS-\allowbreak{}<MAJOR>.\allowbreak{}<MINOR><REVISION>}}. Revision 0 is always written as \texttt{-\allowbreak{}}, followed by \texttt{A}, \texttt{B}, and so on. Hardware identifiers are separate from firmware software versions.

{\def\LTcaptype{none} % do not increment counter
\begin{longtable}[]{@{}
  >{\raggedright\arraybackslash}p{(\linewidth - 4\tabcolsep) * \real{0.2234}}
  >{\raggedright\arraybackslash}p{(\linewidth - 4\tabcolsep) * \real{0.3883}}
  >{\raggedright\arraybackslash}p{(\linewidth - 4\tabcolsep) * \real{0.3883}}@{}}
\toprule\noalign{}
\begin{minipage}[b]{\linewidth}\raggedright
\textbf{Change}
\end{minipage} & \begin{minipage}[b]{\linewidth}\raggedright
\textbf{When to use it}
\end{minipage} & \begin{minipage}[b]{\linewidth}\raggedright
\textbf{Examples}
\end{minipage} \\
\midrule\noalign{}
\endhead
\bottomrule\noalign{}
\endlastfoot
Revision: \texttt{MGS-\allowbreak{}1.\allowbreak{}0-\allowbreak{}} → \texttt{MGS-\allowbreak{}1.\allowbreak{}0A} & Fix or refine the existing design without intentionally adding functionality or breaking compatibility. The same firmware should ideally work across these revisions; document any required workaround. & Correct traces or resistor values, improve signal integrity or decoupling, adjust placement, or fix power behavior. \\
Minor: \texttt{MGS-\allowbreak{}1.\allowbreak{}0B} → \texttt{MGS-\allowbreak{}1.\allowbreak{}1-\allowbreak{}} & Add capabilities while preserving existing external interfaces and essential functionality. Existing firmware should generally retain its existing functions; using new hardware may require updated firmware. & Add a temperature sensor, synchronization input, trigger output, compatible expansion connector, or compatible flash capacity increase. \\
Major: \texttt{MGS-\allowbreak{}1.\allowbreak{}1A} → \texttt{MGS-\allowbreak{}2.\allowbreak{}0-\allowbreak{}} & Break an established interface or change the architecture incompatibly. Existing firmware and connected hardware cannot be assumed compatible. & Change connector pin assignments or supply requirements, replace the MCU family, or adopt a substantially different sensor interface. \\
\end{longtable}
}

Classify changes by their purpose and compatibility impact, rather than their size or component count. For example, adding a companion IC to fix the already specified USB PD dead-battery behavior, with the same firmware and external behavior, is a revision change: \textbf{\texttt{MGS-\allowbreak{}1.\allowbreak{}0-\allowbreak{}} → \texttt{MGS-\allowbreak{}1.\allowbreak{}0A}}. It restores intended functionality. A single incompatible connector pin change can require a major version increment.

Apply these release rules:

\begin{enumerate}
\def\labelenumi{\arabic{enumi}.}
\tightlist
\item
  Give every released PCB layout modification a new identifier. Unreleased development edits do not each require a revision increment.
\item
  Reset the revision to \texttt{-\allowbreak{}} whenever the major or minor number changes: \texttt{MGS-\allowbreak{}1.\allowbreak{}0C} → \texttt{MGS-\allowbreak{}1.\allowbreak{}1-\allowbreak{}}, or \texttt{MGS-\allowbreak{}1.\allowbreak{}1B} → \texttt{MGS-\allowbreak{}2.\allowbreak{}0-\allowbreak{}}.
\item
  Never reuse an assigned identifier, even for an abandoned design.
\item
  Track assembly variants and BOM revisions separately from PCB versions. Different flash populations or electrically equivalent component substitutions on the same layout do not by themselves create a new PCB revision. If a variant requires a changed layout, assign the appropriate new board identifier.
\item
  Maintain a hardware compatibility table recording revision-specific quirks, available capabilities, assembly variants, and firmware requirements. Revisions can differ in actual behavior even when a change is classified as a bug fix. Detect optional features using capabilities or the assembly variant, rather than assuming a minor version alone guarantees their presence.
\end{enumerate}

Keep the firmware identity in {\texttt{common/\allowbreak{}board/\allowbreak{}board\_\allowbreak{}identity.\allowbreak{}h} (\nolinkurl{common/board/board_identity.h})}, the CubeMX USB product string, and the stored hardware fields consistent when releasing a board. See \hyperref[wiki-firmware-parameters--editable-fields]{device parameters} for the identity fields and how to update an existing device.

\chapter{Build, Flash and Debug Workflow}\label{wiki-getting-started-build-flash-debug--build-flash-and-debug-workflow}

{\small\color{MGSMuted}Source: \nolinkurl{wiki/getting-started/build-flash-debug.md}\par}\medskip

\section{Initial setup}\label{wiki-getting-started-build-flash-debug--1-initial-setup}

After cloning the repository, make the shell scripts executable:

\begin{verbatim}
chmod +x scripts/*.sh
\end{verbatim}

You normally only need to do this once.

The scripts automatically try to locate the STM32Cube GNU toolchain and STM32CubeProgrammer tools from common installation locations, including the STM32Cube VS Code bundles under:

\begin{verbatim}
~/.local/share/stm32cube/bundles/
\end{verbatim}

The build scripts prefer the STM32Cube GNU toolchain over Ubuntu\textquotesingle s generic \texttt{/\allowbreak{}usr/\allowbreak{}bin/\allowbreak{}arm-\allowbreak{}none-\allowbreak{}eabi-\allowbreak{}gcc}.

\subsection{Required tools}\label{wiki-getting-started-build-flash-debug--required-tools}

The following tools are required:

\begin{verbatim}
cmake
arm-none-eabi-gcc
arm-none-eabi-g++
arm-none-eabi-objcopy
STM32_SigningTool_CLI
STM32_Programmer_CLI
ST-LINK_gdbserver
\end{verbatim}

\texttt{STM32\_\allowbreak{}SigningTool\_\allowbreak{}CLI} may require installation of the full STM32CubeProgrammer package. The lightweight programmer bundle installed by the VS Code extension may contain \texttt{STM32\_\allowbreak{}Programmer\_\allowbreak{}CLI} without the signing tool.

\section{Building}\label{wiki-getting-started-build-flash-debug--2-building}

The main build script is:

\begin{verbatim}
scripts/build.sh
\end{verbatim}

Syntax:

\begin{verbatim}
./scripts/build.sh --config <Debug|Release> --target <appli|fsbl|extmem|all>
\end{verbatim}

Examples:

\begin{verbatim}
./scripts/build.sh --config Debug --target appli
./scripts/build.sh --config Debug --target fsbl
./scripts/build.sh --config Debug --target extmem
./scripts/build.sh --config Debug --target all
./scripts/build.sh --config Release --target all
\end{verbatim}

\subsection{Build outputs}\label{wiki-getting-started-build-flash-debug--build-outputs}

\subsubsection{Application}\label{wiki-getting-started-build-flash-debug--application}

Building \texttt{appli} produces:

\begin{verbatim}
Appli/build/magnecko-sense-fw_Appli.elf
Appli/build/magnecko-sense-fw_Appli.bin
Appli/build/magnecko-sense-fw_Appli-signed.bin
\end{verbatim}

The signed application image is the file used for programming external flash.

\subsubsection{FSBL}\label{wiki-getting-started-build-flash-debug--fsbl}

Building \texttt{fsbl} produces:

\begin{verbatim}
FSBL/build/magnecko-sense-fw_FSBL.elf
FSBL/build/magnecko-sense-fw_FSBL.bin
FSBL/build/magnecko-sense-fw_FSBL-trusted.bin
\end{verbatim}

The trusted FSBL image is the file used for programming external flash.

\subsubsection{External memory loader}\label{wiki-getting-started-build-flash-debug--external-memory-loader}

Building \texttt{extmem} produces:

\begin{verbatim}
ExtMemLoader/build/magnecko-sense-fw_ExtMemLoader.elf
ExtMemLoader/build/magnecko-sense-fw_ExtMemLoader.stldr
\end{verbatim}

The \texttt{.\allowbreak{}stldr} file is still an ELF file internally. It is copied from the generated ELF and must not be converted to a raw binary.

\section{Building from VS Code}\label{wiki-getting-started-build-flash-debug--3-building-from-vs-code}

Use:

\begin{verbatim}
Terminal -> Run Build Task
\end{verbatim}

or:

\begin{verbatim}
Ctrl+Shift+B
\end{verbatim}

Select:

\begin{verbatim}
Build configuration:
    Debug
    Release
\end{verbatim}

and then:

\begin{verbatim}
Build target:
    appli
    fsbl
    extmem
    all
\end{verbatim}

For normal development, \texttt{Debug} is usually the appropriate configuration.

For a complete firmware package, use:

\begin{verbatim}
Release + all
\end{verbatim}

or:

\begin{verbatim}
Debug + all
\end{verbatim}

depending on the desired build configuration.

\section{External flash layout}\label{wiki-getting-started-build-flash-debug--4-external-flash-layout}

The current external flash programming addresses are:

\begin{verbatim}
FSBL:
    0x90040000

Application Slot A:
    0x90100000

Application Slot B:
    0x90900000
\end{verbatim}

The custom ExtMemLoader exposes the SEMPER flash through the STM32N6 XSPI1 memory-mapped region.

The current boot flow uses:

\begin{verbatim}
BootROM
  -> FSBL2 at 0x90040000
  -> Application Slot A at 0x90100000
\end{verbatim}

\section{Flashing persistent external memory}\label{wiki-getting-started-build-flash-debug--5-flashing-persistent-external-memory}

Persistent programming is handled by:

\begin{verbatim}
scripts/flash.sh
\end{verbatim}

\subsection{Build and flash}\label{wiki-getting-started-build-flash-debug--build-and-flash}

To rebuild the required images before programming:

\begin{verbatim}
./scripts/flash.sh --target all --config Debug --build
\end{verbatim}

Other examples:

\begin{verbatim}
./scripts/flash.sh --target fsbl --config Debug --build
./scripts/flash.sh --target appli-a --config Debug --build
./scripts/flash.sh --target appli-b --config Debug --build
\end{verbatim}

Available flash targets are:

\begin{verbatim}
fsbl
appli-a
appli-b
all
\end{verbatim}

\texttt{all} currently programs:

\begin{verbatim}
FSBL
Application Slot A
\end{verbatim}

\subsection{Flash without rebuilding}\label{wiki-getting-started-build-flash-debug--flash-without-rebuilding}

If the images have already been built:

\begin{verbatim}
./scripts/flash.sh --target all
\end{verbatim}

or:

\begin{verbatim}
./scripts/flash.sh --target fsbl
./scripts/flash.sh --target appli-a
./scripts/flash.sh --target appli-b
\end{verbatim}

The script automatically uses:

\begin{verbatim}
ExtMemLoader/build/magnecko-sense-fw_ExtMemLoader.stldr
\end{verbatim}

with \texttt{STM32\_\allowbreak{}Programmer\_\allowbreak{}CLI}.

Every persistent flash also updates device/build parameters at \texttt{0x90000000}. The script reads the connected board\textquotesingle s existing parameters, preserves editable settings and other slots\textquotesingle{} records, validates the selected image manifests, then programs/verifies firmware followed by parameters. Build records now include compiler/toolchain versions, actual compile flags/definitions and feature options. Existing images keep their recorded Debug/Release configuration. Old manifests require a one-time rebuild; SRAM \texttt{-\allowbreak{}-\allowbreak{}dev} images must be rebuilt for persistent use.

Use \texttt{-\allowbreak{}-\allowbreak{}target parameters} to replace editable settings from the default JSON, or \texttt{-\allowbreak{}-\allowbreak{}parameters /\allowbreak{}path/\allowbreak{}to/\allowbreak{}device.\allowbreak{}json} to apply a specific settings file with any target. Inspect \texttt{build/\allowbreak{}flash\_\allowbreak{}parameters.\allowbreak{}json} or the PC viewer\textquotesingle s \textbf{Device \& firmware parameters} tab. \texttt{parameter\_\allowbreak{}build.\allowbreak{}built\_\allowbreak{}at\_\allowbreak{}utc} refreshes on every parameter build. Flashing FSBL, Appli A, Appli B or all also refreshes each selected slot\textquotesingle s \texttt{build.\allowbreak{}*.\allowbreak{}built\_\allowbreak{}at\_\allowbreak{}utc} and \texttt{programmed\_\allowbreak{}at\_\allowbreak{}utc}; \texttt{image\_\allowbreak{}built\_\allowbreak{}at\_\allowbreak{}utc} retains the image\textquotesingle s original build date. A parameter-only update preserves all firmware records. See \hyperref[wiki-firmware-parameters--device-parameters-and-build-metadata]{parameter management} for complete field and preservation details.

\section{Flashing from VS Code}\label{wiki-getting-started-build-flash-debug--6-flashing-from-vs-code}

Use:

\begin{verbatim}
Terminal -> Run Task
\end{verbatim}

and select:

\begin{verbatim}
Flash: Build + Program
\end{verbatim}

This asks for:

\begin{verbatim}
Debug / Release
\end{verbatim}

and:

\begin{verbatim}
all
fsbl
appli-a
appli-b
\end{verbatim}

Alternatively, use:

\begin{verbatim}
Flash: Program Existing Images
\end{verbatim}

to program already-generated images without rebuilding.

\section{SRAM development mode}\label{wiki-getting-started-build-flash-debug--7-sram-development-mode}

The STM32N6 can run firmware directly from internal SRAM without programming the external flash.

For this mode:

\begin{verbatim}
BOOT1 = 1
\end{verbatim}

The firmware is volatile and is lost after reset or power removal.

The development helper is:

\begin{verbatim}
scripts/run_dev.sh
\end{verbatim}

Supported targets:

\begin{verbatim}
appli
fsbl
\end{verbatim}

\subsection{Build, load into SRAM, run, then detach}\label{wiki-getting-started-build-flash-debug--build-load-into-sram-run-then-detach}

Example:

\begin{verbatim}
./scripts/run_dev.sh     --action run     --target appli     --config Debug     --build
\end{verbatim}

For the FSBL:

\begin{verbatim}
./scripts/run_dev.sh     --action run     --target fsbl     --config Debug     --build
\end{verbatim}

This performs approximately:

\begin{verbatim}
build
-> reset target
-> load ELF into SRAM
-> start execution
-> detach debugger
\end{verbatim}

The MCU continues running after the debugger disconnects.

\section{SRAM development from VS Code}\label{wiki-getting-started-build-flash-debug--8-sram-development-from-vs-code}

Use:

\begin{verbatim}
Terminal -> Run Task
-> Dev: Load + Run SRAM (detach)
\end{verbatim}

Then select:

\begin{verbatim}
Debug / Release
\end{verbatim}

and:

\begin{verbatim}
appli / fsbl
\end{verbatim}

This is useful when the firmware should simply run from SRAM without keeping a debugger attached.

\section{Source-level debugging from SRAM}\label{wiki-getting-started-build-flash-debug--9-source-level-debugging-from-sram}

For interactive source-level debugging, use the VS Code Run and Debug panel.

Available configurations:

\begin{verbatim}
Dev: Debug Application from SRAM
Dev: Debug FSBL from SRAM
\end{verbatim}

These configurations:

\begin{verbatim}
build the selected Debug target
-> load the ELF into SRAM
-> load debug symbols
-> run to main
-> keep the debugger attached
\end{verbatim}

\texttt{BOOT1} must be set to \texttt{1} for this development-mode workflow.

\section{Attach to an already-running target}\label{wiki-getting-started-build-flash-debug--10-attach-to-an-already-running-target}

An already-running application can be debugged without intentionally resetting or reloading it.

Use the VS Code Run and Debug configuration:

\begin{verbatim}
Attach: Application (no reset)
\end{verbatim}

or:

\begin{verbatim}
Attach: FSBL (no reset)
\end{verbatim}

This loads only the ELF symbols and attaches to the target.

Typical workflow:

\begin{verbatim}
1. Load application into SRAM and detach
2. Let the firmware run
3. Later select "Attach: Application (no reset)"
\end{verbatim}

The same attach configuration can also be used after a normal standalone boot from external flash, provided debug access is still available.

\section{Development mode vs normal boot}\label{wiki-getting-started-build-flash-debug--11-development-mode-vs-normal-boot}

\subsection{Development mode}\label{wiki-getting-started-build-flash-debug--development-mode}

Use:

\begin{verbatim}
BOOT1 = 1
\end{verbatim}

Typical workflow:

\begin{verbatim}
build
-> load ELF directly into SRAM
-> debug/run
\end{verbatim}

No external flash programming is required.

The program is lost after reset or power removal.

\subsection{Normal standalone boot}\label{wiki-getting-started-build-flash-debug--normal-standalone-boot}

Program the persistent external flash first, then select the normal boot configuration.

Current boot chain:

\begin{verbatim}
BootROM
-> external SEMPER flash
-> FSBL at 0x90040000
-> application at 0x90100000
\end{verbatim}

After changing boot mode, reset or power-cycle the MCU.

\section{CMake compiler cache}\label{wiki-getting-started-build-flash-debug--12-cmake-compiler-cache}

CMake stores the absolute compiler path in its cache.

If the project was previously configured using Ubuntu\textquotesingle s generic compiler:

\begin{verbatim}
/usr/bin/arm-none-eabi-gcc
\end{verbatim}

instead of the STM32Cube GNU toolchain, remove the existing build directories once:

\begin{verbatim}
rm -rf build/Debug
rm -rf build/Release
rm -rf FSBL/build
rm -rf Appli/build
rm -rf ExtMemLoader/build
\end{verbatim}

Then rebuild:

\begin{verbatim}
./scripts/build.sh --config Debug --target all
\end{verbatim}

During configuration/build, the compiler path should point to the STM32Cube bundle, for example:

\begin{verbatim}
~/.local/share/stm32cube/bundles/gnu-tools-for-stm32/.../bin/arm-none-eabi-gcc
\end{verbatim}

and not:

\begin{verbatim}
/usr/bin/arm-none-eabi-gcc
\end{verbatim}

\section{Optional tool path overrides}\label{wiki-getting-started-build-flash-debug--13-optional-tool-path-overrides}

The scripts normally auto-detect the required STM32 tools.

If necessary, individual paths can be overridden using environment variables.

Example:

\begin{verbatim}
export STM32_GCC_PATH=/path/to/gnu-tools-for-stm32/bin
export STM32_SIGNING_TOOL=/path/to/STM32_SigningTool_CLI
export STM32_PROGRAMMER_CLI=/path/to/STM32_Programmer_CLI
export STLINK_GDBSERVER=/path/to/ST-LINK_gdbserver
\end{verbatim}

For CMake or objcopy:

\begin{verbatim}
export CMAKE_BIN=/path/to/cmake
export OBJCOPY_BIN=/path/to/arm-none-eabi-objcopy
\end{verbatim}

These overrides are optional and should normally not be committed with machine-specific absolute paths.

\section{Quick reference}\label{wiki-getting-started-build-flash-debug--14-quick-reference}

\subsection{Build application}\label{wiki-getting-started-build-flash-debug--build-application}

\begin{verbatim}
./scripts/build.sh --config Debug --target appli
\end{verbatim}

\subsection{Build everything}\label{wiki-getting-started-build-flash-debug--build-everything}

\begin{verbatim}
./scripts/build.sh --config Debug --target all
\end{verbatim}

\subsection{Build release images}\label{wiki-getting-started-build-flash-debug--build-release-images}

\begin{verbatim}
./scripts/build.sh --config Release --target all
\end{verbatim}

\subsection{Program FSBL + Application A}\label{wiki-getting-started-build-flash-debug--program-fsbl--application-a}

\begin{verbatim}
./scripts/flash.sh --target all --config Debug --build
\end{verbatim}

\subsection{Program already-built images}\label{wiki-getting-started-build-flash-debug--program-already-built-images}

\begin{verbatim}
./scripts/flash.sh --target all
\end{verbatim}

\subsection{Run application from SRAM and detach}\label{wiki-getting-started-build-flash-debug--run-application-from-sram-and-detach}

\begin{verbatim}
./scripts/run_dev.sh --action run --target appli --config Debug --build
\end{verbatim}

\subsection{Run FSBL from SRAM and detach}\label{wiki-getting-started-build-flash-debug--run-fsbl-from-sram-and-detach}

\begin{verbatim}
./scripts/run_dev.sh --action run --target fsbl --config Debug --build
\end{verbatim}

\subsection{Make scripts executable after cloning}\label{wiki-getting-started-build-flash-debug--make-scripts-executable-after-cloning}

\begin{verbatim}
chmod +x scripts/*.sh
\end{verbatim}

\section{Recommended everyday workflow}\label{wiki-getting-started-build-flash-debug--15-recommended-everyday-workflow}

For application development:

\begin{verbatim}
BOOT1 = 1
-> Run and Debug
-> Dev: Debug Application from SRAM
\end{verbatim}

For quick standalone SRAM execution:

\begin{verbatim}
BOOT1 = 1
-> Run Task
-> Dev: Load + Run SRAM (detach)
\end{verbatim}

For persistent firmware testing:

\begin{verbatim}
Build / sign
-> Flash: Build + Program
-> select normal boot mode
-> reset / power-cycle
\end{verbatim}

For debugging firmware that is already running:

\begin{verbatim}
Run and Debug
-> Attach: Application (no reset)
\end{verbatim}

\chapter{Electrical ratings and characteristics}\label{wiki-hardware-electrical-ratings--electrical-ratings-and-characteristics}

{\small\color{MGSMuted}Source: \nolinkurl{wiki/hardware/electrical-ratings.md}\par}\medskip

The value cells below are intentionally blank. They are unassigned board-level ratings, not zeroes, typical component ratings, or values inferred from USB-PD profiles. Fill them in from approved design limits and board characterization.

These tables are the shared source for the user guide and the datasheet\textquotesingle s electrical section. Edit this page, then run \texttt{python3 scripts/\allowbreak{}build\_\allowbreak{}docs.\allowbreak{}py} to update both LaTeX documents. The datasheet section between its generated markers must not be edited separately.

\section{Absolute maximum ratings}\label{wiki-hardware-electrical-ratings--absolute-maximum-ratings}

{\def\LTcaptype{none} % do not increment counter
\begin{longtable}[]{@{}
  >{\raggedright\arraybackslash}p{(\linewidth - 8\tabcolsep) * \real{0.3128}}
  >{\raggedright\arraybackslash}p{(\linewidth - 8\tabcolsep) * \real{0.1759}}
  >{\raggedright\arraybackslash}p{(\linewidth - 8\tabcolsep) * \real{0.1704}}
  >{\raggedright\arraybackslash}p{(\linewidth - 8\tabcolsep) * \real{0.1704}}
  >{\raggedright\arraybackslash}p{(\linewidth - 8\tabcolsep) * \real{0.1704}}@{}}
\toprule\noalign{}
\begin{minipage}[b]{\linewidth}\raggedright
\textbf{Parameter}
\end{minipage} & \begin{minipage}[b]{\linewidth}\raggedright
\textbf{Symbol}
\end{minipage} & \begin{minipage}[b]{\linewidth}\raggedright
\textbf{Min}
\end{minipage} & \begin{minipage}[b]{\linewidth}\raggedright
\textbf{Max}
\end{minipage} & \begin{minipage}[b]{\linewidth}\raggedright
\textbf{Unit}
\end{minipage} \\
\midrule\noalign{}
\endhead
\bottomrule\noalign{}
\endlastfoot
Board input voltage & VIN & & & V \\
USB VBUS voltage & VBUS & & & V \\
Exposed signal pin voltage & VIO & & & V \\
Board input current & IIN & & & A \\
Storage temperature & TSTG & & & °C \\
\end{longtable}
}

\section{Recommended operating conditions}\label{wiki-hardware-electrical-ratings--recommended-operating-conditions}

{\def\LTcaptype{none} % do not increment counter
\begin{longtable}[]{@{}
  >{\raggedright\arraybackslash}p{(\linewidth - 10\tabcolsep) * \real{0.2601}}
  >{\raggedright\arraybackslash}p{(\linewidth - 10\tabcolsep) * \real{0.1511}}
  >{\raggedright\arraybackslash}p{(\linewidth - 10\tabcolsep) * \real{0.1472}}
  >{\raggedright\arraybackslash}p{(\linewidth - 10\tabcolsep) * \real{0.1472}}
  >{\raggedright\arraybackslash}p{(\linewidth - 10\tabcolsep) * \real{0.1472}}
  >{\raggedright\arraybackslash}p{(\linewidth - 10\tabcolsep) * \real{0.1472}}@{}}
\toprule\noalign{}
\begin{minipage}[b]{\linewidth}\raggedright
\textbf{Parameter}
\end{minipage} & \begin{minipage}[b]{\linewidth}\raggedright
\textbf{Symbol}
\end{minipage} & \begin{minipage}[b]{\linewidth}\raggedright
\textbf{Min}
\end{minipage} & \begin{minipage}[b]{\linewidth}\raggedright
\textbf{Typ}
\end{minipage} & \begin{minipage}[b]{\linewidth}\raggedright
\textbf{Max}
\end{minipage} & \begin{minipage}[b]{\linewidth}\raggedright
\textbf{Unit}
\end{minipage} \\
\midrule\noalign{}
\endhead
\bottomrule\noalign{}
\endlastfoot
Board input voltage & VIN & & & & V \\
Digital interface voltage & VIO & & & & V \\
Ambient operating temperature & TA & & & & °C \\
Relative humidity & RH & & & & \% \\
\end{longtable}
}

\section{Electrical characteristics}\label{wiki-hardware-electrical-ratings--electrical-characteristics}

{\def\LTcaptype{none} % do not increment counter
\begin{longtable}[]{@{}
  >{\raggedright\arraybackslash}p{(\linewidth - 12\tabcolsep) * \real{0.1986}}
  >{\raggedright\arraybackslash}p{(\linewidth - 12\tabcolsep) * \real{0.1363}}
  >{\raggedright\arraybackslash}p{(\linewidth - 12\tabcolsep) * \real{0.1484}}
  >{\raggedright\arraybackslash}p{(\linewidth - 12\tabcolsep) * \real{0.1292}}
  >{\raggedright\arraybackslash}p{(\linewidth - 12\tabcolsep) * \real{0.1292}}
  >{\raggedright\arraybackslash}p{(\linewidth - 12\tabcolsep) * \real{0.1292}}
  >{\raggedright\arraybackslash}p{(\linewidth - 12\tabcolsep) * \real{0.1292}}@{}}
\toprule\noalign{}
\begin{minipage}[b]{\linewidth}\raggedright
\textbf{Parameter}
\end{minipage} & \begin{minipage}[b]{\linewidth}\raggedright
\textbf{Symbol}
\end{minipage} & \begin{minipage}[b]{\linewidth}\raggedright
\textbf{Conditions}
\end{minipage} & \begin{minipage}[b]{\linewidth}\raggedright
\textbf{Min}
\end{minipage} & \begin{minipage}[b]{\linewidth}\raggedright
\textbf{Typ}
\end{minipage} & \begin{minipage}[b]{\linewidth}\raggedright
\textbf{Max}
\end{minipage} & \begin{minipage}[b]{\linewidth}\raggedright
\textbf{Unit}
\end{minipage} \\
\midrule\noalign{}
\endhead
\bottomrule\noalign{}
\endlastfoot
Input current, ToF idle & IIDLE & & & & & mA \\
Input current, ToF ranging & IRANGE & & & & & mA \\
Input power, ToF idle & PIDLE & & & & & W \\
Input power, ToF ranging & PRANGE & & & & & W \\
USB VBUS inrush current & IINRUSH & & & & & A \\
Supply ripple & VRIPPLE & & & & & mVpp \\
Power-on startup time & TSTART & & & & & ms \\
\end{longtable}
}

Record the PCB revision, assembly variant, firmware, supply conditions, sensor configuration, and measurement method when assigning values. Keep absolute limits, recommended conditions, and measured characteristics distinct.

\chapter{Current external flash layout}\label{wiki-hardware-flash-layout--current-external-flash-layout}

{\small\color{MGSMuted}Source: \nolinkurl{wiki/hardware/flash-layout.md}\par}\medskip

\section{Address conventions}\label{wiki-hardware-flash-layout--address-conventions}

The S28HS512T is a 64 MiB NOR device connected through XSPI1 and XSPIM port 2. Its memory-mapped base is \texttt{0x90000000}; the exclusive end is \texttt{0x94000000}. The flash driver uses device-relative offsets, while CubeProgrammer uses the absolute addresses shown below.

\section{Programming regions}\label{wiki-hardware-flash-layout--programming-regions}

{\def\LTcaptype{none} % do not increment counter
\begin{longtable}[]{@{}
  >{\raggedright\arraybackslash}p{(\linewidth - 6\tabcolsep) * \real{0.3174}}
  >{\raggedright\arraybackslash}p{(\linewidth - 6\tabcolsep) * \real{0.2276}}
  >{\raggedright\arraybackslash}p{(\linewidth - 6\tabcolsep) * \real{0.2200}}
  >{\raggedright\arraybackslash}p{(\linewidth - 6\tabcolsep) * \real{0.2349}}@{}}
\toprule\noalign{}
\begin{minipage}[b]{\linewidth}\raggedright
\textbf{Purpose}
\end{minipage} & \begin{minipage}[b]{\linewidth}\raggedright
\textbf{Flash offset}
\end{minipage} & \begin{minipage}[b]{\linewidth}\raggedright
\textbf{CPU address}
\end{minipage} & \begin{minipage}[b]{\linewidth}\raggedright
\textbf{Reserved size}
\end{minipage} \\
\midrule\noalign{}
\endhead
\bottomrule\noalign{}
\endlastfoot
Parameter JSON and padding & \texttt{0x00000000} & \texttt{0x90000000} & 256 KiB \\
FSBL2 & \texttt{0x00040000} & \texttt{0x90040000} & 512 KiB \\
Reserved & \texttt{0x000C0000} & \texttt{0x900C0000} & 256 KiB \\
Application A & \texttt{0x00100000} & \texttt{0x90100000} & 8 MiB \\
Application B & \texttt{0x00900000} & \texttt{0x90900000} & 8 MiB \\
Update metadata copy A & \texttt{0x01100000} & \texttt{0x91100000} & 256 KiB \\
Update metadata copy B & \texttt{0x01140000} & \texttt{0x91140000} & 256 KiB \\
Future data & \texttt{0x01180000} & \texttt{0x91180000} & 46.5 MiB \\
\end{longtable}
}

Application flash slots are larger than the current execution allocation. An LRUN application image, including its STM32 header and padding, must fit in 512 KiB. FSBL copies it to \texttt{0x34000000}; the vector table starts at \texttt{0x34000400}.

\section{Parameters and the boot header}\label{wiki-hardware-flash-layout--parameters-and-the-boot-header}

The current parameter packer writes JSON at offset zero, followed by a NUL terminator and erased-value padding across the first 256 KiB. The application reads at most the first 16 KiB. JSON is not a valid STM32 FSBL1 header, so the board uses its FSBL2 image at offset \texttt{0x00040000}.

The \hyperref[wiki-reference-flash-layout-proposal--stm32n6--s28hs512t-external-flash-memory-map]{original proposal} described a separate 4 KiB guard followed by individual configuration sectors. That is design history; the current \hyperref[wiki-firmware-parameters--device-parameters-and-build-metadata]{parameter workflow} replaces the complete first 256 KiB region. It does not preserve separate metadata sectors inside it.

\section{A/B updates and metadata}\label{wiki-hardware-flash-layout--ab-updates-and-metadata}

The application updater chooses the inactive application slot. The confirmed image, parameters, and FSBL are excluded from its erase targets. Two separate 256 KiB sectors hold alternating committed update records, including image hashes, confirmed/pending slots, trial state, and persistent status flags.

These metadata reservations occupy part of the former future/data region so that parameter programming cannot erase boot state. See \hyperref[wiki-firmware-firmware-updates--application-firmware-updates-and-ab-recovery]{firmware updates and A/B recovery} for the commit sequence, watchdog, validation rules, and direct-SWD reprovisioning steps.

\section{Erase geometry}\label{wiki-hardware-flash-layout--erase-geometry}

{\def\LTcaptype{none} % do not increment counter
\begin{longtable}[]{@{}
  >{\raggedright\arraybackslash}p{(\linewidth - 2\tabcolsep) * \real{0.4686}}
  >{\raggedright\arraybackslash}p{(\linewidth - 2\tabcolsep) * \real{0.5314}}@{}}
\toprule\noalign{}
\begin{minipage}[b]{\linewidth}\raggedright
\textbf{Offset range}
\end{minipage} & \begin{minipage}[b]{\linewidth}\raggedright
\textbf{Physical geometry}
\end{minipage} \\
\midrule\noalign{}
\endhead
\bottomrule\noalign{}
\endlastfoot
\texttt{0x00000000}--\texttt{0x0001FFFF} & 32 × 4 KiB parameter sectors \\
\texttt{0x00020000}--\texttt{0x0003FFFF} & One 128 KiB transition region \\
\texttt{0x00040000}--\texttt{0x03FFFFFF} & 255 × 256 KiB main sectors \\
\end{longtable}
}

The external loader must use this hybrid geometry when replacing parameters. The \hyperref[wiki-firmware-parameters--troubleshooting-parameter-programming]{parameter troubleshooting guide} explains the earlier uniform-erase failure and rebuilding the corrected loader.

\section{Implementation references}\label{wiki-hardware-flash-layout--implementation-references}

\begin{itemize}
\tightlist
\item
  {Update slot, metadata, and image limits (\nolinkurl{common/firmware/fw_update.h})}
\item
  {Update flash-access bounds (\nolinkurl{common/firmware/fw_flash.c})}
\item
  {Persistent programming addresses (\nolinkurl{scripts/flash.sh})}
\item
  {Parameter packing and size limits (\nolinkurl{scripts/pack_parameters.py})}
\item
  {Application SRAM linker layout (\nolinkurl{Appli/STM32N657XX_LRUN.ld})}
\end{itemize}

\chapter{Boot and application status LED}\label{wiki-firmware-status-led--boot-and-application-status-led}

{\small\color{MGSMuted}Source: \nolinkurl{wiki/firmware/status-led.md}\par}\medskip

The same RGB LED indicates progress through ROM-loaded FSBL and application startup, then follows the acknowledged ToF sensor state. Both firmware images need the updated LED controller to show the complete sequence.

{\def\LTcaptype{none} % do not increment counter
\begin{longtable}[]{@{}
  >{\raggedright\arraybackslash}p{(\linewidth - 4\tabcolsep) * \real{0.2693}}
  >{\raggedright\arraybackslash}p{(\linewidth - 4\tabcolsep) * \real{0.3043}}
  >{\raggedright\arraybackslash}p{(\linewidth - 4\tabcolsep) * \real{0.4263}}@{}}
\toprule\noalign{}
\begin{minipage}[b]{\linewidth}\raggedright
\textbf{State}
\end{minipage} & \begin{minipage}[b]{\linewidth}\raggedright
\textbf{LED pattern}
\end{minipage} & \begin{minipage}[b]{\linewidth}\raggedright
\textbf{Trigger}
\end{minipage} \\
\midrule\noalign{}
\endhead
\bottomrule\noalign{}
\endlastfoot
FSBL started & Amber blink, about 2 Hz & FSBL has configured GPIO/I2C, before flash initialization. \\
FSBL loading & Amber blink, about 5 Hz & Reading A/B metadata, validating a slot and loading it into SRAM. \\
FSBL handoff & Solid magenta & Copied application passed SHA-256 verification; the branch is being prepared. \\
FSBL error & Counted amber burst, then 2 s dark & FSBL error handler; count identifies the boot stage below. \\
Application started & Blue blink, about 2 Hz & Application has configured GPIO/I2C, before the remaining peripheral initialization. \\
Application core init finished & Cyan double pulse, then a pause & Application entry thread starts after RTOS and middleware initialization completes. Sensor initialization can still be in progress. \\
Ranging & Green breathing, about 1.2 s per cycle & ToF successfully enters streaming/ranging state. \\
Idle & Dim blue breathing, about 2.4 s per cycle & ToF successfully stops, or a sensor reset finishes with ranging stopped. \\
ToF powered down & Dim steady amber & Sensor XSHUT shutdown completed. \\
ToF resetting & Cyan blink, about 5 Hz & Reset command is reinitializing the sensor. \\
ToF fault & Red double pulse, then a pause & Latched initialization, capture or ranging fault; explicit Reset ToF is required. \\
Fatal application initialization error & Counted red, blue, or magenta burst, then 2 s dark & Application error handler; color and count identify the startup stage below. \\
\end{longtable}
}

Early boot transitions can be brief; no artificial boot delays are introduced. Green means ToF ranging was started, independently of a USB/Ethernet client. Idle and power-down apply to ToF only: IMU, networking and USB can keep running. Stop/Start/Reset/Power down update the LED after the corresponding state change; rejected commands leave the current state indication intact.

Patterns run autonomously in the LED chip. They continue without CPU animation timers, and can remain visible if the CPU later hangs. A surviving pattern is the last reported state, not proof that the processor is still making progress. No indication can mean failure before GPIO/I2C initialization or an unavailable LED device; magenta only confirms preparation for handoff, not successful application execution. See \hyperref[wiki-firmware-firmware-updates--diagnosing-standalone-boot]{standalone boot diagnostics}.

\section{Counted startup failure codes}\label{wiki-firmware-status-led--counted-startup-failure-codes}

Each failure burst has \textbf{N flashes, 200 ms on / 200 ms off, followed by 2 s dark}, repeating autonomously. Count a complete burst after the long gap. Normal continuous amber/blue startup blinks have no long gap. These codes require the new FSBL and application images; an old image can still show an old pattern.

\textbf{Amber: FSBL failure before application entry}

{\def\LTcaptype{none} % do not increment counter
\begin{longtable}[]{@{}
  >{\raggedright\arraybackslash}p{(\linewidth - 2\tabcolsep) * \real{0.2784}}
  >{\raggedright\arraybackslash}p{(\linewidth - 2\tabcolsep) * \real{0.7216}}@{}}
\toprule\noalign{}
\begin{minipage}[b]{\linewidth}\raggedright
\textbf{Count}
\end{minipage} & \begin{minipage}[b]{\linewidth}\raggedright
\textbf{Stage active when the FSBL error handler was reached}
\end{minipage} \\
\midrule\noalign{}
\endhead
\bottomrule\noalign{}
\endlastfoot
1 & LED setup / DMA setup before flash initialization \\
2 & XSPI controller or external flash initialization \\
3 & Reading/opening the two A/B metadata records \\
4 & Validating/selecting a boot slot (including initial image adoption and metadata writes), or starting its trial watchdog \\
5 & Copying the selected image from flash into SRAM \\
6 & Verifying the SRAM copy\textquotesingle s SHA-256 \\
7 & Preparing application handoff (reserved if this stage reports an error) \\
\end{longtable}
}

For amber 3 or 4 after direct SWD replacement of an application, check the \hyperref[wiki-firmware-firmware-updates--flash-allocation-and-limits]{A/B metadata reprovisioning notes}. Existing descriptors can still refer to the old image hash; replacing a binary alone does not update those descriptors. Use \textbf{Flash: Reprovision Release Baseline}, or \texttt{scripts/\allowbreak{}flash.\allowbreak{}sh -\allowbreak{}-\allowbreak{}target all -\allowbreak{}-\allowbreak{}config Release -\allowbreak{}-\allowbreak{}build -\allowbreak{}-\allowbreak{}reset-\allowbreak{}boot-\allowbreak{}metadata}, to back up and reset only those boot records after verifying the baseline. Do not erase unrelated device data.

\textbf{Red: application peripheral initialization}

{\def\LTcaptype{none} % do not increment counter
\begin{longtable}[]{@{}
  >{\raggedright\arraybackslash}p{(\linewidth - 2\tabcolsep) * \real{0.3396}}
  >{\raggedright\arraybackslash}p{(\linewidth - 2\tabcolsep) * \real{0.6604}}@{}}
\toprule\noalign{}
\begin{minipage}[b]{\linewidth}\raggedright
\textbf{Count}
\end{minipage} & \begin{minipage}[b]{\linewidth}\raggedright
\textbf{Stage}
\end{minipage} \\
\midrule\noalign{}
\endhead
\bottomrule\noalign{}
\endlastfoot
1 & HAL / tick initialization \\
2 & System or common peripheral clocks \\
3 & GPIO \\
4 & I2C2 \\
5 & GPDMA1 \\
6 & DCMIPP / CSI capture setup \\
7 & Ethernet MAC / PHY setup \\
8 & UCPD1 \\
9 & USB OTG HS peripheral \\
10 & SPI4 \\
11 & I3C2 \\
12 & XSPI1 \\
13 & System isolation / RIF setup \\
14 & USB-PD pre-OS setup \\
\end{longtable}
}

Errors before application I2C/LED initialization cannot program a new code; the previous FSBL pattern can remain. Codes identify the active initialization call, including nested MSP/clock errors; they do not identify its underlying cause.

\textbf{Blue: application startup services}

{\def\LTcaptype{none} % do not increment counter
\begin{longtable}[]{@{}
  >{\raggedright\arraybackslash}p{(\linewidth - 2\tabcolsep) * \real{0.3158}}
  >{\raggedright\arraybackslash}p{(\linewidth - 2\tabcolsep) * \real{0.6842}}@{}}
\toprule\noalign{}
\begin{minipage}[b]{\linewidth}\raggedright
\textbf{Count}
\end{minipage} & \begin{minipage}[b]{\linewidth}\raggedright
\textbf{Stage}
\end{minipage} \\
\midrule\noalign{}
\endhead
\bottomrule\noalign{}
\endlastfoot
1 & Reading device parameters \\
2 & Entering ThreadX before middleware initialization \\
3 & Runtime: application entry thread has started \\
\end{longtable}
}

\textbf{Magenta: RTOS pools / middleware setup}

{\def\LTcaptype{none} % do not increment counter
\begin{longtable}[]{@{}
  >{\raggedright\arraybackslash}p{(\linewidth - 2\tabcolsep) * \real{0.3338}}
  >{\raggedright\arraybackslash}p{(\linewidth - 2\tabcolsep) * \real{0.6662}}@{}}
\toprule\noalign{}
\begin{minipage}[b]{\linewidth}\raggedright
\textbf{Count}
\end{minipage} & \begin{minipage}[b]{\linewidth}\raggedright
\textbf{Stage}
\end{minipage} \\
\midrule\noalign{}
\endhead
\bottomrule\noalign{}
\endlastfoot
1 & ThreadX application memory pool \\
2 & Application thread / service creation \\
3 & NetX Duo memory pool \\
4 & NetX Duo initialization \\
5 & USBX memory pool \\
6 & USBX initialization \\
7 & USB-PD memory pool \\
8 & USB-PD initialization \\
\end{longtable}
}

When debug access is available, \texttt{app\_\allowbreak{}boot\_\allowbreak{}stage}, \texttt{app\_\allowbreak{}boot\_\allowbreak{}failed\_\allowbreak{}stage}, and \texttt{app\_\allowbreak{}boot\_\allowbreak{}error\_\allowbreak{}return\_\allowbreak{}address} retain the current checkpoint, failing checkpoint, and return address of the caller of \texttt{Error\_\allowbreak{}Handler()}. These are live RAM symbols, not a persistent reset log. FSBL retains \texttt{fw\_\allowbreak{}boot\_\allowbreak{}stage} / \texttt{fw\_\allowbreak{}boot\_\allowbreak{}result}.

\subsection{Earlier continuous-red ambiguity}\label{wiki-firmware-status-led--earlier-continuous-red-ambiguity}

The earlier double-pulse preset used infinite \textbf{pattern} repetition, preventing its final pause from being reached. An FSBL error or ToF fault could therefore look like the application\textquotesingle s continuous red 5 Hz fatal indication. Continuous red on those images does not prove that the jump reached the application. The driver now runs the double-pulse pattern once and repeats the complete \textbf{engine}, including its pause. Counted codes also use finite pattern repetition followed by the engine\textquotesingle s complete-cycle repeat, as described in \href{https://www.ti.com/lit/gpn/lp5815}{TI LP5815 sections 7.3.4.1--7.3.4.2}.

\section{Driver behavior}\label{wiki-firmware-status-led--driver-behavior}

The register map matches \href{https://www.ti.com/lit/gpn/lp5815}{TI LP5815}. Timing codes, infinite repetition, engine routing and commands follow its autonomous animation registers. The existing channel current settings and red/green/blue output mapping are retained. Legacy animation names remain available; lifecycle code uses the explicit FSBL/application names.

\texttt{LEDController\_\allowbreak{}Init()} runs once after I2C setup in each firmware image. The header declares one shared controller instance. Every animation is constructed from cleared configuration, with explicitly encoded register bytes, so earlier patterns and compiler bitfield layout cannot alter later states.

\texttt{LEDController\_\allowbreak{}SetAnimation()} returns the HAL transfer status. Repeating a successful request generates no I2C traffic and does not restart the animation. Failed repeated requests retry at most every 500 ms; the sensor owner thread services retries even while idle or waiting for a frame. A changed state can be attempted immediately. A concurrent writer receives \texttt{HAL\_\allowbreak{}BUSY}; the driver masks interrupts only while guarding its ownership flag and restores the caller\textquotesingle s interrupt state before I2C transfers. Calls belong in startup or thread context. \texttt{led\_\allowbreak{}controller.\allowbreak{}last\_\allowbreak{}status} and \texttt{animation\_\allowbreak{}valid} expose unsuccessful writes when debugger access is available.

\section{Validation}\label{wiki-firmware-status-led--validation}

\begin{verbatim}
python3 -m unittest discover -s test_scripts -p test_led_controller.py -v
python3 -m unittest discover -s test_scripts -p test_sensor_control.py -v
\end{verbatim}

Native tests check actual configuration/engine/pattern register bytes, all preset transitions, infinite playback, duplicate suppression, failed transfer steps, retry timing and wraparound, concurrent calls, and ToF state mapping. An independent waveform model checks that double-pulse pauses actually execute, that each diagnostic burst contains the requested number of flashes, and that the full 2 s gap is dark. Tests cover every application startup-code mapping. The sensor harness verifies LED state follows Stop/Start/Reset/Power down/fault acknowledgements. Cross-build FSBL and Appli before a board check; confirm colors and cadence on the physical LED after flashing the images.

\chapter{Device parameters and build metadata}\label{wiki-firmware-parameters--device-parameters-and-build-metadata}

{\small\color{MGSMuted}Source: \nolinkurl{wiki/firmware/parameters.md}\par}\medskip

\section{Programming parameters}\label{wiki-firmware-parameters--programming-parameters}

For a complete build and flash:

\begin{verbatim}
./scripts/flash.sh --target all --config Release --build
\end{verbatim}

For an application-only update:

\begin{verbatim}
./scripts/flash.sh --target appli-a --config Release --build
\end{verbatim}

Before programming, the script reads the connected board\textquotesingle s first 16 KiB of parameters. Existing identity, firmware version labels, calibration and custom settings are preserved; new fields get defaults from the template. An erased board starts from the template. The script updates only the selected firmware slots\textquotesingle{} build records and verifies every selected image before writing/verifying the parameters last. A readback or validation failure stops the operation before any writes. A firmware write failure stops before committing the new parameters; an interrupted operation can leave partially updated firmware, so rerun it.

\section{Replacing editable settings}\label{wiki-firmware-parameters--replacing-editable-settings}

To intentionally replace editable settings, edit the template and run:

\begin{verbatim}
./scripts/flash.sh --target parameters
\end{verbatim}

Or provide a device-specific file with a firmware update:

\begin{verbatim}
./scripts/flash.sh --target appli-a --parameters /path/to/settings.json
\end{verbatim}

Both commands replace editable settings with the selected file while retaining other slots\textquotesingle{} build records from the connected board. A parameters-only update needs the external loader but no local firmware binaries or manifests. Invalid stored JSON requires an explicit settings replacement; a failed/short readback always stops the operation.

\section{Editable fields}\label{wiki-firmware-parameters--editable-fields}

Editable sections:

\begin{itemize}
\tightlist
\item
  \texttt{device}: product, PCB identifier, board version, name, description, serial number, hardware/PCB/assembly revisions.
\item
  \texttt{firmware}: manually maintained version, release number and protocol version.
\item
  \texttt{manufacturing}: date, batch and test timestamp.
\item
  \texttt{calibration}: confirmed calibration values, initially empty.
\end{itemize}

Use \texttt{null} for unassigned values. These fields are stored and displayed by the PC app; firmware does not apply them as runtime configuration.

The default product is \texttt{Magnecko Sense}, PCB identifier \texttt{MGS}, hardware version \texttt{1.\allowbreak{}0}, and PCB revision \texttt{-\allowbreak{}} (revision 0), giving board version \texttt{MGS-\allowbreak{}1.\allowbreak{}0-\allowbreak{}}. Revisions A and B are written \texttt{MGS-\allowbreak{}1.\allowbreak{}0A} and \texttt{MGS-\allowbreak{}1.\allowbreak{}0B}. These are hardware identifiers; \texttt{firmware.\allowbreak{}version} remains a separate software version.

When changing a board revision, keep \texttt{device.\allowbreak{}hardware\_\allowbreak{}revision}, \texttt{device.\allowbreak{}pcb\_\allowbreak{}revision}, and \texttt{device.\allowbreak{}board\_\allowbreak{}version} consistent with {\texttt{common/\allowbreak{}board/\allowbreak{}board\_\allowbreak{}identity.\allowbreak{}h} (\nolinkurl{common/board/board_identity.h})} and the USB product string in {\texttt{magnecko-\allowbreak{}sense-\allowbreak{}fw.\allowbreak{}ioc} (\nolinkurl{magnecko-sense-fw.ioc})}. The shared header supplies the firmware\textquotesingle s USB, sensor interface, and external loader identities. After CubeMX regeneration, keep \texttt{USBD\_\allowbreak{}PRODUCT\_\allowbreak{}STRING} mapped to \texttt{MGS\_\allowbreak{}USB\_\allowbreak{}PRODUCT\_\allowbreak{}STRING} in \texttt{ux\_\allowbreak{}device\_\allowbreak{}descriptors.\allowbreak{}h}.

Previously programmed identity values are preserved during ordinary firmware flashes. To apply the new branding to an existing board, merge the new identity fields into its device-specific settings file and program it with \texttt{-\allowbreak{}-\allowbreak{}target parameters -\allowbreak{}-\allowbreak{}parameters /\allowbreak{}path/\allowbreak{}to/\allowbreak{}settings.\allowbreak{}json}, retaining its serial number, calibration, and other per-device values.

\section{Generated metadata}\label{wiki-firmware-parameters--generated-metadata}

Generated sections (do not maintain these in the editable file):

\begin{itemize}
\tightlist
\item
  \texttt{parameter\_\allowbreak{}build}: \texttt{built\_\allowbreak{}at\_\allowbreak{}utc} refreshes on \textbf{every parameter build}, even standalone \texttt{pack\_\allowbreak{}parameters.\allowbreak{}py} calls without \texttt{-\allowbreak{}-\allowbreak{}target}; \texttt{source} names the editable JSON file. This timestamp appears under \textbf{parameter\_build} in the Device Monitor\textquotesingle s parameter tree and in \texttt{build/\allowbreak{}flash\_\allowbreak{}parameters.\allowbreak{}json}.
\item
  \texttt{repository}: Git commit, description, branch, dirty flag and submodule revisions at parameter-generation time.
\item
  \texttt{build.\allowbreak{}fsbl}, \texttt{build.\allowbreak{}application} (slot A), \texttt{build.\allowbreak{}application\_\allowbreak{}b} (slot B): Git state at packaging, image packaging timestamp (\texttt{image\_\allowbreak{}built\_\allowbreak{}at\_\allowbreak{}utc}), Debug/Release configuration, image filename/size/SHA-256, compiler name/version/target, CMake version/generator, toolchain filename, CMake options (\texttt{DEV\_\allowbreak{}MODE}, \texttt{VL53L9\_\allowbreak{}SENSOR\_\allowbreak{}ONLY} where applicable), actual compile flags/definitions for \texttt{Core/\allowbreak{}Src/\allowbreak{}main.\allowbreak{}c}, and preprocessed ThreadX settings. Each selected slot\textquotesingle s \texttt{built\_\allowbreak{}at\_\allowbreak{}utc} and \texttt{programmed\_\allowbreak{}at\_\allowbreak{}utc} refresh to the current host UTC time for every flash, including when programming the same existing image again. Other slots retain their records. Flashing also records each selected slot\textquotesingle s address.
\item
  \texttt{flash}: parameter-generation timestamp, selected target, updated image labels, connection mode and parameter-region address/size. This describes the latest completed parameter update; timestamps come from the host clock.
\end{itemize}

\section{Timestamps and image provenance}\label{wiki-firmware-parameters--timestamps-and-image-provenance}

A parameter-only rebuild updates \texttt{parameter\_\allowbreak{}build.\allowbreak{}built\_\allowbreak{}at\_\allowbreak{}utc} without changing \texttt{build.\allowbreak{}fsbl.\allowbreak{}built\_\allowbreak{}at\_\allowbreak{}utc}, \texttt{build.\allowbreak{}application.\allowbreak{}built\_\allowbreak{}at\_\allowbreak{}utc} or slot B\textquotesingle s date. When flashing Appli, FSBL or both, each selected firmware record\textquotesingle s \texttt{built\_\allowbreak{}at\_\allowbreak{}utc} refreshes too. \texttt{image\_\allowbreak{}built\_\allowbreak{}at\_\allowbreak{}utc} preserves the original build/packaging time from the verified image manifest. Building parameters does not rebuild firmware. When preparing a flash operation, \texttt{parameter\_\allowbreak{}build.\allowbreak{}built\_\allowbreak{}at\_\allowbreak{}utc}, the selected slots\textquotesingle{} \texttt{built\_\allowbreak{}at\_\allowbreak{}utc} / \texttt{programmed\_\allowbreak{}at\_\allowbreak{}utc}, and \texttt{flash.\allowbreak{}parameters\_\allowbreak{}generated\_\allowbreak{}at\_\allowbreak{}utc} share the same UTC timestamp. Firmware is verified before the new parameter block is committed; a failed firmware write leaves the board\textquotesingle s old records intact.

Image build details come from each packaged image\textquotesingle s adjacent \texttt{.\allowbreak{}manifest.\allowbreak{}json}; flashing existing images retains their original configuration and image provenance while refreshing the selected slots\textquotesingle{} record/flash timestamps. The image name, byte count and SHA-256 must match the manifest. Old manifests without the expanded build settings require a one-time rebuild using \texttt{scripts/\allowbreak{}build.\allowbreak{}sh} or \texttt{flash.\allowbreak{}sh -\allowbreak{}-\allowbreak{}build}. SRAM \texttt{DEV\_\allowbreak{}MODE=ON} application images must be rebuilt without \texttt{-\allowbreak{}-\allowbreak{}dev} before persistent flashing. Unselected local images are never used to claim what is installed in other slots. Older on-board records remain as recorded until those slots are flashed again.

\section{Storage layout}\label{wiki-firmware-parameters--storage-layout}

Layout: address \texttt{0x90000000}, reserved size \texttt{0x40000} (256 KiB). Schema version 2 uses compact UTF-8 JSON, one NUL byte and \texttt{0xFF} padding, still within the existing 16 KiB firmware reader limit. Updating parameters replaces this reserved block; FSBL starts at \texttt{0x90040000}, application A at \texttt{0x90100000}, B at \texttt{0x90900000}.

\section{Inspecting stored records}\label{wiki-firmware-parameters--inspecting-stored-records}

Inspect the generated, readable JSON in \texttt{build/\allowbreak{}flash\_\allowbreak{}parameters.\allowbreak{}json}; the padded flash image is \texttt{build/\allowbreak{}flash\_\allowbreak{}parameters.\allowbreak{}bin}. The PC viewer shows the expanded sections automatically under \textbf{Device \& firmware parameters}. Compiler paths, user/host names, credentials and signing keys are not added to these records.

\section{Troubleshooting parameter programming}\label{wiki-firmware-parameters--troubleshooting-parameter-programming}

Rebuild the custom external loader before retrying the 256 KiB parameter image:

\begin{verbatim}
./scripts/build.sh --target extmem --config Release
\end{verbatim}

Use the newly generated \texttt{ExtMemLoader/\allowbreak{}build/\allowbreak{}magnecko-\allowbreak{}sense-\allowbreak{}fw\_\allowbreak{}ExtMemLoader.\allowbreak{}stldr}, including when selecting the loader manually in CubeProgrammer. The first 256 KiB consists of \textbf{32 × 4 KiB + 1 × 128 KiB} sectors. The previous uniform-sector wrapper used 256 KiB erase commands for this hybrid area: it could leave the 4 KiB overlay unerased and extend into the FSBL at \texttt{0x90040000}. The board override now walks the advertised physical geometry and bounds the inclusive erase range. Loader data cache is disabled because SWD updates programming buffers between calls.

If an update was already attempted with the old loader, rebuild and program the full firmware to restore any affected FSBL content, rather than updating only parameters:

\begin{verbatim}
./scripts/flash.sh --target all --config Release --build
\end{verbatim}

Native regressions cover parameter-only boundaries and the full device range. Actual erase/program/verify still requires checking on the board; a Sector{[}0{]} message by itself does not uniquely identify the hardware failure.

\chapter{Application firmware updates and A/B recovery}\label{wiki-firmware-firmware-updates--application-firmware-updates-and-ab-recovery}

{\small\color{MGSMuted}Source: \nolinkurl{wiki/firmware/firmware-updates.md}\par}\medskip

\section{First installation}\label{wiki-firmware-firmware-updates--first-installation}

Install \textbf{both the new FSBL and this DFU-capable application} using the normal SWD workflow, without \texttt{-\allowbreak{}-\allowbreak{}dev}:

\begin{verbatim}
./scripts/flash.sh --target all --config Release --build
\end{verbatim}

Boot normally from flash and verify USB/Ethernet operation. On the first boot, FSBL validates existing STM32 v2.3 images, adopts valid A/B images, and creates boot metadata. Slot A is preferred for this initial migration. Installing the DFU-capable baseline is necessary so that rollback also restores update access. Thereafter normal updates use the uploader, which never programs FSBL.

Generated source integration is confined to existing \texttt{USER CODE} sections. The \texttt{.\allowbreak{}inc} implementations are included once from each project\textquotesingle s \texttt{main.\allowbreak{}c} to preserve CubeMX source lists and regeneration without editing generated code.

\section{Uploading and inspecting status}\label{wiki-firmware-firmware-updates--uploading-and-inspecting-status}

Firmware updates are available directly in \textbf{Device Monitor → Firmware update}. Select a packaged application, enter its release identifier, and press Upload image. The tab shows acknowledged progress, verification, A/B slot status, and persistent warnings. It can reset a staged image and monitor confirmation or fallback. USB streaming is paused and resumed automatically. See the \hyperref[wiki-tools-device-monitor--firmware-updates]{viewer instructions} for all controls.

The command-line client remains available for scripts and terminal use. Use the packaged application binary from a \textbf{non-DEV\_MODE} build:

\begin{verbatim}
./scripts/build.sh --target appli --config Release
python3 scripts/firmware_update.py --tcp 192.168.8.20 \
  --image Appli/build/magnecko-sense-fw_Appli-signed.bin --version 2
\end{verbatim}

For USB, close Device Monitor/other serial owners first:

\begin{verbatim}
python3 scripts/firmware_update.py --usb /dev/ttyACM0 \
  --image Appli/build/magnecko-sense-fw_Appli-signed.bin --version 2
\end{verbatim}

The USB uploader requires \texttt{pyserial}. \texttt{-\allowbreak{}-\allowbreak{}version} is a uint32 release identifier, not an anti-rollback counter. The updater selects the inactive slot itself. The client validates the wrapper/checksum and sends a SHA-256 digest; the board verifies each programmed chunk and the entire stored image before committing it. By default the client requests reset after verification. \texttt{-\allowbreak{}-\allowbreak{}no-\allowbreak{}reset} leaves the image staged; a later \texttt{-\allowbreak{}-\allowbreak{}reset} starts its trial boot.

After reset, reconnect after roughly ten seconds and inspect status:

\begin{verbatim}
python3 scripts/firmware_update.py --tcp 192.168.8.20
python3 scripts/firmware_update.py --usb /dev/ttyACM0
\end{verbatim}

For a failed startup, allow roughly 32 seconds plus boot time for the watchdog and fallback. The USB device node may change after reset; a \texttt{/\allowbreak{}dev/\allowbreak{}serial/\allowbreak{}by-\allowbreak{}id/\allowbreak{}} path is preferable when available. Slot numbers are 0 = A, 1 = B; \texttt{4294967295} means no slot. Successful activation has \texttt{running == confirmed}, no pending slot, and \texttt{boot\_\allowbreak{}confirmed} set. A fallback keeps \texttt{fallback\_\allowbreak{}warning} set and records the cause in \texttt{last\_\allowbreak{}error}. Warnings remain until explicitly cleared:

\begin{verbatim}
python3 scripts/firmware_update.py --tcp 192.168.8.20 --clear-warnings
python3 scripts/firmware_update.py --tcp 192.168.8.20 --abort
python3 scripts/firmware_update.py --tcp 192.168.8.20 --reset
\end{verbatim}

\texttt{-\allowbreak{}-\allowbreak{}abort} cancels an in-progress transfer, not an already verified pending image. A lost connection aborts an incomplete upload; inactivity for 60 seconds also aborts it. A power interruption is recorded as an upload failure on next boot. Uploads restart from the beginning; chunks are not blindly retried or resumed. If an acknowledgement is lost, inspect status before choosing abort or reset. Only one USB/TCP session may own an upload. USB telemetry is suppressed during an upload, and a packet mutex prevents interleaving DFU replies and telemetry. Existing USB \texttt{CTRL} framing continues to work outside DFU operations.

\section{Flash allocation and limits}\label{wiki-firmware-firmware-updates--flash-allocation-and-limits}

All offsets below are relative to XSPI1\textquotesingle s \texttt{0x90000000} bank.

{\def\LTcaptype{none} % do not increment counter
\begin{longtable}[]{@{}
  >{\raggedright\arraybackslash}p{(\linewidth - 4\tabcolsep) * \real{0.4084}}
  >{\raggedright\arraybackslash}p{(\linewidth - 4\tabcolsep) * \real{0.2570}}
  >{\raggedright\arraybackslash}p{(\linewidth - 4\tabcolsep) * \real{0.3346}}@{}}
\toprule\noalign{}
\begin{minipage}[b]{\linewidth}\raggedright
\textbf{Region}
\end{minipage} & \begin{minipage}[b]{\linewidth}\raggedright
\textbf{Offset}
\end{minipage} & \begin{minipage}[b]{\linewidth}\raggedright
\textbf{Reserved size}
\end{minipage} \\
\midrule\noalign{}
\endhead
\bottomrule\noalign{}
\endlastfoot
Application A & \texttt{0x00100000} & 8 MiB \\
Application B & \texttt{0x00900000} & 8 MiB \\
Update metadata copy A & \texttt{0x01100000} & 256 KiB \\
Update metadata copy B & \texttt{0x01140000} & 256 KiB \\
Remaining future/data space & \texttt{0x01180000} & Until end of flash \\
\end{longtable}
}

The two metadata sectors are \textbf{new reservations in the previous future/data area}, rather than the older proposed 4 KiB metadata offsets in the parameter area. \texttt{flash.\allowbreak{}sh} currently replaces the first 256 KiB parameter region, so using that area would erase boot state during parameter programming. Factory data, parameters, FSBL and the confirmed application are never DFU erase targets. The hardware adapter also restricts writes and erases to these explicit ranges.

Although flash slots are 8 MiB, the current LRUN code allocation is 512 KiB. Uploaded images must be at most \texttt{0x80000} bytes including the STM32 v2.3 header and alignment padding. Vectors are at offset \texttt{0x400}; the signed file is copied to \texttt{0x34000000}, with code starting at \texttt{0x34000400}. Raw \texttt{.\allowbreak{}bin}/ELF files and images for other load layouts are rejected.

\textbf{SWD reprogramming after A/B metadata exists:} replacing a slot directly does not update its stored SHA-256 descriptor. To reprovision a baseline with the programming script, run the VS Code task \textbf{Flash: Reprovision Release Baseline}, or use:

\begin{verbatim}
./scripts/flash.sh --target all --config Release --build --reset-boot-metadata
\end{verbatim}

The option is valid only with \texttt{all} or \texttt{appli-\allowbreak{}a}. It validates the application header/vectors/checksum and build manifest, backs up exactly the two metadata sectors to \texttt{build/\allowbreak{}boot-\allowbreak{}metadata-\allowbreak{}backup.\allowbreak{}*.\allowbreak{}bin}, then programs/verifies the selected firmware and parameter JSON. Only after those operations succeed does it replace \texttt{0x91100000}--\texttt{0x9117FFFF} with erased bytes. Programming verification and an independent full 512 KiB readback must both succeed. No mass erase is used; device identity/calibration are preserved by the ordinary parameter merge. The backup is diagnostic history; restoring it alone does not restore an older application.

This discards update/trial history and causes validated-image adoption again, preferring A, on the next normal boot. Return to normal flash boot mode and reset/power-cycle after programming. The usual \textbf{Build + Program} task retains boot metadata; use DFU for routine A/B updates, or this explicit reprovisioning task when deliberately replacing the slot A baseline over SWD. \texttt{appli-\allowbreak{}b} is rejected with the reset option because first-image adoption prefers A.

Four amber flashes identify stage 4 (slot selection/validation or trial-watchdog startup). Stale image size/hash descriptors after direct SWD replacement can produce this code even though the newly programmed image verifies correctly. The reprovisioning regression reproduces this case with the real boot policy. If the code persists after a verified metadata reset, investigate flash image readback, metadata erase/program operations, or trial-watchdog setup; the code alone does not distinguish those causes.

Parameter programming alone does not require or erase these sectors. \texttt{PARM} JSON build metadata is not rewritten by DFU; current DFU versions/flags come from \texttt{DFUS}.

\section{Diagnosing standalone boot}\label{wiki-firmware-firmware-updates--diagnosing-standalone-boot}

The ROM copies FSBL from external flash to SRAM; loading the FSBL ELF directly in development mode bypasses that flash-loading step and its inherited hardware state. ROM uses \textbf{XSPI1 on XSPIM port 2} for this board\textquotesingle s serial NOR. FSBL must reset XSPI1 as well as XSPIM before reinitializing the flash controller.

FSBL drives the existing I2C LED controller before initializing flash, so these indications do not require debug access with BOOT1 low:

{\def\LTcaptype{none} % do not increment counter
\begin{longtable}[]{@{}
  >{\raggedright\arraybackslash}p{(\linewidth - 2\tabcolsep) * \real{0.4335}}
  >{\raggedright\arraybackslash}p{(\linewidth - 2\tabcolsep) * \real{0.5665}}@{}}
\toprule\noalign{}
\begin{minipage}[b]{\linewidth}\raggedright
\textbf{Indication}
\end{minipage} & \begin{minipage}[b]{\linewidth}\raggedright
\textbf{Meaning}
\end{minipage} \\
\midrule\noalign{}
\endhead
\bottomrule\noalign{}
\endlastfoot
Amber blinking (2 Hz, then 5 Hz) & FSBL reached GPIO/I2C initialization, then is checking/loading an application. \\
Counted amber burst, then 2 s dark & FSBL error; the count identifies flash initialization, metadata/slot validation, or SRAM copy/verification. See the LED code table. \\
Solid magenta & The SRAM application copy passed SHA-256 verification and FSBL is preparing the branch. Persistent magenta without the application\textquotesingle s blue startup indication narrows the failure to handoff or early application startup. \\
Blue blinking & Application reached GPIO/I2C initialization, before its remaining peripherals. \\
Cyan double pulse & Application entry thread started after RTOS/middleware setup. Sensor initialization may still be running. \\
Green breathing & ToF entered ranging state. \\
Counted red, blue, or magenta burst, then 2 s dark & Application error; color and count identify the failing peripheral or RTOS/middleware setup. \\
No FSBL indication & Failure may be in ROM loading, FSBL early clock/GPIO/I2C setup, or LED communication. This alone does not identify the failing component. \\
\end{longtable}
}

See the \hyperref[wiki-firmware-status-led--counted-startup-failure-codes]{complete LED legend and counted fault codes} for the full decoder. Older double-pulse presets repeated their inner pattern infinitely and skipped the pause, so continuous red on those builds could also be an FSBL error; it did not establish successful application entry. The early boot transitions may be brief. The LED controller generates its patterns independently; faults after the last successful LED write can leave that indication displayed. The LED writes are diagnostic and their acknowledgement does not gate boot. A trial watchdog is not refreshed in the FSBL error handler, preserving automatic rollback.

When debug access is available, \texttt{fw\_\allowbreak{}boot\_\allowbreak{}stage} records 1 = LED setup, 2 = flash initialization, 3 = metadata read, 4 = slot validation/selection, 5 = SRAM copy, 6 = copy verification, 7 = handoff preparation. \texttt{fw\_\allowbreak{}boot\_\allowbreak{}result} records the returned updater result if selection/loading fails.

The secure handoff stops SysTick, disables and clears pending NVIC interrupts, cleans/disables D-cache, invalidates/disables I-cache, disables the inherited MPU, and installs VTOR at \texttt{0x34000400}. The final assembly branch clears stack limits and interrupt masks, selects privileged MSP, installs the application\textquotesingle s stack, and branches without returning through the FSBL\textquotesingle s old C stack. Clock settings are inherited by the persistent application; it owns subsequent MPU configuration. See \href{https://www.st.com/resource/en/user_manual/um3234-how-to-proceed-with-boot-rom-on-stm32n6-mcus-stmicroelectronics.pdf}{ST\textquotesingle s boot ROM manual, sections 3.2.2 and 3.4.5} for the differences in ROM cache/MPU state and serial NOR initialization.

To build and install only the updated FSBL (the flashing script also refreshes the parameter record):

\begin{verbatim}
./scripts/build.sh --target fsbl --config Release
./scripts/flash.sh --target fsbl --config Release
\end{verbatim}

Then select normal flash boot (BOOT1 low) and fully power-cycle the board, including its flash supply, for the first test. Record the LED sequence. A CPU reset and a complete flash power cycle are different tests: the running flash driver changes volatile SPI/OPI mode, while ROM expects legacy SPI access.

\section{Persistent state and boot behavior}\label{wiki-firmware-firmware-updates--persistent-state-and-boot-behavior}

Each metadata sector holds a 124-byte record: format/version/length, generation, confirmed slot, pending slot, trial-attempt flag, status flags, last error, and both slot descriptors (size, release identifier and SHA-256). Integers are LE32. A CRC32 covers bytes 0--115; offset 116 holds that CRC and offset 120 the final commit word. The replacement sector is erased/programmed/read back first, and its commit word is programmed \textbf{last}. The previous valid copy remains intact throughout. The newest valid generation is selected with rollover-aware comparison.

\begin{enumerate}
\def\labelenumi{\arabic{enumi}.}
\tightlist
\item
  \texttt{BEGIN} invalidates the inactive descriptor in durable metadata \textbf{before} erasing that slot. The confirmed slot remains untouched.
\item
  \texttt{DATA} accepts only sequential, aligned, bounded chunks with a matching CRC32; each chunk is programmed and read back.
\item
  \texttt{FINISH} rereads the entire image, checks SHA-256, STM32 header/payload checksum, file length, stack bounds/alignment and Thumb entry point. Only then are \texttt{image\_\allowbreak{}uploaded}, \texttt{image\_\allowbreak{}verified} and the pending slot committed.
\item
  FSBL validates the selected image again on every boot. Before jumping to a candidate it commits \texttt{trial = 1}, starts the independent watchdog, loads SRAM, and verifies the copied image\textquotesingle s SHA-256.
\item
  The application confirms the trial only after its dedicated update server and USB receive task have run for five seconds, and its loaded image matches the descriptor. Connections/cables are not required. A sensor fault does not prevent confirmation when the updater itself is available.
\item
  A reset before confirmation (including watchdog reset after an initialization hang or hard fault) makes FSBL reject the candidate, boot the previous image, and retain a fallback warning. A failed confirmation metadata write also keeps the old image recoverable. After confirmation, the DFU worker refreshes IWDG.
\item
  A corrupted confirmed image is rejected using its stored SHA-256; FSBL tries the other recorded valid image and records a warning. If neither is valid, FSBL records \texttt{no\_\allowbreak{}valid\_\allowbreak{}image} when possible and stops rather than jumping to corrupt code; SWD recovery is then required.
\end{enumerate}

The trial window is approximately 32 seconds at nominal 32 kHz LSI; oscillator variation affects the exact timeout. A debugger halt during trial can trigger rollback. Automatic reset/recovery requires normal flash boot mode; with BOOT1 set for development, reset returns to the ROM development flow instead of autonomously restarting FSBL. Five-second confirmation tests updater availability, not every sensor or future application function. Runtime failures after confirmation can reset through IWDG, but do not automatically revoke an already confirmed image unless its stored image fails validation on the next boot.

Image hashes and checksums detect corruption; the existing unauthenticated USB/ TCP transport model is retained. Despite the packaging filename \texttt{-\allowbreak{}signed.\allowbreak{}bin}, the current build uses SigningTool \texttt{-\allowbreak{}nk}; this updater does not authenticate a publisher or verify a cryptographic signing key. Deploy on trusted interfaces.

\section{Wire format}\label{wiki-firmware-firmware-updates--wire-format}

The existing common 16-byte header is used: \texttt{<4sIIHH>} = magic, request ID, payload bytes, version, operation. Requests use \texttt{DFUC}, nonzero ID, version 1. Responses use \texttt{DFUS}, the matching ID/operation, version 1 and 48 payload bytes. One request is outstanding per connection. USB multiplexes \texttt{DFUC} and \texttt{CTRL} without scanning for magic inside accepted payloads; TCP 50003 accepts \texttt{DFUC}.

{\def\LTcaptype{none} % do not increment counter
\begin{longtable}[]{@{}
  >{\raggedright\arraybackslash}p{(\linewidth - 4\tabcolsep) * \real{0.3473}}
  >{\raggedright\arraybackslash}p{(\linewidth - 4\tabcolsep) * \real{0.1750}}
  >{\raggedright\arraybackslash}p{(\linewidth - 4\tabcolsep) * \real{0.4777}}@{}}
\toprule\noalign{}
\begin{minipage}[b]{\linewidth}\raggedright
\textbf{Operation}
\end{minipage} & \begin{minipage}[b]{\linewidth}\raggedright
\textbf{Number}
\end{minipage} & \begin{minipage}[b]{\linewidth}\raggedright
\textbf{Payload}
\end{minipage} \\
\midrule\noalign{}
\endhead
\bottomrule\noalign{}
\endlastfoot
Status & 1 & Empty \\
Begin & 2 & \texttt{<II32s>}: image size, release identifier, SHA-256 \\
Data & 3 & \texttt{<II>}: byte offset, CRC32 of data; followed by 4--512 data bytes \\
Finish & 4 & Empty \\
Abort in-progress upload & 5 & Empty \\
Reset into verified pending image & 6 & Empty \\
Clear sticky warnings & 7 & Empty \\
\end{longtable}
}

CRC32 uses the standard reflected polynomial \texttt{0xEDB88320}, compatible with Python \texttt{zlib.\allowbreak{}crc32}. Image/chunk sizes are multiples of four. The board picks the slot; no command accepts a physical flash address. Reset is scheduled 500 ms after successfully sending its reply. DEV\_MODE builds accept status/legacy commands but reject update mutations and never confirm a flash trial.

\texttt{DFUS} payload is \texttt{<12I>}:

{\def\LTcaptype{none} % do not increment counter
\begin{longtable}[]{@{}
  >{\raggedright\arraybackslash}p{(\linewidth - 2\tabcolsep) * \real{0.3438}}
  >{\raggedright\arraybackslash}p{(\linewidth - 2\tabcolsep) * \real{0.6562}}@{}}
\toprule\noalign{}
\begin{minipage}[b]{\linewidth}\raggedright
\textbf{Word}
\end{minipage} & \begin{minipage}[b]{\linewidth}\raggedright
\textbf{Field}
\end{minipage} \\
\midrule\noalign{}
\endhead
\bottomrule\noalign{}
\endlastfoot
0 & Result \\
1 & Persistent flags \\
2 & Running slot \\
3 & Confirmed slot \\
4 & Pending slot \\
5 & Received bytes in current upload \\
6 & Total bytes in current upload \\
7 & Last persistent error \\
8 & Metadata generation \\
9 & Slot A release identifier \\
10 & Slot B release identifier \\
11 & Maximum chunk data size (512) \\
\end{longtable}
}

Results: 0 OK, 1 bad request, 2 busy/other owner, 3 flash I/O failure, 4 invalid image, 5 SHA/CRC mismatch, 6 unexpected offset/length, 7 not ready, 8 previous trial failed. The status result also reports startup validation or missing A/B metadata, so a legacy FSBL cannot silently accept staged updates.

{\def\LTcaptype{none} % do not increment counter
\begin{longtable}[]{@{}
  >{\raggedright\arraybackslash}p{(\linewidth - 2\tabcolsep) * \real{0.3427}}
  >{\raggedright\arraybackslash}p{(\linewidth - 2\tabcolsep) * \real{0.6573}}@{}}
\toprule\noalign{}
\begin{minipage}[b]{\linewidth}\raggedright
\textbf{Flag bit}
\end{minipage} & \begin{minipage}[b]{\linewidth}\raggedright
\textbf{Meaning}
\end{minipage} \\
\midrule\noalign{}
\endhead
\bottomrule\noalign{}
\endlastfoot
0 & Upload in progress \\
1 & Complete image uploaded \\
2 & Stored image verified \\
3 & Candidate boot confirmed \\
4 & Sticky fallback warning \\
5 & Sticky upload failure \\
6 & Sticky no-valid-image warning \\
7 & Metadata initialized/recovered by image adoption \\
8 & Sticky aborted upload \\
\end{longtable}
}

\section{Validation}\label{wiki-firmware-firmware-updates--validation}

Run host tests against the actual portable C policy/parser plus the Python client:

\begin{verbatim}
python3 -m unittest discover -s test_scripts -p test_firmware_update.py -v
python3 -m unittest discover -s test_scripts -p test_viewer_firmware_update.py -v
\end{verbatim}

Tests model NOR one-to-zero programming, sequential updates in both directions, header/vector/checksum/SHA failures, additive-checksum collisions, wrong offsets, CRC failures, interrupted transfers, unconfirmed boots, corrupted confirmed/ pending images, generation rollover, and interrupted metadata writes/erases. ASan/UBSan cover the host C harness. Cross-build both Appli and FSBL. Viewer tests cover progress, exclusive USB ownership and automatic reconnection, staging, cancellation, lost acknowledgements, confirmation, and fallback. With a working display they also exercise the actual tab and window-close behavior; those GUI tests are skipped when a display cannot be opened.

Hardware acceptance remains necessary:

\begin{itemize}
\tightlist
\item
  Provision and boot the new baseline once; inspect initial status.
\item
  Upload A → B over USB, reset, then verify B confirmation and stored flags.
\item
  Upload B → A over TCP, reset, then verify A confirmation.
\item
  Interrupt transfers and power during slot erase/program and metadata commits; the confirmed image must still boot and accept a new update.
\item
  Upload a structurally valid test build that deliberately hangs before updater startup; verify watchdog reset, previous-image boot, and fallback warning.
\item
  Corrupt the pending image before reset and a confirmed image with an intact backup; verify rejection and fallback without executing the corrupted image.
\item
  Disconnect/reconnect USB and Ethernet mid-upload and check ownership, abort flags, telemetry recovery and a successful subsequent upload.
\end{itemize}

\chapter{Magnecko Sense wire protocol}\label{wiki-protocols-wire-protocol--magnecko-sense-wire-protocol}

{\small\color{MGSMuted}Source: \nolinkurl{wiki/protocols/wire-protocol.md}\par}\medskip

\section{Connections and lifecycle}\label{wiki-protocols-wire-protocol--connections-and-lifecycle}

{\def\LTcaptype{none} % do not increment counter
\begin{longtable}[]{@{}
  >{\raggedright\arraybackslash}p{(\linewidth - 4\tabcolsep) * \real{0.2160}}
  >{\raggedright\arraybackslash}p{(\linewidth - 4\tabcolsep) * \real{0.3510}}
  >{\raggedright\arraybackslash}p{(\linewidth - 4\tabcolsep) * \real{0.4330}}@{}}
\toprule\noalign{}
\begin{minipage}[b]{\linewidth}\raggedright
\textbf{Transport}
\end{minipage} & \begin{minipage}[b]{\linewidth}\raggedright
\textbf{Endpoint}
\end{minipage} & \begin{minipage}[b]{\linewidth}\raggedright
\textbf{Direction}
\end{minipage} \\
\midrule\noalign{}
\endhead
\bottomrule\noalign{}
\endlastfoot
Ethernet & TCP port \textbf{50002} (default board IP \texttt{192.\allowbreak{}168.\allowbreak{}8.\allowbreak{}20}) & Framed commands to board; interleaved telemetry and replies from board \\
USB & CDC ACM, typically \texttt{/\allowbreak{}dev/\allowbreak{}ttyACM0} & Identical framing in each direction; assert DTR \\
Legacy Ethernet & TCP port 50000 & Echo service; does \textbf{not} accept the commands below \\
\end{longtable}
}

USB baud rate is ignored (host uses 115200). One TCP stream client and one USB serial owner are supported; both transports can operate simultaneously. Open a connection and read packets; no handshake is required. Ranging starts at boot with the compiled \texttt{CONF\_\allowbreak{}USECASE} profile. \texttt{GET\_\allowbreak{}CONFIG} or periodic \texttt{TOFS} reports reveal the actual running configuration; do not assume particular defaults.

Ethernet supports cable insertion after boot and unplug/replug during operation. The network thread checks the PHY every 200 ms and explicitly enables/disables the MAC/DMA driver; PHY link alone does not mean the interface is ready. A cable cycle invalidates old TCP sessions, including incomplete handshakes, so clients must reconnect. The static IPv4 configuration is retained. USB and sensor acquisition continue during Ethernet outages.

\texttt{STOP} stops sensor ranging and the CSI receiver, affecting depth output on \textbf{both} transports. IMU, parameters, CPU profiling and sensor status continue. \texttt{START} restarts stopped ranging. Disconnecting does not change ranging state. Settings and start/stop state are volatile. A board reset restores compiled defaults and starts ranging. \texttt{RESET} (operation 5) resets only the ToF sensor, restores compiled defaults and leaves ranging \textbf{stopped}; press Start afterward. IMU, USB/Ethernet connections and CPU profiling continue during sensor reset. \texttt{POWER\_\allowbreak{}DOWN} (operation 6) stops capture and holds XSHUT low. State becomes \textbf{powered down}; this is sensor shutdown with supply rails still connected. Use \texttt{RESET}, then \texttt{START}, to range again. Start and Set while powered down return result 7. Stop and Power down are idempotent while off. Shutdown is also accepted when a fault is latched; it preserves the error until Reset. The flash parameter JSON is identity/build metadata, separate from these runtime sensor settings.

All I3C sensor operations run in the ToF thread. A complete depth packet already owned by the transport is finished before a Stop acknowledgement, sensor reset, power-down or configuration change. Successful Stop acknowledgement therefore follows the last depth packet on that connection, until a later Start (possibly from the other transport). Faults remain latched until an explicit \texttt{RESET} or board reset; there is no automatic reset or retry of faulted emission. \texttt{RESET} stops CSI capture, clears its pending frame event, toggles the sensor\textquotesingle s XSHUT pin, reassigns its I3C address, reloads the sensor firmware patch and applies the compiled profile. Success clears the fault and acknowledges stopped state with default settings. Failure returns Sensor operation failed and keeps a fault latched. Driver validation and sensor protection remain active.

\section{Common 16-byte header}\label{wiki-protocols-wire-protocol--common-16-byte-header}

All integers are \textbf{little-endian}. Floats are IEEE-754 binary32, little-endian. There is no structure padding, CRC, terminator or alignment between packets. Python header format: \texttt{struct.\allowbreak{}Struct('<4sIIHH')}.

{\def\LTcaptype{none} % do not increment counter
\begin{longtable}[]{@{}
  >{\raggedright\arraybackslash}p{(\linewidth - 6\tabcolsep) * \real{0.2059}}
  >{\raggedright\arraybackslash}p{(\linewidth - 6\tabcolsep) * \real{0.2256}}
  >{\raggedright\arraybackslash}p{(\linewidth - 6\tabcolsep) * \real{0.2059}}
  >{\raggedright\arraybackslash}p{(\linewidth - 6\tabcolsep) * \real{0.3626}}@{}}
\toprule\noalign{}
\begin{minipage}[b]{\linewidth}\raggedright
\textbf{Byte offset}
\end{minipage} & \begin{minipage}[b]{\linewidth}\raggedright
\textbf{Size}
\end{minipage} & \begin{minipage}[b]{\linewidth}\raggedright
\textbf{Type}
\end{minipage} & \begin{minipage}[b]{\linewidth}\raggedright
\textbf{Meaning}
\end{minipage} \\
\midrule\noalign{}
\endhead
\bottomrule\noalign{}
\endlastfoot
0 & 4 & ASCII bytes & Magic identifying packet type \\
4 & 4 & uint32 & Sequence number, or command request ID \\
8 & 4 & uint32 & Payload byte length, excluding the header \\
12 & 2 & uint16 & Type-specific field A \\
14 & 2 & uint16 & Type-specific field B \\
16 & payload length & bytes & Payload \\
\end{longtable}
}

TCP reads and USB transfers may split or combine packets arbitrarily. Accumulate bytes, validate the magic/header/length, then consume exactly \texttt{16 + length} bytes. Never search for magic inside a valid packet\textquotesingle s payload. Preserve up to three trailing bytes while searching for a new magic after an invalid header. Clear partial packets at every reconnect. Sequence arithmetic wraps modulo 2³².

{\def\LTcaptype{none} % do not increment counter
\begin{longtable}[]{@{}
  >{\raggedright\arraybackslash}p{(\linewidth - 8\tabcolsep) * \real{0.1307}}
  >{\raggedright\arraybackslash}p{(\linewidth - 8\tabcolsep) * \real{0.1583}}
  >{\raggedright\arraybackslash}p{(\linewidth - 8\tabcolsep) * \real{0.2535}}
  >{\raggedright\arraybackslash}p{(\linewidth - 8\tabcolsep) * \real{0.2465}}
  >{\raggedright\arraybackslash}p{(\linewidth - 8\tabcolsep) * \real{0.2111}}@{}}
\toprule\noalign{}
\begin{minipage}[b]{\linewidth}\raggedright
\textbf{Magic}
\end{minipage} & \begin{minipage}[b]{\linewidth}\raggedright
\textbf{Direction}
\end{minipage} & \begin{minipage}[b]{\linewidth}\raggedright
\textbf{Header A}
\end{minipage} & \begin{minipage}[b]{\linewidth}\raggedright
\textbf{Header B}
\end{minipage} & \begin{minipage}[b]{\linewidth}\raggedright
\textbf{Payload bytes}
\end{minipage} \\
\midrule\noalign{}
\endhead
\bottomrule\noalign{}
\endlastfoot
\texttt{CTRL} & Host → board & Format version \textbf{1 or 2} & Operation & 0 or 12 \\
\texttt{TOFS} & Board → host & Format version \textbf{2} (legacy 1 also decoded by PC) & Replied-to operation; 0 for periodic status & 36 (legacy 32) \\
\texttt{VL9F} & Board → host & CSI width & CSI height & width × height, at most 14900 \\
\texttt{IMUF} & Board → host & \textbf{10} float channels & Format version \textbf{1} & 44 \\
\texttt{PARM} & Board → host & Format version \textbf{1} & Reserved \textbf{0} & 1--16384 UTF-8 bytes \\
\texttt{PROF} & Board → host & Format version \textbf{1} & Thread count, 0--64 & 44 + 48 × thread count \\
\end{longtable}
}

Camera, IMU and profiler sequence counters are independent, start at zero and continue across connection and ToF start/stop cycles. \texttt{PARM} uses sequence zero. \texttt{TOFS} replies use the request\textquotesingle s ID; periodic status uses ID zero. Do not include \texttt{TOFS}/\texttt{PARM} in sensor sequence-gap calculations.

\section{\texorpdfstring{Commands: \texttt{CTRL}}{Commands: }}\label{wiki-protocols-wire-protocol--commands-ctrl}

Use a \textbf{nonzero uint32 request ID}, field A = 1 or 2, field B = operation. The monitor sends Set as v2 and empty commands as v1. Firmware accepts both. Update firmware and monitor together to use the new modes/shutdown; old firmware does not recognize v2 Set commands. Allow one outstanding command per connection. Match replies by request ID \textbf{and} operation. The Device Monitor waits up to eight seconds; a timeout/disconnect leaves command outcome uncertain. Read current state before retrying a mutating command. Do not blindly resend commands on reconnect.

{\def\LTcaptype{none} % do not increment counter
\begin{longtable}[]{@{}
  >{\raggedright\arraybackslash}p{(\linewidth - 6\tabcolsep) * \real{0.1675}}
  >{\raggedright\arraybackslash}p{(\linewidth - 6\tabcolsep) * \real{0.1783}}
  >{\raggedright\arraybackslash}p{(\linewidth - 6\tabcolsep) * \real{0.1933}}
  >{\raggedright\arraybackslash}p{(\linewidth - 6\tabcolsep) * \real{0.4608}}@{}}
\toprule\noalign{}
\begin{minipage}[b]{\linewidth}\raggedright
\textbf{Operation}
\end{minipage} & \begin{minipage}[b]{\linewidth}\raggedright
\textbf{Name}
\end{minipage} & \begin{minipage}[b]{\linewidth}\raggedright
\textbf{Payload}
\end{minipage} & \begin{minipage}[b]{\linewidth}\raggedright
\textbf{Behavior}
\end{minipage} \\
\midrule\noalign{}
\endhead
\bottomrule\noalign{}
\endlastfoot
1 & \texttt{GET\_\allowbreak{}CONFIG} & Empty & Return current settings/state \\
2 & \texttt{START} & Empty & Start stopped ranging; already streaming succeeds without restarting \\
3 & \texttt{STOP} & Empty & Stop ranging/capture; already stopped succeeds \\
4 & \texttt{SET\_\allowbreak{}CONFIG} & 12 bytes below & Apply settings while stopped; remain stopped \\
5 & \texttt{RESET} & Empty & Reset/reinitialize the sensor from streaming, stopped, powered-down or faulted state; restore compiled defaults and remain stopped \\
6 & \texttt{POWER\_\allowbreak{}DOWN} & Empty & Stop CSI and assert XSHUT; leave sensor off until Reset \\
\end{longtable}
}

\texttt{SET\_\allowbreak{}CONFIG} payload format: \texttt{<IHHBBH}.

{\def\LTcaptype{none} % do not increment counter
\begin{longtable}[]{@{}
  >{\raggedright\arraybackslash}p{(\linewidth - 6\tabcolsep) * \real{0.2726}}
  >{\raggedright\arraybackslash}p{(\linewidth - 6\tabcolsep) * \real{0.1476}}
  >{\raggedright\arraybackslash}p{(\linewidth - 6\tabcolsep) * \real{0.2063}}
  >{\raggedright\arraybackslash}p{(\linewidth - 6\tabcolsep) * \real{0.3735}}@{}}
\toprule\noalign{}
\begin{minipage}[b]{\linewidth}\raggedright
\textbf{Payload offset (absolute offset)}
\end{minipage} & \begin{minipage}[b]{\linewidth}\raggedright
\textbf{Size}
\end{minipage} & \begin{minipage}[b]{\linewidth}\raggedright
\textbf{Field}
\end{minipage} & \begin{minipage}[b]{\linewidth}\raggedright
\textbf{Valid values}
\end{minipage} \\
\midrule\noalign{}
\endhead
\bottomrule\noalign{}
\endlastfoot
0 (16) & 4 & Frame period, µs & 10000--1000000, inclusive \\
4 (20) & 2 & Exposure, ms & Integer 1--30; exposure × 1000 ≤ frame period \\
6 (22) & 2 & Binning & 2, 4, 6, 8, 12, 24 \\
8 (24) & 1 & Power mode & 0 regular, 1 low, 2 ultra low \\
9 (25) & 1 & Range context & 0 short, 1 long \\
10 (26) & 2 & Flags & v1: \textbf{0} (DSS enabled); v2 bit 0: \textbf{disable DSS}, all other bits zero \\
\end{longtable}
}

The datasheet\textquotesingle s ranging modes correspond to these settings:

{\def\LTcaptype{none} % do not increment counter
\begin{longtable}[]{@{}
  >{\raggedright\arraybackslash}p{(\linewidth - 6\tabcolsep) * \real{0.3140}}
  >{\raggedright\arraybackslash}p{(\linewidth - 6\tabcolsep) * \real{0.2112}}
  >{\raggedright\arraybackslash}p{(\linewidth - 6\tabcolsep) * \real{0.2526}}
  >{\raggedright\arraybackslash}p{(\linewidth - 6\tabcolsep) * \real{0.2222}}@{}}
\toprule\noalign{}
\begin{minipage}[b]{\linewidth}\raggedright
\textbf{Ranging mode}
\end{minipage} & \begin{minipage}[b]{\linewidth}\raggedright
\textbf{Context}
\end{minipage} & \begin{minipage}[b]{\linewidth}\raggedright
\textbf{DSS enabled}
\end{minipage} & \begin{minipage}[b]{\linewidth}\raggedright
\textbf{v2 flags}
\end{minipage} \\
\midrule\noalign{}
\endhead
\bottomrule\noalign{}
\endlastfoot
Precision (no DSS) & 0 short & false & 1 \\
Precision & 0 short & true & 0 \\
Ambient & 1 long & true & 0 \\
\end{longtable}
}

Long context with DSS disabled is rejected by the control protocol. Gaming now uses Precision (no DSS), regular power, binning 2, 10000 µs period and the user\textquotesingle s 1 ms exposure. Applying binning first restores driver DSS defaults; the explicit DSS selection is written and read back afterwards.

Sync mode stays autonomous; CSI stays RAW8, VC0, one lane, 1 Gbit/s. The firmware updates both the sensor\textquotesingle s ranging settings and its CSI frame geometry when binning changes. Passing these host-side bounds does not guarantee that the sensor accepts every combination; check the acknowledgement and state.

Command parser storage is bounded to 28 bytes, with eight queued commands and eight queued replies. Invalid magic, version, or payload length above 12 is resynchronized without an acknowledgement. Recognized frames with an invalid operation, payload layout, reserved value or configuration receive an error. Queue saturation produces a best-effort Busy reply; a flooding client can also lose replies. A transport session change discards partial commands and stale replies; already accepted commands may still execute. Do not pipeline requests.

\subsection{Exact command examples}\label{wiki-protocols-wire-protocol--exact-command-examples}

Stop, request ID \texttt{0x12345678} (16 bytes):

\begin{verbatim}
43 54 52 4c  78 56 34 12  00 00 00 00  01 00 03 00
C  T  R  L  request ID    payload = 0  v1    STOP
\end{verbatim}

Reset ToF, request ID 1 (16 bytes, no payload):

\begin{verbatim}
43 54 52 4c  01 00 00 00  00 00 00 00  01 00 05 00
\end{verbatim}

Wait for a \texttt{TOFS} response with request ID 1, operation 5, result 0 and state 1 (stopped). Reset discards runtime settings; the reply contains the restored configuration. Send Start separately to resume ranging.

Set 20 ms frame period, 2 ms exposure, binning 4, low power, short context, request ID 7 (28 bytes, v2, DSS enabled):

\begin{verbatim}
43 54 52 4c  07 00 00 00  0c 00 00 00  02 00 04 00
20 4e 00 00  02 00 04 00  01 00 00 00
\end{verbatim}

Python (on an already-open USB serial connection or TCP stream socket):

\begin{verbatim}
from Sensor_Control import build_command, STOP, SET_CONFIG, START

stop = build_command(1, STOP)
settings = build_command(2, SET_CONFIG, dict(
    period_us=20000, exposure_ms=2, binning=4, power=1, context=0, dss_enabled=True))
start = build_command(3, START)
# Send stop; wait for matching successful TOFS.
# Send settings; wait for matching successful TOFS.
# Send start; wait for matching successful TOFS.
# TCP: connection.sendall(command); USB: connection.write(command).
\end{verbatim}

\section{\texorpdfstring{Sensor status and acknowledgements: \texttt{TOFS}}{Sensor status and acknowledgements: }}\label{wiki-protocols-wire-protocol--sensor-status-and-acknowledgements-tofs}

Every command reply contains a complete state snapshot. Replies go only to the requesting transport/session, interleaved with telemetry by the single transport writer. Unsolicited snapshots are sent to both links about once per second, including while stopped, powered down or faulted. Current v2 payload format: \texttt{<IIiIHHBBBBHHIHH} (36 bytes). The first 32 bytes retain the v1 layout. The PC treats v1 status as DSS enabled.

{\def\LTcaptype{none} % do not increment counter
\begin{longtable}[]{@{}
  >{\raggedright\arraybackslash}p{(\linewidth - 6\tabcolsep) * \real{0.2022}}
  >{\raggedright\arraybackslash}p{(\linewidth - 6\tabcolsep) * \real{0.1505}}
  >{\raggedright\arraybackslash}p{(\linewidth - 6\tabcolsep) * \real{0.1547}}
  >{\raggedright\arraybackslash}p{(\linewidth - 6\tabcolsep) * \real{0.4926}}@{}}
\toprule\noalign{}
\begin{minipage}[b]{\linewidth}\raggedright
\textbf{Payload offset}
\end{minipage} & \begin{minipage}[b]{\linewidth}\raggedright
\textbf{Size}
\end{minipage} & \begin{minipage}[b]{\linewidth}\raggedright
\textbf{Type}
\end{minipage} & \begin{minipage}[b]{\linewidth}\raggedright
\textbf{Field}
\end{minipage} \\
\midrule\noalign{}
\endhead
\bottomrule\noalign{}
\endlastfoot
0 & 4 & uint32 & MCU uptime timestamp, ms \\
4 & 4 & uint32 & Configuration/state revision; increments on state or configuration changes \\
8 & 4 & int32 & Latched error: 0 healthy; sensor firmware code or negative application/driver error (\texttt{-\allowbreak{}1} for generic failures) \\
12 & 4 & uint32 & Frame period, µs \\
16 & 2 & uint16 & Exposure, ms \\
18 & 2 & uint16 & Binning \\
20 & 1 & uint8 & Power mode, as above \\
21 & 1 & uint8 & Range context, as above \\
22 & 1 & uint8 & State: 0 initializing, 1 stopped, 2 streaming, 3 fault, 4 powered down \\
23 & 1 & uint8 & Command result, below; periodic status uses 0 \\
24 & 2 & uint16 & CSI width in bytes \\
26 & 2 & uint16 & CSI height in rows, including status row \\
28 & 4 & uint32 & Frames captured since STM32 boot, modulo 2³²; preserved across sensor resets \\
32 & 2 & uint16 & v2 flags: bit 0 = DSS disabled; all other bits zero \\
34 & 2 & uint16 & v2 reserved: zero \\
\end{longtable}
}

Settings/dimensions can be zero during initialization or after initialization failure. A configuration failure latches a fault; do not assume partially written hardware settings were rolled back. The last successfully applied settings remain in the snapshot. Result 0 in a periodic report is not a health indication; inspect state/error. While stopped, faulted or powered down, the monitor clears the depth display.

{\def\LTcaptype{none} % do not increment counter
\begin{longtable}[]{@{}
  >{\raggedright\arraybackslash}p{(\linewidth - 2\tabcolsep) * \real{0.3069}}
  >{\raggedright\arraybackslash}p{(\linewidth - 2\tabcolsep) * \real{0.6931}}@{}}
\toprule\noalign{}
\begin{minipage}[b]{\linewidth}\raggedright
\textbf{Result}
\end{minipage} & \begin{minipage}[b]{\linewidth}\raggedright
\textbf{Meaning}
\end{minipage} \\
\midrule\noalign{}
\endhead
\bottomrule\noalign{}
\endlastfoot
0 & Success \\
1 & Unknown command \\
2 & Invalid request/configuration \\
3 & Busy / command queue full \\
4 & Sensor/capture operation failed; inspect latched state/error \\
5 & Configuration requires stopped ranging \\
6 & Sensor fault is latched; explicit \texttt{RESET} or board reset required \\
7 & Sensor is powered down; use \texttt{RESET} before Start/Set \\
\end{longtable}
}

\section{\texorpdfstring{Depth frames: \texttt{VL9F}}{Depth frames: }}\label{wiki-protocols-wire-protocol--depth-frames-vl9f}

The payload is a tightly packed CSI RAW8 dump, including sensor padding and the final \textbf{100-byte metadata row}. Header width/height describe transport bytes, not the number of depth zones. Sequence is the board\textquotesingle s captured-frame counter, independent of the sensor metadata\textquotesingle s own frame counter, which can restart on sensor reset.

{\def\LTcaptype{none} % do not increment counter
\begin{longtable}[]{@{}
  >{\raggedright\arraybackslash}p{(\linewidth - 8\tabcolsep) * \real{0.1552}}
  >{\raggedright\arraybackslash}p{(\linewidth - 8\tabcolsep) * \real{0.1874}}
  >{\raggedright\arraybackslash}p{(\linewidth - 8\tabcolsep) * \real{0.2685}}
  >{\raggedright\arraybackslash}p{(\linewidth - 8\tabcolsep) * \real{0.2015}}
  >{\raggedright\arraybackslash}p{(\linewidth - 8\tabcolsep) * \real{0.1874}}@{}}
\toprule\noalign{}
\begin{minipage}[b]{\linewidth}\raggedright
\textbf{Binning}
\end{minipage} & \begin{minipage}[b]{\linewidth}\raggedright
\textbf{Raw zone grid}
\end{minipage} & \begin{minipage}[b]{\linewidth}\raggedright
\textbf{Displayed grid after default crop}
\end{minipage} & \begin{minipage}[b]{\linewidth}\raggedright
\textbf{CSI bytes × rows}
\end{minipage} & \begin{minipage}[b]{\linewidth}\raggedright
\textbf{Payload bytes}
\end{minipage} \\
\midrule\noalign{}
\endhead
\bottomrule\noalign{}
\endlastfoot
2 & 54×42 & 54×42 & 100×149 & 14900 \\
4 & 24×24 & 24×20 & 100×39 & 3900 \\
6 & 18×14 & 18×14 & 100×18 & 1800 \\
8 & 12×10 & 12×10 & 100×9 & 900 \\
12 & 8×8 & 8×6 & 100×6 & 600 \\
24 & 4×4 & 4×4 & 100×3 & 300 \\
\end{longtable}
}

For \texttt{N = raw\_\allowbreak{}width × raw\_\allowbreak{}height}, payload planes are:

{\def\LTcaptype{none} % do not increment counter
\begin{longtable}[]{@{}
  >{\raggedright\arraybackslash}p{(\linewidth - 4\tabcolsep) * \real{0.2360}}
  >{\raggedright\arraybackslash}p{(\linewidth - 4\tabcolsep) * \real{0.3001}}
  >{\raggedright\arraybackslash}p{(\linewidth - 4\tabcolsep) * \real{0.4639}}@{}}
\toprule\noalign{}
\begin{minipage}[b]{\linewidth}\raggedright
\textbf{Payload offset}
\end{minipage} & \begin{minipage}[b]{\linewidth}\raggedright
\textbf{Length}
\end{minipage} & \begin{minipage}[b]{\linewidth}\raggedright
\textbf{Content}
\end{minipage} \\
\midrule\noalign{}
\endhead
\bottomrule\noalign{}
\endlastfoot
0 & 2N & Little-endian uint16 radial distances: bits 0--14 = mm, bit 15 = ranging-step selector \\
2N & 2N & uint16 amplitude plane \\
4N & 2N & uint16 ambient plane \\
6N & D = N/2 if DSS enabled; otherwise 0 & Packed DSS lookup data (absent with DSS disabled) \\
6N + D & Remaining bytes before final row & CSI padding; fixed geometry above supports either DSS setting \\
payload\_length − 100 & 100 & Sensor metadata \\
\end{longtable}
}

Zero distance means a missing return in the monitor. Bit 15 is not an invalidity flag. Decode/crop from metadata rather than guessing from the CSI row count. Within the final metadata row (offsets relative to that row):

{\def\LTcaptype{none} % do not increment counter
\begin{longtable}[]{@{}
  >{\raggedright\arraybackslash}p{(\linewidth - 4\tabcolsep) * \real{0.1787}}
  >{\raggedright\arraybackslash}p{(\linewidth - 4\tabcolsep) * \real{0.1787}}
  >{\raggedright\arraybackslash}p{(\linewidth - 4\tabcolsep) * \real{0.6425}}@{}}
\toprule\noalign{}
\begin{minipage}[b]{\linewidth}\raggedright
\textbf{Offset}
\end{minipage} & \begin{minipage}[b]{\linewidth}\raggedright
\textbf{Type}
\end{minipage} & \begin{minipage}[b]{\linewidth}\raggedright
\textbf{Meaning}
\end{minipage} \\
\midrule\noalign{}
\endhead
\bottomrule\noalign{}
\endlastfoot
0 & uint32 & Sensor frame counter \\
4 & int16 & Sensor die temperature, °C \\
6 & uint16 & Laser-driver temperature register (raw; no host conversion) \\
58 & uint16 & Dynamic flags; DSS = \texttt{(word >> 4) \& 3} (0 disabled, 1 long, 2 short), binning = \texttt{(word >> 6) \& 0x1f}, context = \texttt{(word >> 11) \& 1} \\
60 & uint16 & Sensor error code \\
62 & uint8 & Sensor error status \\
68 & uint32 & Reported sensor frame period, µs (does not prove delivered FPS) \\
72 & uint32 & Crop: width bits 0--5, height 6--11, x offset 12--17, y offset 18--23, enabled bit 24 \\
\end{longtable}
}

\texttt{vl53l9\_\allowbreak{}utils.\allowbreak{}h} defines the remaining sensor-specific metadata fields. Raw radial ranges are not ST\textquotesingle s calibrated/confidence-filtered or perpendicular-depth output.

\section{\texorpdfstring{IMU: \texttt{IMUF}}{IMU: }}\label{wiki-protocols-wire-protocol--imu-imuf}

Payload \texttt{<I10f>} is 44 bytes; the complete packet is 60 bytes.

{\def\LTcaptype{none} % do not increment counter
\begin{longtable}[]{@{}
  >{\raggedright\arraybackslash}p{(\linewidth - 4\tabcolsep) * \real{0.2715}}
  >{\raggedright\arraybackslash}p{(\linewidth - 4\tabcolsep) * \real{0.2452}}
  >{\raggedright\arraybackslash}p{(\linewidth - 4\tabcolsep) * \real{0.4833}}@{}}
\toprule\noalign{}
\begin{minipage}[b]{\linewidth}\raggedright
\textbf{Payload offset}
\end{minipage} & \begin{minipage}[b]{\linewidth}\raggedright
\textbf{Type}
\end{minipage} & \begin{minipage}[b]{\linewidth}\raggedright
\textbf{Meaning}
\end{minipage} \\
\midrule\noalign{}
\endhead
\bottomrule\noalign{}
\endlastfoot
0 & uint32 & MCU acquisition timestamp, ms \\
4 & 3 × float32 & Acceleration X/Y/Z, mg \\
16 & 3 × float32 & Gyroscope X/Y/Z, mdps (divide by 1000 for °/s) \\
28 & 4 × float32 & SFLP game quaternion X/Y/Z/W \\
\end{longtable}
}

Sampling is published at 100 Hz, backed by 120 Hz sensor/SFLP updates. A bounded 16-packet queue drops samples when full; sequence gaps expose losses. Timestamps wrap modulo 2³². SFLP has no compass-heading reference.

\section{\texorpdfstring{Device/build parameters: \texttt{PARM}}{Device/build parameters: }}\label{wiki-protocols-wire-protocol--devicebuild-parameters-parm}

Payload is UTF-8 JSON, excluding the flash\textquotesingle s NUL terminator/padding. Repeated about every two seconds for late opens/reconnects. Parameter \texttt{schema\_\allowbreak{}version} is independent of wire format version 1. The board forwards stored metadata; it does not apply this JSON to the runtime sensor controls. Unreadable or missing parameters use a JSON object containing \texttt{parameter\_\allowbreak{}error}.

\texttt{parameter\_\allowbreak{}build.\allowbreak{}built\_\allowbreak{}at\_\allowbreak{}utc} records when the parameter artifact was generated. Each Appli/FSBL flash also refreshes the selected slot\textquotesingle s \texttt{build.\allowbreak{}<slot>.\allowbreak{}built\_\allowbreak{}at\_\allowbreak{}utc} and \texttt{programmed\_\allowbreak{}at\_\allowbreak{}utc}; \texttt{image\_\allowbreak{}built\_\allowbreak{}at\_\allowbreak{}utc} retains the manifest\textquotesingle s original image build date. Unflashed slots and all firmware dates during a parameter-only update stay unchanged. The application reads the JSON once at startup, so a new parameter image appears in \texttt{PARM} after a board restart.

See \hyperref[wiki-firmware-parameters--device-parameters-and-build-metadata]{parameter management} for fields, flash layout and the metadata update workflow.

\section{\texorpdfstring{CPU profiling: \texttt{PROF}}{CPU profiling: }}\label{wiki-protocols-wire-protocol--cpu-profiling-prof}

Reports arrive about once per second, with actual intervals lengthened by transport backpressure. Payload header \texttt{<IIQQQQHH>} is 44 bytes:

{\def\LTcaptype{none} % do not increment counter
\begin{longtable}[]{@{}
  >{\raggedright\arraybackslash}p{(\linewidth - 4\tabcolsep) * \real{0.2650}}
  >{\raggedright\arraybackslash}p{(\linewidth - 4\tabcolsep) * \real{0.1845}}
  >{\raggedright\arraybackslash}p{(\linewidth - 4\tabcolsep) * \real{0.5505}}@{}}
\toprule\noalign{}
\begin{minipage}[b]{\linewidth}\raggedright
\textbf{Payload offset}
\end{minipage} & \begin{minipage}[b]{\linewidth}\raggedright
\textbf{Size}
\end{minipage} & \begin{minipage}[b]{\linewidth}\raggedright
\textbf{Field}
\end{minipage} \\
\midrule\noalign{}
\endhead
\bottomrule\noalign{}
\endlastfoot
0 & 4 & MCU timestamp, ms \\
4 & 4 & CPU clock frequency, Hz \\
8 & 8 & Interval cycles \\
16 & 8 & Idle cycles \\
24 & 8 & Interrupt cycles (nested interrupts counted once) \\
32 & 8 & Scheduler cycles \\
40 & 2 & Thread count (must match common header field B) \\
42 & 2 & Flags: bit 0 = thread list truncated; other bits reserved zero \\
\end{longtable}
}

Each following thread record \texttt{<IIQ32s>} is 48 bytes:

{\def\LTcaptype{none} % do not increment counter
\begin{longtable}[]{@{}
  >{\raggedright\arraybackslash}p{(\linewidth - 4\tabcolsep) * \real{0.2886}}
  >{\raggedright\arraybackslash}p{(\linewidth - 4\tabcolsep) * \real{0.2025}}
  >{\raggedright\arraybackslash}p{(\linewidth - 4\tabcolsep) * \real{0.5090}}@{}}
\toprule\noalign{}
\begin{minipage}[b]{\linewidth}\raggedright
\textbf{Record offset}
\end{minipage} & \begin{minipage}[b]{\linewidth}\raggedright
\textbf{Size}
\end{minipage} & \begin{minipage}[b]{\linewidth}\raggedright
\textbf{Field}
\end{minipage} \\
\midrule\noalign{}
\endhead
\bottomrule\noalign{}
\endlastfoot
0 & 4 & Thread ID (control-block address) \\
4 & 4 & ThreadX priority \\
8 & 8 & Thread CPU cycles in this interval \\
16 & 32 & NUL-terminated name, up to 31 UTF-8 bytes \\
\end{longtable}
}

\texttt{milliseconds = cycles × 1000 /\allowbreak{} clock\_\allowbreak{}hz}, \texttt{CPU share = cycles × 100 /\allowbreak{} interval\_\allowbreak{}cycles}, \texttt{CPU busy = 100 − idle share}. Blocked/sleeping time is not charged to threads. Any residual interval cycles represent omitted/deleted threads. Up to 64 threads are reported. Reports are independent intervals; a new connection needs no previous report to compute percentages.

\section{Backpressure and recovery}\label{wiki-protocols-wire-protocol--backpressure-and-recovery}

The single transport writer sends complete packets to each active link. TCP allocation/send waits are bounded to 250 ms per 1460-byte segment; USB writes retain a one-second timeout. A slow link can delay camera, replies and reports on both links. A partial/failed TCP send closes that session. ToF commands wait for the current camera packet to be consumed; USB command writes use a 0.5-second host timeout. The PC receiver\textquotesingle s eight-second acknowledgement timeout is separate from these transport timeouts.

After copying a completed frame into the transport packet, the firmware immediately rearms the separate DMA capture buffer. Capture overlaps transmission; the transport packet is never overwritten while in use. The next completed snapshot waits for consumption of the previous packet. Slow readers can still cause skipped sensor frames. The monitor estimates sensor FPS from changes in the sensor counter divided by host arrival time (averaged over at least 0.5 s). Buffering/jitter affects this estimate; compare it with received Lidar FPS and reported period. Sensor reset, stop, shutdown and reconnect restart the timing window. Temperature and mode readouts show the most recent frame and its age, including after shutdown.

With no readers, camera packets are discarded and capture is rearmed while ranging remains enabled. Stopping ToF prevents new capture. The IMU producer is independent and can drop queued samples under backpressure. Neither a reconnect nor command timeout resets the sensor or changes a latched fault.

\chapter{Device Monitor}\label{wiki-tools-device-monitor--device-monitor}

{\small\color{MGSMuted}Source: \nolinkurl{wiki/tools/device-monitor.md}\par}\medskip

\section{Installation}\label{wiki-tools-device-monitor--installation}

Ubuntu/Debian dependencies (use Python with Tk support):

\begin{verbatim}
sudo apt install python3-tk python3-pil.imagetk
python3 -m pip install numpy matplotlib pyserial Pillow
\end{verbatim}

Run the viewer without sudo so it uses your user\textquotesingle s Python packages and display session. Ethernet connections do not need root privileges.

\section{Connecting}\label{wiki-tools-device-monitor--connecting}

Launch from VS Code task \textbf{Monitor: Board Viewer}, or:

\begin{verbatim}
python3 test_scripts/Board_Viewer.py --host 192.168.8.20 --connect
python3 test_scripts/Board_Viewer.py --port /dev/ttyACM0 --connect
\end{verbatim}

Select USB/Ethernet, enter the serial device or host and press Connect. The viewer reconnects automatically after a link interruption. To switch endpoints, press Disconnect and wait briefly before Connect. Close other programs using that USB serial port or TCP streaming connection; the board supports one TCP stream client.

\section{Live sensors}\label{wiki-tools-device-monitor--live-sensors}

The top bar shows measured receive throughput (Mbit/s and KiB/s), lidar FPS, IMU Hz, received bytes, missing sequence numbers and parser errors. These are host measurements, not the PHY\textquotesingle s negotiated link speed. Live sensors shows the decoded lidar zone grid and radial depths in mm, accelerometer in mg, gyro in degrees/s, and SFLP quaternion XYZW. \texttt{-\allowbreak{}-\allowbreak{}depth-\allowbreak{}max} controls the heatmap range in mm.

Gray depth pixels represent a received distance of \textbf{0 mm}, which the monitor treats as a missing return. Positive distances, including very small values, use the color scale; values above \texttt{-\allowbreak{}-\allowbreak{}depth-\allowbreak{}max} use its upper endpoint color. The text below the heatmap counts zero returns and shows the nonzero min/max. Point to a pixel to see its row, column and received distance. The distance is the packet\textquotesingle s lower 15 bits; bit 15 selects the ranging step. Earlier versions displayed masked zero returns as white/transparent pixels. An entirely gray heatmap can also mean no frame or an invalid frame; the accompanying text distinguishes those cases. These readouts show what reached the host and do not establish why the sensor produced a zero return.

\section{Viewer performance}\label{wiki-tools-device-monitor--viewer-performance}

The live display defaults to \textbf{20 FPS}. Use \texttt{-\allowbreak{}-\allowbreak{}fps} (1--60) to change the display limit; 10 FPS reduces CPU use on battery, while 30 FPS provides smoother motion:

\begin{verbatim}
python3 test_scripts/Board_Viewer.py --host 192.168.8.20 --connect --fps 10
\end{verbatim}

Reception runs independently of the display rate. The viewer renders the latest depth frame and up to 1000 IMU samples within the latest ten seconds, without queuing old frames for display. Packet counters and sensor frame-rate estimates continue at the incoming rate. Hidden live plots and the hidden CPU donut do not redraw; returning to a tab shows the latest data. Minimizing pauses plot updates. Unchanged telemetry does not trigger a plot redraw.

The renderer caches axes, ticks, grids, legends and the colorbar, then updates the heatmap and traces using Matplotlib blitting. IMU scales expand immediately for new peaks and shrink at most every two seconds, avoiding a full redraw for small fluctuations. Resizing and binning changes rebuild the cache automatically. Status and control widgets refresh separately at 5 Hz. Existing Python packages are sufficient; no additional GUI toolkit or compiler is needed.

Compare the former full-redraw path with the cached renderer on your machine:

\begin{verbatim}
python3 test_scripts/benchmark_viewer.py
python3 test_scripts/benchmark_viewer.py --width 1600 --height 900
\end{verbatim}

The benchmark uses synthetic 54×42 depth frames and 1000 IMU samples, reports median/p95 frame time and CPU time, and needs neither a board nor a display. It measures Agg rendering and excludes Tk\textquotesingle s transfer of pixels to the display; actual GUI performance also depends on window size and display scaling.

\section{Sensor controls}\label{wiki-tools-device-monitor--sensor-controls}

Start ToF / Stop ToF / Reset ToF controls are available above Live sensors and in the Sensor configuration tab. The tab reads the board\textquotesingle s actual settings and lets you edit frame period, exposure, binning, power mode and range context. Stop ranging, wait for acknowledgement, apply settings, then Start ToF. Configuration changes are volatile until reboot; IMU, CPU and metadata reports continue while ToF is stopped. Read from device refreshes the edit fields; periodic status updates do not overwrite unsaved edits. Results and latched sensor faults appear below the controls. Reset ToF resets only the sensor, reloads its firmware and compiled defaults, clears a recoverable fault and leaves ranging stopped; press Start afterward. Runtime edits are discarded and the form reloads the confirmed default settings. IMU, profiling and the connection stay active. A failed reset reports an error and leaves the fault latched. Buttons wait for device acknowledgements, and connection loss clears pending commands without replaying them. A timeout leaves the outcome uncertain; read current state before retrying.

The commands use the same USB serial connection or TCP streaming socket (50002), so no second serial owner or control connection is required. Rebuild/flash the updated application to enable controls. Older firmware can still be viewed but cannot acknowledge these commands. See the \hyperref[wiki-protocols-wire-protocol--magnecko-sense-wire-protocol]{wire protocol README} for exact headers, payload offsets, result codes and byte examples.

\section{Firmware updates}\label{wiki-tools-device-monitor--firmware-updates}

Open \textbf{Firmware update} to update the application over the selected USB or Ethernet connection. Install the \hyperref[wiki-firmware-firmware-updates--first-installation]{DFU-capable FSBL and application baseline} once using SWD before using this tab.

\begin{enumerate}
\def\labelenumi{\arabic{enumi}.}
\tightlist
\item
  Select the device in the connection bar. Ethernet uses the separate update TCP port, \textbf{50003} by default; the normal streaming port remains \textbf{50002}. The update port is editable in the tab or through \texttt{-\allowbreak{}-\allowbreak{}dfu-\allowbreak{}port}.
\item
  Choose \textbf{Read status} to inspect running, confirmed, and pending slots, release identifiers, upload/verification flags, and persistent warnings.
\item
  Browse to \texttt{Appli/\allowbreak{}build/\allowbreak{}magnecko-\allowbreak{}sense-\allowbreak{}fw\_\allowbreak{}Appli-\allowbreak{}signed.\allowbreak{}bin} from a normal flash build and enter its unsigned 32-bit release identifier. Raw application binaries, corrupt images, and images marked \texttt{DEV\_\allowbreak{}MODE} in their manifest are rejected.
\item
  Press \textbf{Upload image}. The progress bar advances after each acknowledged chunk. The device chooses the inactive slot and verifies the entire image.
\item
  With \textbf{Reset after upload and verify boot confirmation} selected, the viewer requests one reset and queries status for up to 60 seconds. It reports either confirmed activation, fallback to the previous image, or an unconfirmed outcome requiring another status read. Uncheck the option to leave the image staged; \textbf{Reset staged image} activates it later.
\end{enumerate}

The window stays responsive during transfers. USB streaming pauses until the updater releases the serial port, then reconnects automatically if it was active before the operation. Ethernet streaming can continue on its own socket. Device selection and sensor commands are locked during an update operation. Disconnect cancels the operation and prevents automatic USB reconnection.

\textbf{Cancel operation} waits for the current acknowledgement, then aborts an unfinished upload. An image already verified remains staged. During boot monitoring, cancellation stops status polling without sending another reset. \textbf{Abort incomplete upload} clears an unfinished transfer shown by Read status; \textbf{Clear warnings} explicitly clears persistent diagnostics. A failed or lost acknowledgement is never blindly replayed: read status before retrying.

If USB re-enumerates under a different device name after reset, select the new serial port and read status. A stable \texttt{/\allowbreak{}dev/\allowbreak{}serial/\allowbreak{}by-\allowbreak{}id/\allowbreak{}} path avoids this on systems that provide one. Upload verification and successful boot are displayed separately. The parameter JSON tab retains the last SWD-generated provenance; current A/B update state and release identifiers are shown in Firmware update.

\section{Device and firmware parameters}\label{wiki-tools-device-monitor--device-and-firmware-parameters}

Device \& firmware parameters shows the complete flash JSON in an expandable tree, including serial/revisions, compiler and CMake versions, build options, compile flags/definitions, Git/image hashes and per-slot flash records. Save received JSON exports it. Every persistent flash refreshes these records automatically; see \hyperref[wiki-firmware-parameters--device-parameters-and-build-metadata]{parameter management} for identity preservation and explicit settings updates. An unprogrammed/unreadable parameter block is reported as \texttt{parameter\_\allowbreak{}error}.

\section{CPU profiler}\label{wiki-tools-device-monitor--cpu-profiler}

CPU profiler shows every ThreadX thread (including middleware threads), its priority, CPU milliseconds and percentage of the latest reporting interval. Idle, interrupts and scheduler overhead have separate rows. The busy meter is 100\% minus idle; sleeping or blocked threads consume no CPU time. Reports arrive about once per second over either transport. The actual interval is displayed because slow transports can delay a report. Report age and a stale indicator make stopped updates visible. Reconnecting clears the previous session\textquotesingle s data. The live donut chart beside the table shows CPU shares with consistent colors and overall CPU usage in the center. Hover over a slice for its time and exact percentage. Small slices (below 1\%, or beyond the seven largest activities) are combined as Other activities; the table keeps every thread\textquotesingle s individual values. Drag the divider to resize the chart and table.

\section{Profiling implementation}\label{wiki-tools-device-monitor--profiling-implementation}

Firmware uses the Cortex-M55 DWT cycle counter and ThreadX execution hooks, with a small idle-entry hook in \texttt{tx\_\allowbreak{}thread\_\allowbreak{}schedule.\allowbreak{}S}. Peripheral ISRs and the ThreadX tick bracket interrupt processing, including nested interrupts. The 1 kHz tick extends the 32-bit cycle counter into 64-bit interval totals; counters are snapped atomically, then serialized by the existing transport thread. This measures the application\textquotesingle s current running-core/busy-idle configuration.

\section{Parameter packets and timestamps}\label{wiki-tools-device-monitor--parameter-packets-and-timestamps}

Firmware reads the first 16 KiB of the reserved parameter block once at startup, then sends the stored JSON contents as a \texttt{PARM} packet every two seconds on USB and Ethernet. The header is \texttt{<4sIIHH}: magic PARM, sequence 0, JSON byte length, format version 1, reserved 0. The NUL terminator is not transmitted. Repetition allows late connections to receive identity without restarting the board. JSON parsing happens on the host; firmware forwards the stored settings and does not apply them to networking or sensors. The packer enforces the 16 KiB limit.

Under \textbf{Device \& firmware parameters → parameter\_build → built\_at\_utc}, the timestamp records the most recent parameter build. It refreshes independently of the per-firmware dates under \textbf{build}. A parameter-only build leaves those firmware dates unchanged. Flash the generated parameter block and restart the board to see its new timestamp; firmware caches the parameter JSON at startup. Flashing Appli/FSBL also refreshes the selected firmware record\textquotesingle s \texttt{built\_\allowbreak{}at\_\allowbreak{}utc} and \texttt{programmed\_\allowbreak{}at\_\allowbreak{}utc}. Its \texttt{image\_\allowbreak{}built\_\allowbreak{}at\_\allowbreak{}utc} preserves the original image build time; unflashed slots retain their previous records.

Rebuild/flash before use (firmware flashes also update the parameters):

\begin{verbatim}
./scripts/flash.sh --target all --config Release --build
\end{verbatim}

No GUI dependencies are needed by the existing USB\_Debug/Ethernet\_Debug parsers beyond their own plotting dependencies; those viewers skip PARM packets and print CPU measurements as \texttt{[CPU]} status lines. Reconstruction ignores profiler and parameter packets while preserving them in its raw stream recording.

\section{Sensor diagnostics}\label{wiki-tools-device-monitor--sensor-diagnostics}

The Sensor configuration tab also offers \textbf{Power down} (XSHUT shutdown) and a \textbf{Ranging mode} selector: Precision (no DSS), Precision, Ambient. Stop before applying settings. To wake a powered-down sensor, use Reset ToF, then Start ToF; reset restores compiled defaults. These controls need the matching firmware.

The live and configuration tabs show die temperature, estimated sensor FPS, frame-reported mode/period, frame counter and error. Compare the sensor estimate with received Lidar FPS in the status bar. Last-frame age makes stale readings visible after Stop/Power down. See \hyperref[wiki-tools-streaming-and-diagnostics--gaming-mode-frame-rate-and-temperature]{timing diagnostics}.

\chapter{Streaming and sensor diagnostics}\label{wiki-tools-streaming-and-diagnostics--streaming-and-sensor-diagnostics}

{\small\color{MGSMuted}Source: \nolinkurl{wiki/tools/streaming-and-diagnostics.md}\par}\medskip

\section{Ethernet streaming}\label{wiki-tools-streaming-and-diagnostics--ethernet-streaming}

The application sends identical camera and IMU packets over USB CDC and Ethernet TCP. To view depth, acceleration, gyro and SFLP over Ethernet, run:

\begin{verbatim}
python3 test_scripts/Ethernet_Debug.py --host 192.168.8.20
\end{verbatim}

TCP port defaults to 50002 (\texttt{-\allowbreak{}-\allowbreak{}tcp-\allowbreak{}port} changes the client port). The server forwards telemetry when one client connects; no handshake is required. ToF starts at boot, and can be stopped/restarted from the Device Monitor on USB or TCP 50002. The viewer reconnects automatically and resets partial packets on reconnect. Ethernet can be plugged in after boot or unplugged/replugged while running. The network thread polls the PHY every 200 ms, retries MAC/DMA enable when link appears, and publishes IP readiness only after enable succeeds. Link removal withdraws readiness and stops MAC/DMA. Negotiation time is additional to polling. TCP listeners discard old sessions and incomplete handshakes on a cable cycle; the Device Monitor reconnects automatically. ToF/IMU operation and USB continue. The IPv4 address remains static: configure the PC\textquotesingle s Ethernet adapter on \texttt{192.\allowbreak{}168.\allowbreak{}8.\allowbreak{}0/\allowbreak{}24} with a different host address (for example \texttt{192.\allowbreak{}168.\allowbreak{}8.\allowbreak{}10/\allowbreak{}24}). The legacy echo server remains on port 50000; the throughput server and its benchmark scripts have been removed. The normal full application remains Release.

\section{USB streaming}\label{wiki-tools-streaming-and-diagnostics--usb-streaming}

For USB, run:

\begin{verbatim}
python3 test_scripts/USB_Debug.py --port /dev/ttyACM1
\end{verbatim}

Omit \texttt{-\allowbreak{}-\allowbreak{}port} to discover devices and follow changed ttyACM numbers on reconnect. The plot window shows a ToF depth heatmap plus acceleration xyz (mg), gyro xyz (deg/s), and SFLP quaternion xyzw over the most recent 10 seconds. IMU plots use the MCU timestamps and clear when the device reconnects. Plots target 30 FPS (\texttt{-\allowbreak{}-\allowbreak{}fps 60} requests 60). The depth grid is decoded from the latest complete frame using status-row binning and crop metadata, rather than CSI transport dimensions (e.g. 100×149 carries a 54×42 zone grid). Little-endian depth words contain radial distance in millimetres in bits 0--14; bit 15 selects the ranging step and is not an invalid-data flag. Zero returns appear gray. The colour scale defaults to 0--4000 mm (\texttt{-\allowbreak{}-\allowbreak{}depth-\allowbreak{}max 9000} changes its maximum); hover to see a zone\textquotesingle s depth. This displays raw radial distances without ST\textquotesingle s calibration, confidence filtering, or radial-to-perpendicular postprocessing. See \href{https://www.st.com/resource/en/technical_note/tn1641-vl53l9cx-postprocessing-pipeline-stmicroelectronics.pdf\#page=5}{TN1641, section 2.3}. The IMU time axis is relative to the latest sample, and acceleration/gyro ranges expand when needed.

\section{Packet framing and profiling}\label{wiki-tools-streaming-and-diagnostics--packet-framing-and-profiling}

Both packet headers use \texttt{<4sIIHH} (little-endian): magic, sequence number, payload length, and two type-specific fields. Sequence numbers are independent per sensor.

{\def\LTcaptype{none} % do not increment counter
\begin{longtable}[]{@{}
  >{\raggedright\arraybackslash}p{(\linewidth - 4\tabcolsep) * \real{0.2135}}
  >{\raggedright\arraybackslash}p{(\linewidth - 4\tabcolsep) * \real{0.3450}}
  >{\raggedright\arraybackslash}p{(\linewidth - 4\tabcolsep) * \real{0.4415}}@{}}
\toprule\noalign{}
\begin{minipage}[b]{\linewidth}\raggedright
\textbf{Magic}
\end{minipage} & \begin{minipage}[b]{\linewidth}\raggedright
\textbf{Header fields 12/14}
\end{minipage} & \begin{minipage}[b]{\linewidth}\raggedright
\textbf{Payload}
\end{minipage} \\
\midrule\noalign{}
\endhead
\bottomrule\noalign{}
\endlastfoot
\texttt{VL9F} & width / height & width × height RAW8 bytes, including sensor status row \\
\texttt{IMUF} & 10 float channels / format version 1 & \texttt{<I10f}: MCU timestamp in milliseconds, acc xyz in mg, gyro xyz in mdps, SFLP quaternion xyzw \\
\texttt{PARM} & format version 1 / reserved 0 & UTF-8 device/firmware parameter JSON, repeated every two seconds \\
\texttt{PROF} & format version 1 / thread count & CPU profiling interval, sent about once per second \\
\texttt{TOFS} & format version 2 / acknowledged operation (0 for periodic) & 36-byte sensor configuration/state, DSS flags and command result; PC also accepts v1 (32 bytes) \\
\texttt{CTRL} (host → board) & format version 1 or 2 / operation & Empty for Get/Start/Stop/Reset/Power down; 12-byte settings for Set (v2 adds DSS flags) \\
\end{longtable}
}

\texttt{PROF} has its own sequence counter. Its payload starts with \texttt{<IIQQQQHH}: MCU timestamp (ms), core clock (Hz), interval cycles, idle cycles, interrupt cycles, scheduler cycles, thread count, flags. Each thread record is \texttt{<IIQ32s}: thread ID (control-block address), priority, interval CPU cycles and a NUL-terminated UTF-8 name (up to 31 bytes). All numbers are little-endian. Up to 64 threads are reported; flags bit 0 indicates a truncated thread list. Unreported/removed-thread cycles appear as \texttt{Other /\allowbreak{} removed threads} on the PC. CPU time in ms is \texttt{cycles * 1000 /\allowbreak{} core\_\allowbreak{}clock}; CPU percentage is \texttt{cycles * 100 /\allowbreak{} interval\_\allowbreak{}cycles}. Report windows continue across connections, so receiving one complete report is enough to display usage. See the \hyperref[wiki-tools-device-monitor--device-monitor]{desktop viewer} for the CPU profiler tab.

\section{Capture and IMU timing}\label{wiki-tools-streaming-and-diagnostics--capture-and-imu-timing}

Camera capture uses DCMIPP PIPE0\textquotesingle s dump path to preserve RAW8 bytes without ISP processing or padded row pitch. Watch \texttt{hdcmipp.\allowbreak{}PipeState[0]} for this capture. PIPE0 setup lives in \texttt{DCMIPP\_\allowbreak{}Init 2}\textquotesingle s CubeMX USER CODE block in \texttt{dcmipp.\allowbreak{}c}, so it survives regeneration even while CubeMX also configures PIPE1.

IMU packets are 60 bytes total. Acc, gyro, and SFLP run at 120 Hz; the thread publishes the latest readings at 100 Hz after all three have produced data. The timestamp records host acquisition time, not an IMU hardware timestamp. SFLP is the game rotation quaternion, so it does not provide a compass heading. Transports use a bounded 16-packet queue: sampling continues during backpressure and full queues drop samples. Sequence gaps expose lost samples at the host.

When stationary, acceleration magnitude should be about 1000 mg, gyro near zero, and quaternion magnitude near one. The host should report about 100 IMU packets/s when the transport keeps up. IMU streaming continues after a VL53L9 sensor fault. IMU sampling has priority over the USB application thread. FIFO overruns flush the FIFO and resume after fresh acc/gyro/SFLP readings; they no longer terminate streaming. A sequence gap flags the discontinuity.

\section{Backpressure and link recovery}\label{wiki-tools-streaming-and-diagnostics--backpressure-and-link-recovery}

A single transport thread owns packet consumption and sends complete packets to USB (if CDC/DTR is active) and the TCP client (if connected). With neither active, frames are discarded and capture is rearmed. This avoids two threads competing for the same camera buffer or splitting IMU packets between clients. Both connections can be used together; slow readers can reduce camera capture rate. TCP allocation/send waits are bounded to 250 ms per segment; a failed/partial write closes that session so a future connection starts at a fresh packet.

To check cable recovery on hardware, boot without Ethernet, wait more than five seconds, then insert the cable. Once link negotiates, verify \texttt{ping 192.\allowbreak{}168.\allowbreak{}8.\allowbreak{}20} and connect the monitor to TCP 50002. Unplug during streaming, reinsert, and confirm that the monitor reconnects without resetting the board. Also exercise unplugging while the PC is attempting a connection. Native regressions cover late insertion, failed enable/disable retries, an already enabled link at boot, half-open TCP cleanup and a full cable cycle during a blocked transport write. USB writes retain their existing one-second timeout. Sampling is not blocked by the queue, so backpressure can drop IMU samples. One TCP client is supported.

\section{Isolating sensor faults}\label{wiki-tools-streaming-and-diagnostics--isolating-sensor-faults}

The following comparisons can isolate an intermittent VL53L9 fault using one board. Reset/reboot the sensor between attempts and record its firmware code and all five laser-driver bytes. The VL9F sequence counts frames captured for USB; it is not the sensor\textquotesingle s own frame counter or a fault code. Repeated runs may fail at different sequence numbers without changing the underlying fault.

\begin{enumerate}
\def\labelenumi{\arabic{enumi}.}
\item
  Run without the Python receiver (or USB host). Camera packets are discarded while disconnected, so this checks whether host traffic changes the fault.
\item
  Keep exposure, binning, power mode and CSI bitrate unchanged, but change the frame period in the selected profile in \texttt{Appli/\allowbreak{}Libs/\allowbreak{}53L9/\allowbreak{}interface/\allowbreak{}vl53l9/\allowbreak{}vl53l9\_\allowbreak{}utils.\allowbreak{}c}, then rebuild.

  Flash the resulting ELF and compare repeated runs by elapsed time as well as frame count. Improvement suggests a duty-cycle/timing dependency; it does not establish that the laser-driver path is healthy.
\item
  Run the selected profile with DCMIPP capture disabled:

\begin{verbatim}
cmake -S Appli -B Appli/build -DVL53L9_SENSOR_ONLY=ON
cmake --build Appli/build -j 4
\end{verbatim}

  The sensor still emits CSI frames, but no VL9F packets are produced. IMU packets continue. The stream thread polls sensor status once a second. A recurring \texttt{0x0F00} / \texttt{laser\_\allowbreak{}driver[3] = 0x80} in this mode establishes that STM32 frame capture and camera USB transfers are not required to trigger it. This test leaves the sensor\textquotesingle s CSI transmitter enabled and does not isolate every possible electrical interaction with the board.
\end{enumerate}

Restore normal settings afterward:

\begin{verbatim}
cmake -S Appli -B Appli/build -DVL53L9_SENSOR_ONLY=OFF
cmake --build Appli/build -j 4
\end{verbatim}

The diagnostic options do not change laser safety checks or automatically retry faulted emission. The firmware preserves sensor fault status, stops capture and latches a fault. The command service stays alive to report the fault; Start is rejected until an explicit Reset ToF (or board reset). Reset ToF restores compiled defaults and leaves ranging stopped; Start must be requested separately. IMU and profiler reports continue during sensor reset. ST\textquotesingle s error handling and bit definitions are in \href{https://www.st.com/resource/en/user_manual/um3683-programing-guide-stmicroelectronics.pdf\#page=16}{UM3683, sections 2.4--2.4.2}.

\section{Gaming mode, frame rate and temperature}\label{wiki-tools-streaming-and-diagnostics--gaming-mode-frame-rate-and-temperature}

The monitor\textquotesingle s \textbf{Sensor configuration} tab exposes Precision (no DSS), Precision and Ambient. These map to short/DSS off, short/DSS short and long/DSS long. Previously Gaming selected long context and \texttt{vl53l9\_\allowbreak{}set\_\allowbreak{}binning} enabled DSS; it therefore used Ambient rather than the datasheet\textquotesingle s Gaming mode. Gaming now selects Precision (no DSS), 54×42, regular power and a 10 ms requested period, keeping the local 1 ms exposure. Other presets retain DSS.

ST\textquotesingle s \href{https://www.st.com/resource/en/datasheet/vl53l9cx.pdf}{VL53L9CX datasheet, Table 9} uses \textbf{4 ms} exposure for its 100 fps Gaming example. A 1 ms setting is a deliberate shorter exposure. Low/ultralow modes turn blocks off between frames; regular mode keeps them powered. Exposure alone therefore does not determine power or temperature. The \href{https://www.st.com/resource/en/user_manual/um3683-programing-guide-stmicroelectronics.pdf}{programming guide, Tables 22--23} describes context/DSS mapping and processing time: full-resolution Precision DSS computation is listed as 13.5 ms, running in parallel with exposure. A readable 10 ms period register alone cannot establish actual frame rate.

CSI snapshots previously waited for the transport to finish before rearming. The next capture is now armed immediately after copying into the separate transport packet. This removes a source of alternate-frame skipping; transport backpressure can still limit received FPS. The monitor displays both received Lidar FPS and \textbf{estimated sensor FPS} from the sensor frame counter. A sensor estimate near 100 with received FPS near 50 points to capture/transport loss; if both are near 50, investigate sensor timing/mode and measurement jitter. The GUI redraw rate is separate from packet reception.

The monitor also displays die temperature in °C, the mode and period reported in the frame, the frame counter, firmware error and the age of the last frame. Temperature is retained as a last-frame reading after stopping/shutdown; it is not a live reading while off. Confirm temperature/error details on hardware before concluding that heating caused a fault. Protection remains enabled.

\textbf{Power down} holds XSHUT low and keeps IMU/network/profiler operation running. Supply rails remain connected. To wake: \textbf{Reset ToF → Start ToF}. Reset restores compiled defaults and leaves emission stopped. Stop pauses ranging without asserting XSHUT. The full wire format is in \hyperref[wiki-protocols-wire-protocol--magnecko-sense-wire-protocol]{protocol/README.md}.

\chapter{Scene reconstruction}\label{wiki-tools-scene-reconstruction--scene-reconstruction}

{\small\color{MGSMuted}Source: \nolinkurl{wiki/tools/scene-reconstruction.md}\par}\medskip

\section{Recording a scan}\label{wiki-tools-scene-reconstruction--recording-a-scan}

With uncertain IMU/lidar axes, record slow rotations about at least two axes while viewing a static scene with depth structure (corners/furniture, not just a flat wall). The following records five seconds after both streams arrive, tries to estimate the mounting rotation, and reconstructs the scene:

\begin{verbatim}
python3 test_scripts/Reconstruct_Scene.py --port /dev/ttyACM1 --duration 5 \
  --estimate-mounting --output scans/scan01 --preview
\end{verbatim}

Mounting estimation compares unseeded lidar ICP rotations with SFLP relative rotations and solves the rotation equation A X = X B. It rejects insufficient motion or inconsistent fits. It cannot guarantee correctness on ambiguous scenes, and it cannot infer mounting translation from unconstrained handheld motion. Translation defaults to zero unless supplied in a calibration JSON. Inspect the estimated rotation and validate it with known physical motion before trusting it.

\section{Output artifacts}\label{wiki-tools-scene-reconstruction--output-artifacts}

Output files:

\begin{itemize}
\tightlist
\item
  \texttt{usb\_\allowbreak{}stream.\allowbreak{}bin}: original received USB bytes.
\item
  \texttt{capture.\allowbreak{}npz}: depth payloads/sequence numbers, all IMU readings (MCU timestamp, acc mg, gyro mdps, quaternion xyzw), and host arrival times. Load without pickle.
\item
  \texttt{depth\_\allowbreak{}video.\allowbreak{}gif}: depth visualization; maximum colour scale is 4000 mm.
\item
  \texttt{estimated\_\allowbreak{}calibration.\allowbreak{}json}: rotation fit, when estimation succeeds.
\item
  \texttt{scene.\allowbreak{}ply}: reconstructed XYZ point cloud in metres, viewable in point-cloud tools.
\item
  \texttt{poses.\allowbreak{}npz}: accepted \texttt{world\_\allowbreak{}from\_\allowbreak{}lidar} 4×4 transforms and their timestamps.
\item
  \texttt{report.\allowbreak{}json}: accepted/rejected frames, ICP overlap/RMSE, calibration assumptions.
\end{itemize}

\section{Replaying a capture}\label{wiki-tools-scene-reconstruction--replaying-a-capture}

The output directory must be empty. If rotation estimation fails, the recording and depth GIF remain available. Reprocess without reconnecting or recapturing:

\begin{verbatim}
python3 test_scripts/Reconstruct_Scene.py --input scans/scan01/capture.npz \
  --calibration mounting.json --output scans/reprocessed --preview
\end{verbatim}

\section{Calibration files}\label{wiki-tools-scene-reconstruction--calibration-files}

A calibration file uses this schema (identity below is an example, not a measured calibration for this board):

\begin{verbatim}
{
  "rotation_lidar_to_imu": [[1, 0, 0], [0, 1, 0], [0, 0, 1]],
  "translation_lidar_to_imu_m": [0, 0, 0],
  "horizontal_fov_deg": 55,
  "vertical_fov_deg": 42,
  "invert_sflp_quaternion": false,
  "mounting_calibrated": true
}
\end{verbatim}

\section{Coordinate conventions}\label{wiki-tools-scene-reconstruction--coordinate-conventions}

The transform is \texttt{point\_\allowbreak{}imu = rotation\_\allowbreak{}lidar\_\allowbreak{}to\_\allowbreak{}imu @ point\_\allowbreak{}lidar + translation}. Translation is the lidar origin\textquotesingle s position relative to the IMU origin, expressed in \textbf{IMU axes}, in metres. Both coordinate systems must be right-handed. The nominal lidar projection uses X toward increasing columns, Y toward increasing rows (down), and Z forward. The expected SFLP convention is IMU-to-world xyzw; \texttt{invert\_\allowbreak{}sflp\_\allowbreak{}quaternion} transposes that rotation if your convention is reversed.

The provided PCB locations, IMU \texttt{(10, 10)} mm and lidar \texttt{(33.\allowbreak{}4, -\allowbreak{}10)} mm, imply a PCB displacement \texttt{(23.\allowbreak{}4, -\allowbreak{}20)} mm. This is not yet an IMU-coordinate translation: its axes must first be confirmed. Do not infer an axis inversion from placement alone. A 180° Z rotation maps X and Y to their negatives while preserving Z; negating all three axes is a reflection, not a valid mounting rotation.

\section{Board mounting calibration}\label{wiki-tools-scene-reconstruction--board-mounting-calibration}

\texttt{test\_\allowbreak{}scripts/\allowbreak{}mounting.\allowbreak{}json} now includes the rotation inferred from the reported static-gravity observations: upright +IMU X, clockwise 90° (viewed from behind) +IMU Y, pointing the lidar upwards −IMU Z. With lidar X right, Y down, Z forward, the mapping is lidar X → −IMU Y, lidar Y → −IMU X, lidar Z → −IMU Z. Its matrix is right-handed. This assumes static measurements with the displayed image aligned to those lidar image axes. Translation is explicitly provisional zero until the PCB-to-IMU axes and Z offset are known; nominal FoV rays remain approximate.

Use that rotation without estimating a new one:

\begin{verbatim}
python3 test_scripts/Reconstruct_Scene.py --port /dev/ttyACM1 --duration 5 \
  --calibration test_scripts/mounting.json --output scans/scan02 --preview
\end{verbatim}

\section{Calibrated projection rays}\label{wiki-tools-scene-reconstruction--calibrated-projection-rays}

If calibrated per-zone rays are available, add \texttt{"ray\_\allowbreak{}directions\_\allowbreak{}npy":\allowbreak{} "rays.\allowbreak{}npy"}. That file must contain an H×W×3 array for the \textbf{uncropped raw grid}, in the same zone order as the received depth, with directions expressed in lidar axes. Rays are normalized by the script, and crop metadata selects the corresponding subset. A size mismatch is rejected. Otherwise the script uses an approximate pinhole ray model from the nominal 55°×42° FoV, documented in \href{https://www.st.com/resource/en/datasheet/vl53l9cx.pdf}{ST\textquotesingle s datasheet}.

\section{Approximate mode}\label{wiki-tools-scene-reconstruction--approximate-mode}

For a deliberately approximate first experiment, \texttt{-\allowbreak{}-\allowbreak{}approximate} assumes aligned IMU/lidar axes and zero offset. It should not be used to interpret unknown mounting as calibrated geometry. Configuration can be refined and the same capture replayed.

\section{Reconstruction pipeline}\label{wiki-tools-scene-reconstruction--reconstruction-pipeline}

Reconstruction converts the raw \textbf{radial} range to metres along normalized rays. It applies crop/binning metadata, rejects faulted frames and ranges outside 0.05--8.8 m, and uses SFLP (the IMU\textquotesingle s fused acc/gyro orientation) as a rotation prior. Trimmed point-to-point ICP estimates translation and refines rotation against the previous accepted frame. Poor overlap, residuals or large corrections are rejected. Accepted frames are merged into 2 cm voxels (\texttt{-\allowbreak{}-\allowbreak{}voxel} changes this); reconstruction selects up to 10 depth frames/s (\texttt{-\allowbreak{}-\allowbreak{}fps} changes this) while recording retains all packets. The IMU lever-arm offset contributes to predicted lidar displacement.

\section{Accuracy limits}\label{wiki-tools-scene-reconstruction--accuracy-limits}

Accuracy limits:

\begin{itemize}
\tightlist
\item
  There is no shared depth/IMU timestamp in VL9F. Association uses USB arrival times, which include batching and transport latency. Move slowly.
\item
  SFLP supplies relative attitude without a compass; yaw and ICP translation can drift. Accelerometer readings are saved, not double-integrated for position.
\item
  Raw depth has not received ST\textquotesingle s calibration/confidence/postprocessing pipeline.
\item
  ICP cannot reliably determine motion in every scene, particularly flat, repetitive or moving scenes. There is no loop closure, pose-graph optimization, or guaranteed SLAM accuracy. A PLY point cloud is produced, not a watertight mesh.
\end{itemize}

\chapter{Board bring-up and system test plan}\label{wiki-hardware-bring-up-and-testing--magnecko-sense-mgs--bring-up--system-test-protocol}

{\small\color{MGSMuted}Source: \nolinkurl{wiki/hardware/bring-up-and-testing.md}\par}\medskip

\begin{quote}
This is a test plan, not a record of completed validation. It includes historical and planned UDP/FoE checks. Current behavior is described in the \hyperref[wiki-protocols-wire-protocol--magnecko-sense-wire-protocol]{wire protocol} and \hyperref[wiki-firmware-firmware-updates--application-firmware-updates-and-ab-recovery]{firmware update guide}.
\end{quote}

\section{Test Record}\label{wiki-hardware-bring-up-and-testing--0-test-record}

\begin{itemize}
\tightlist
\item
  \textbf{Board serial / ID:} \_\_\_\_\_\_\_\_\_\_\_\_\_\_\_\_\_\_\_\_\_\_
\item
  \textbf{PCB revision:} \_\_\_\_\_\_\_\_\_\_\_\_\_\_\_\_\_\_\_\_\_\_\_\_\_\_\_
\item
  \textbf{Assembly revision:} \_\_\_\_\_\_\_\_\_\_\_\_\_\_\_\_\_\_\_\_\_\_
\item
  \textbf{Firmware commit / Git hash:} \_\_\_\_\_\_\_\_\_\_\_\_\_\_\_\_\_\_\_\_\_\_
\item
  \textbf{Date:} \_\_\_\_\_\_\_\_\_\_\_\_\_\_\_\_\_\_\_\_\_\_
\item
  \textbf{Tester:} \_\_\_\_\_\_\_\_\_\_\_\_\_\_\_\_\_\_\_\_\_\_
\end{itemize}

\subsubsection{Test equipment}\label{wiki-hardware-bring-up-and-testing--test-equipment}

\begin{itemize}
\tightlist
\item[$\square$]
  Bench supply with adjustable current limit
\item[$\square$]
  DMM
\item[$\square$]
  Oscilloscope
\item[$\square$]
  Logic analyzer
\item[$\square$]
  ST-Link / debugger
\item[$\square$]
  Ethernet-capable PC
\item[$\square$]
  Known-good Ethernet cable
\item[$\square$]
  USB-C cable / PD source
\item[$\square$]
  Optional thermal camera / IR thermometer
\item[$\square$]
  Optional differential probe / high-bandwidth probes
\end{itemize}

\subsubsection{Global rule}\label{wiki-hardware-bring-up-and-testing--global-rule}

If a test fails:

\begin{itemize}
\tightlist
\item[$\square$]
  Stop before progressing to dependent tests
\item[$\square$]
  Record the failure
\item[$\square$]
  Record measured voltages/currents/signals
\item[$\square$]
  Save relevant scope captures / logs
\item[$\square$]
  Fix or understand the failure before continuing
\end{itemize}

\section{Unpowered Inspection}\label{wiki-hardware-bring-up-and-testing--1-unpowered-inspection}

\subsection{Visual inspection}\label{wiki-hardware-bring-up-and-testing--11-visual-inspection}

\begin{itemize}
\tightlist
\item[$\square$]
  Correct PCB revision received
\item[$\square$]
  Correct component population
\item[$\square$]
  No visibly rotated ICs
\item[$\square$]
  No solder bridges
\item[$\square$]
  No tombstoned passives
\item[$\square$]
  No damaged connectors
\item[$\square$]
  SEMPER flash orientation correct
\item[$\square$]
  STM32N657 orientation correct
\item[$\square$]
  ADIN1300 orientation correct
\item[$\square$]
  LSM6DSV16X orientation correct
\item[$\square$]
  VL53L9 orientation correct
\item[$\square$]
  USB-C connector mechanically sound
\item[$\square$]
  Ethernet connector / magnetics mechanically sound
\item[$\square$]
  Debug connector accessible
\item[$\square$]
  BOOT / reset controls accessible
\end{itemize}

\textbf{Notes:}

\subsection{Unpowered short-circuit checks}\label{wiki-hardware-bring-up-and-testing--12-unpowered-short-circuit-checks}

Measure resistance from each major rail to GND.

{\def\LTcaptype{none} % do not increment counter
\begin{longtable}[]{@{}
  >{\raggedright\arraybackslash}p{(\linewidth - 6\tabcolsep) * \real{0.3086}}
  >{\raggedleft\arraybackslash}p{(\linewidth - 6\tabcolsep) * \real{0.2736}}
  >{\raggedleft\arraybackslash}p{(\linewidth - 6\tabcolsep) * \real{0.2239}}
  >{\centering\arraybackslash}p{(\linewidth - 6\tabcolsep) * \real{0.1940}}@{}}
\toprule\noalign{}
\begin{minipage}[b]{\linewidth}\raggedright
\textbf{Rail}
\end{minipage} & \begin{minipage}[b]{\linewidth}\raggedleft
\textbf{Expected}
\end{minipage} & \begin{minipage}[b]{\linewidth}\raggedleft
\textbf{Measured}
\end{minipage} & \begin{minipage}[b]{\linewidth}\centering
\textbf{Pass}
\end{minipage} \\
\midrule\noalign{}
\endhead
\bottomrule\noalign{}
\endlastfoot
Main input / VBUS & No hard short & \_\_\_\_\_\_ Ω & {[} {]} \\
5 V & No hard short & \_\_\_\_\_\_ Ω & {[} {]} \\
3.3 V & No hard short & \_\_\_\_\_\_ Ω & {[} {]} \\
1.8 V & No hard short & \_\_\_\_\_\_ Ω & {[} {]} \\
MCU core rail(s) & No hard short & \_\_\_\_\_\_ Ω & {[} {]} \\
Analog rail(s) & No hard short & \_\_\_\_\_\_ Ω & {[} {]} \\
Other: \_\_\_\_\_\_\_\_\_\_ & No hard short & \_\_\_\_\_\_ Ω & {[} {]} \\
\end{longtable}
}

\begin{quote}
Low resistance is not automatically a fault on core rails. Compare against the schematic/BOM and expected load.
\end{quote}

\begin{itemize}
\tightlist
\item[$\square$]
  No unexpected short between adjacent power rails
\item[$\square$]
  No short from Ethernet shield/chassis domain to digital ground unless intentionally designed
\item[$\square$]
  USB VBUS not unintentionally shorted to internal rails
\end{itemize}

\section{First Power-Up}\label{wiki-hardware-bring-up-and-testing--2-first-power-up}

\subsection{Bench-supply preparation}\label{wiki-hardware-bring-up-and-testing--21-bench-supply-preparation}

Before connecting the board:

\begin{itemize}
\tightlist
\item[$\square$]
  Set supply voltage to expected board input
\item[$\square$]
  Set conservative current limit
\item[$\square$]
  Verify polarity
\item[$\square$]
  Disable output
\item[$\square$]
  Connect DMM / scope to main rail(s)
\end{itemize}

Initial current limit: \_\_\_\_\_\_\_\_\_\_ A\\
Input voltage: \_\_\_\_\_\_\_\_\_\_ V

\subsection{First power-on}\label{wiki-hardware-bring-up-and-testing--22-first-power-on}

\begin{itemize}
\tightlist
\item[$\square$]
  Power enabled
\item[$\square$]
  No immediate current-limit hit
\item[$\square$]
  No audible/visible abnormality
\item[$\square$]
  No component heating rapidly
\item[$\square$]
  Input current appears reasonable
\end{itemize}

Measured input current: \_\_\_\_\_\_\_\_\_\_ mA

\subsubsection{Stop immediately if}\label{wiki-hardware-bring-up-and-testing--stop-immediately-if}

\begin{itemize}
\tightlist
\item[$\square$]
  Current limit is reached unexpectedly
\item[$\square$]
  Any regulator output is grossly incorrect
\item[$\square$]
  Any component heats rapidly
\item[$\square$]
  Supply oscillates / hiccups unexpectedly
\end{itemize}

\section{Power Rails}\label{wiki-hardware-bring-up-and-testing--3-power-rails}

Measure every rail before connecting the debugger.

{\def\LTcaptype{none} % do not increment counter
\begin{longtable}[]{@{}
  >{\raggedright\arraybackslash}p{(\linewidth - 8\tabcolsep) * \real{0.2302}}
  >{\raggedleft\arraybackslash}p{(\linewidth - 8\tabcolsep) * \real{0.1923}}
  >{\raggedleft\arraybackslash}p{(\linewidth - 8\tabcolsep) * \real{0.1923}}
  >{\raggedleft\arraybackslash}p{(\linewidth - 8\tabcolsep) * \real{0.2158}}
  >{\centering\arraybackslash}p{(\linewidth - 8\tabcolsep) * \real{0.1694}}@{}}
\toprule\noalign{}
\begin{minipage}[b]{\linewidth}\raggedright
\textbf{Rail}
\end{minipage} & \begin{minipage}[b]{\linewidth}\raggedleft
\textbf{Nominal}
\end{minipage} & \begin{minipage}[b]{\linewidth}\raggedleft
\textbf{Measured}
\end{minipage} & \begin{minipage}[b]{\linewidth}\raggedleft
\textbf{Ripple}
\end{minipage} & \begin{minipage}[b]{\linewidth}\centering
\textbf{Pass}
\end{minipage} \\
\midrule\noalign{}
\endhead
\bottomrule\noalign{}
\endlastfoot
Main input & \_\_\_\_\_\_ V & \_\_\_\_\_\_ V & \_\_\_\_\_\_ mVpp & {[} {]} \\
5 V & \_\_\_\_\_\_ V & \_\_\_\_\_\_ V & \_\_\_\_\_\_ mVpp & {[} {]} \\
3.3 V & \_\_\_\_\_\_ V & \_\_\_\_\_\_ V & \_\_\_\_\_\_ mVpp & {[} {]} \\
1.8 V & \_\_\_\_\_\_ V & \_\_\_\_\_\_ V & \_\_\_\_\_\_ mVpp & {[} {]} \\
MCU core & \_\_\_\_\_\_ V & \_\_\_\_\_\_ V & \_\_\_\_\_\_ mVpp & {[} {]} \\
PHY supply & \_\_\_\_\_\_ V & \_\_\_\_\_\_ V & \_\_\_\_\_\_ mVpp & {[} {]} \\
Flash supply & \_\_\_\_\_\_ V & \_\_\_\_\_\_ V & \_\_\_\_\_\_ mVpp & {[} {]} \\
Sensor supply & \_\_\_\_\_\_ V & \_\_\_\_\_\_ V & \_\_\_\_\_\_ mVpp & {[} {]} \\
\end{longtable}
}

\begin{itemize}
\tightlist
\item[$\square$]
  Rail sequencing is plausible
\item[$\square$]
  No regulator remains in hiccup mode
\item[$\square$]
  Power-good signals behave as expected
\item[$\square$]
  Reset supervisors / enables behave correctly
\item[$\square$]
  Rail ripple is acceptable
\end{itemize}

\textbf{Scope captures saved:} {[} {]} Yes {[} {]} No

\section{Clocks, Reset, Boot Pins}\label{wiki-hardware-bring-up-and-testing--4-clocks-reset-boot-pins}

\subsection{Clocks}\label{wiki-hardware-bring-up-and-testing--41-clocks}

\begin{itemize}
\tightlist
\item[$\square$]
  STM32 clock source present
\item[$\square$]
  ADIN1300 reference clock present
\item[$\square$]
  XSPI clock only active when expected
\item[$\square$]
  USB reference clock configuration valid
\item[$\square$]
  Other required oscillators present
\end{itemize}

{\def\LTcaptype{none} % do not increment counter
\begin{longtable}[]{@{}
  >{\raggedright\arraybackslash}p{(\linewidth - 6\tabcolsep) * \real{0.2365}}
  >{\raggedleft\arraybackslash}p{(\linewidth - 6\tabcolsep) * \real{0.2747}}
  >{\raggedleft\arraybackslash}p{(\linewidth - 6\tabcolsep) * \real{0.2747}}
  >{\centering\arraybackslash}p{(\linewidth - 6\tabcolsep) * \real{0.2141}}@{}}
\toprule\noalign{}
\begin{minipage}[b]{\linewidth}\raggedright
\textbf{Clock}
\end{minipage} & \begin{minipage}[b]{\linewidth}\raggedleft
\textbf{Expected}
\end{minipage} & \begin{minipage}[b]{\linewidth}\raggedleft
\textbf{Measured}
\end{minipage} & \begin{minipage}[b]{\linewidth}\centering
\textbf{Pass}
\end{minipage} \\
\midrule\noalign{}
\endhead
\bottomrule\noalign{}
\endlastfoot
MCU ref & \_\_\_\_\_\_ MHz & \_\_\_\_\_\_ MHz & {[} {]} \\
PHY ref & \_\_\_\_\_\_ MHz & \_\_\_\_\_\_ MHz & {[} {]} \\
Other & \_\_\_\_\_\_ MHz & \_\_\_\_\_\_ MHz & {[} {]} \\
\end{longtable}
}

\subsection{Reset and boot configuration}\label{wiki-hardware-bring-up-and-testing--42-reset-and-boot-configuration}

\begin{itemize}
\tightlist
\item[$\square$]
  NRST idle level correct
\item[$\square$]
  NRST toggles correctly from debugger
\item[$\square$]
  BOOT0 level correct
\item[$\square$]
  BOOT1 level correct
\item[$\square$]
  Development boot can be selected
\item[$\square$]
  Normal flash boot can be selected
\end{itemize}

\section{Debug Access / MCU Bring-Up}\label{wiki-hardware-bring-up-and-testing--5-debug-access--mcu-bring-up}

\begin{itemize}
\tightlist
\item[$\square$]
  ST-Link detects target
\item[$\square$]
  Correct STM32N657 device identified
\item[$\square$]
  Core can be halted
\item[$\square$]
  Registers can be read
\item[$\square$]
  SRAM can be written/read
\item[$\square$]
  Reset + reconnect works reliably
\item[$\square$]
  Minimal firmware reaches \texttt{main()}
\item[$\square$]
  System clock configured
\item[$\square$]
  ThreadX tick works
\item[$\square$]
  GPIO test works
\item[$\square$]
  Status LED toggles
\item[$\square$]
  SWO/UART logging works
\end{itemize}

Suggested first log:

\begin{verbatim}
Boot
Clock OK
ThreadX OK
\end{verbatim}

\section{ThreadX / Software Foundation}\label{wiki-hardware-bring-up-and-testing--6-threadx--software-foundation}

\begin{itemize}
\tightlist
\item[$\square$]
  \texttt{tx\_\allowbreak{}app\_\allowbreak{}thread} starts
\item[$\square$]
  NetX application thread starts
\item[$\square$]
  Control thread starts
\item[$\square$]
  Thread stacks are valid
\item[$\square$]
  No immediate stack overflow
\item[$\square$]
  Event flags work
\item[$\square$]
  Threads sleep/wake correctly
\item[$\square$]
  System remains stable for at least 10 minutes
\end{itemize}

{\def\LTcaptype{none} % do not increment counter
\begin{longtable}[]{@{}
  >{\raggedright\arraybackslash}p{(\linewidth - 8\tabcolsep) * \real{0.1978}}
  >{\raggedleft\arraybackslash}p{(\linewidth - 8\tabcolsep) * \real{0.2259}}
  >{\raggedleft\arraybackslash}p{(\linewidth - 8\tabcolsep) * \real{0.2076}}
  >{\raggedleft\arraybackslash}p{(\linewidth - 8\tabcolsep) * \real{0.1875}}
  >{\centering\arraybackslash}p{(\linewidth - 8\tabcolsep) * \real{0.1813}}@{}}
\toprule\noalign{}
\begin{minipage}[b]{\linewidth}\raggedright
\textbf{Thread}
\end{minipage} & \begin{minipage}[b]{\linewidth}\raggedleft
\textbf{Stack size}
\end{minipage} & \begin{minipage}[b]{\linewidth}\raggedleft
\textbf{Max used}
\end{minipage} & \begin{minipage}[b]{\linewidth}\raggedleft
\textbf{Margin}
\end{minipage} & \begin{minipage}[b]{\linewidth}\centering
\textbf{Pass}
\end{minipage} \\
\midrule\noalign{}
\endhead
\bottomrule\noalign{}
\endlastfoot
App & \_\_\_\_\_\_ & \_\_\_\_\_\_ & \_\_\_\_\_\_ & {[} {]} \\
NetX & \_\_\_\_\_\_ & \_\_\_\_\_\_ & \_\_\_\_\_\_ & {[} {]} \\
Control & \_\_\_\_\_\_ & \_\_\_\_\_\_ & \_\_\_\_\_\_ & {[} {]} \\
Stream & \_\_\_\_\_\_ & \_\_\_\_\_\_ & \_\_\_\_\_\_ & {[} {]} \\
\end{longtable}
}

\section{SEMPER Flash}\label{wiki-hardware-bring-up-and-testing--7-semper-flash}

\subsection{Physical interface}\label{wiki-hardware-bring-up-and-testing--71-physical-interface}

\begin{itemize}
\tightlist
\item[$\square$]
  XSPI1 connected through XSPIM Port 2
\item[$\square$]
  NCS activity observed
\item[$\square$]
  CLK activity observed
\item[$\square$]
  IO0/IO1 activity observed
\item[$\square$]
  DQS activity later in OPI/DTR mode
\end{itemize}

\subsection{Legacy SPI access}\label{wiki-hardware-bring-up-and-testing--72-legacy-spi-access}

\begin{itemize}
\tightlist
\item[$\square$]
  Reset sequence succeeds
\item[$\square$]
  Flash becomes ready
\item[$\square$]
  JEDEC/device ID correct
\item[$\square$]
  CFR registers readable
\item[$\square$]
  Factory/hybrid geometry matches expectation
\end{itemize}

JEDEC/device ID: \_\_\_\_\_\_\_\_\_\_\_\_\_\_\_\_\_\_\_\_\_\_

\subsection{Mode switching}\label{wiki-hardware-bring-up-and-testing--73-mode-switching}

\begin{itemize}
\tightlist
\item[$\square$]
  SPI STR works
\item[$\square$]
  CFR2 latency configuration verified
\item[$\square$]
  CFR3 latency configuration verified
\item[$\square$]
  CFR5 reserved bits preserved correctly
\item[$\square$]
  Enter OPI STR works, if used
\item[$\square$]
  Enter OPI DTR works
\item[$\square$]
  CFR5 readback matches requested mode
\item[$\square$]
  Exit OPI returns reliably to SPI STR
\item[$\square$]
  Repeated mode transitions work
\end{itemize}

\subsection{Basic flash operations}\label{wiki-hardware-bring-up-and-testing--74-basic-flash-operations}

Use a sacrificial test region only.

\begin{itemize}
\tightlist
\item[$\square$]
  Erased read returns \texttt{0xFF}
\item[$\square$]
  256-byte page program works
\item[$\square$]
  Readback matches
\item[$\square$]
  Cross-page write splitting works
\item[$\square$]
  4 KiB erase works in parameter region
\item[$\square$]
  256 KiB erase works in main region
\item[$\square$]
  Busy polling works
\item[$\square$]
  Program-error handling works
\item[$\square$]
  Erase-error handling works
\item[$\square$]
  Timeout handling works
\end{itemize}

\section{ExtMemLoader}\label{wiki-hardware-bring-up-and-testing--8-extmemloader}

\begin{itemize}
\tightlist
\item[$\square$]
  \texttt{.\allowbreak{}stldr} generated
\item[$\square$]
  STM32CubeProgrammer loads loader
\item[$\square$]
  \texttt{StorageInfo} recognized
\item[$\square$]
  Base address = \texttt{0x90000000}
\item[$\square$]
  Device size = 64 MiB
\item[$\square$]
  Geometry reported correctly
\end{itemize}

Expected geometry:

\begin{verbatim}
32 × 4 KiB
1 × 128 KiB
255 × 256 KiB
\end{verbatim}

\begin{itemize}
\tightlist
\item[$\square$]
  Read external flash
\item[$\square$]
  Program test pattern
\item[$\square$]
  Verify test pattern
\item[$\square$]
  Sector erase works
\item[$\square$]
  Full address range accessible
\item[$\square$]
  Invalid addresses rejected, not wrapped
\end{itemize}

\section{Boot Chain}\label{wiki-hardware-bring-up-and-testing--9-boot-chain}

\subsection{Guard / parameter region}\label{wiki-hardware-bring-up-and-testing--91-guard--parameter-region}

\begin{itemize}
\tightlist
\item[$\square$]
  \texttt{0x00000000–0x00000FFF} remains invalid as FSBL1
\item[$\square$]
  Parameter sectors begin at \texttt{0x00001000}
\item[$\square$]
  BootROM rejects offset \texttt{0x00000000} as FSBL1
\end{itemize}

\subsection{FSBL2 boot}\label{wiki-hardware-bring-up-and-testing--92-fsbl2-boot}

\begin{itemize}
\tightlist
\item[$\square$]
  Signed FSBL programmed at flash offset \texttt{0x00040000}
\item[$\square$]
  CPU address corresponds to \texttt{0x90040000}
\item[$\square$]
  BootROM finds FSBL2
\item[$\square$]
  FSBL executes correctly
\item[$\square$]
  FSBL configures SEMPER runtime mode
\end{itemize}

\subsection{Application A}\label{wiki-hardware-bring-up-and-testing--93-application-a}

\begin{itemize}
\tightlist
\item[$\square$]
  Application A programmed at \texttt{0x00100000}
\item[$\square$]
  FSBL loads Application A into SRAM
\item[$\square$]
  Vector table correct
\item[$\square$]
  MSP/VTOR transition correct
\item[$\square$]
  Cold boot reliable
\item[$\square$]
  Warm reset reliable
\end{itemize}

\section{Ethernet PHY Bring-Up}\label{wiki-hardware-bring-up-and-testing--10-ethernet-phy-bring-up}

\subsection{MDIO}\label{wiki-hardware-bring-up-and-testing--101-mdio}

\begin{itemize}
\tightlist
\item[$\square$]
  PHY reset works
\item[$\square$]
  PHY ID1 readable
\item[$\square$]
  PHY ID2 readable
\item[$\square$]
  BMCR readable
\item[$\square$]
  BMSR readable
\item[$\square$]
  Extended registers readable/writable
\item[$\square$]
  ADIN1300 configuration verified
\end{itemize}

\subsection{Physical link}\label{wiki-hardware-bring-up-and-testing--102-physical-link}

With cable disconnected:

\begin{itemize}
\tightlist
\item[$\square$]
  Link reports DOWN
\end{itemize}

Plug cable:

\begin{itemize}
\tightlist
\item[$\square$]
  Link reports UP
\item[$\square$]
  1 Gbit/s reported when expected
\item[$\square$]
  Full duplex reported
\item[$\square$]
  No repeated renegotiation
\item[$\square$]
  LEDs behave as expected
\end{itemize}

Remove cable:

\begin{itemize}
\tightlist
\item[$\square$]
  Link reports DOWN
\item[$\square$]
  Application keeps running
\end{itemize}

Reconnect:

\begin{itemize}
\tightlist
\item[$\square$]
  Link returns UP without reset
\end{itemize}

\section{NetX Duo Basic Networking}\label{wiki-hardware-bring-up-and-testing--11-netx-duo-basic-networking}

Use static IP first.

Board IP: \_\_\_\_\_\_\_\_\_\_\_\_\_\_\_\_\_\_\_\_\_\_\\
Host IP: \_\_\_\_\_\_\_\_\_\_\_\_\_\_\_\_\_\_\_\_\_\_\\
Netmask: \_\_\_\_\_\_\_\_\_\_\_\_\_\_\_\_\_\_\_\_\_\_

\begin{itemize}
\tightlist
\item[$\square$]
  \texttt{NX\_\allowbreak{}IP\_\allowbreak{}LINK\_\allowbreak{}ENABLED} follows cable state
\item[$\square$]
  \texttt{APP\_\allowbreak{}EVT\_\allowbreak{}ETH\_\allowbreak{}LINK\_\allowbreak{}UP} follows cable state
\item[$\square$]
  \texttt{APP\_\allowbreak{}EVT\_\allowbreak{}ETH\_\allowbreak{}IP\_\allowbreak{}READY} follows usable network state
\item[$\square$]
  ARP works
\item[$\square$]
  ICMP works
\item[$\square$]
  Board responds to ping
\item[$\square$]
  Replugging cable restores ping automatically
\item[$\square$]
  100 cable unplug/replug cycles without reset
\item[$\square$]
  30 min continuous ping without failure
\end{itemize}

\section{TCP Control Channel}\label{wiki-hardware-bring-up-and-testing--12-tcp-control-channel}

\subsection{TCP echo}\label{wiki-hardware-bring-up-and-testing--121-tcp-echo}

\begin{itemize}
\tightlist
\item[$\square$]
  TCP server listens
\item[$\square$]
  PC connects
\item[$\square$]
  STM32 receives data
\item[$\square$]
  Echo returned
\item[$\square$]
  Client disconnect handled
\item[$\square$]
  Reconnect handled
\item[$\square$]
  Cable disconnect during active session handled
\item[$\square$]
  Cable reconnect allows new session
\end{itemize}

\subsection{Control protocol}\label{wiki-hardware-bring-up-and-testing--122-control-protocol}

\begin{itemize}
\tightlist
\item[$\square$]
  \texttt{PING}
\item[$\square$]
  \texttt{GET\_\allowbreak{}DEVICE\_\allowbreak{}INFO}
\item[$\square$]
  \texttt{GET\_\allowbreak{}STATUS}
\end{itemize}

\texttt{GET\_\allowbreak{}DEVICE\_\allowbreak{}INFO} should eventually expose:

\begin{itemize}
\tightlist
\item
  firmware version
\item
  Git hash
\item
  board serial
\item
  hardware revision
\item
  MAC
\item
  protocol version
\end{itemize}

\texttt{GET\_\allowbreak{}STATUS} should expose:

\begin{itemize}
\tightlist
\item
  uptime
\item
  Ethernet state
\item
  SEMPER state
\item
  IMU state
\item
  ToF state
\item
  USB state
\item
  error flags
\item
  stream counters
\end{itemize}

\section{USB-C / USB-PD}\label{wiki-hardware-bring-up-and-testing--13-usb-c--usb-pd}

\subsection{PD}\label{wiki-hardware-bring-up-and-testing--131-pd}

\begin{itemize}
\tightlist
\item[$\square$]
  Attach detected
\item[$\square$]
  Default 5 V safe
\item[$\square$]
  PD negotiation starts
\item[$\square$]
  Requested PDO/RDO correct
\item[$\square$]
  Target voltage achieved
\item[$\square$]
  Detach detected
\item[$\square$]
  Reattach works
\end{itemize}

Negotiated voltage: \_\_\_\_\_\_ V\\
Negotiated current: \_\_\_\_\_\_ A

\subsection{USB data}\label{wiki-hardware-bring-up-and-testing--132-usb-data}

\begin{itemize}
\tightlist
\item[$\square$]
  USBX initializes
\item[$\square$]
  PCD starts appropriately
\item[$\square$]
  Device enumerates
\item[$\square$]
  CDC ACM works
\item[$\square$]
  USB unplug/replug works
\item[$\square$]
  Ethernet unaffected by USB state
\item[$\square$]
  USB works with Ethernet disconnected
\end{itemize}

\section{Shared Control Layer}\label{wiki-hardware-bring-up-and-testing--14-shared-control-layer}

\begin{itemize}
\tightlist
\item[$\square$]
  Same command decoder used by TCP and USB
\item[$\square$]
  \texttt{PING} works via Ethernet
\item[$\square$]
  \texttt{PING} works via USB
\item[$\square$]
  \texttt{GET\_\allowbreak{}DEVICE\_\allowbreak{}INFO} identical on both
\item[$\square$]
  \texttt{GET\_\allowbreak{}STATUS} identical on both
\item[$\square$]
  Concurrent Ethernet + USB behavior defined
\end{itemize}

\section{IMU -- LSM6DSV16X}\label{wiki-hardware-bring-up-and-testing--15-imu--lsm6dsv16x}

\begin{itemize}
\tightlist
\item[$\square$]
  WHO\_AM\_I correct
\item[$\square$]
  Register read/write works
\item[$\square$]
  Reset works
\item[$\square$]
  Accelerometer starts
\item[$\square$]
  Gyroscope starts
\item[$\square$]
  FIFO works if used
\item[$\square$]
  Interrupt works if used
\item[$\square$]
  Static acceleration ≈ 1 g
\item[$\square$]
  Axis orientation matches board coordinates
\item[$\square$]
  Gyro near zero when stationary
\item[$\square$]
  Sample rate correct
\item[$\square$]
  Sensor fusion/game rotation works
\item[$\square$]
  Timestamping consistent
\end{itemize}

\section{VL53L9 / MIPI CSI-2}\label{wiki-hardware-bring-up-and-testing--16-vl53l9--mipi-csi-2}

\begin{itemize}
\tightlist
\item[$\square$]
  Sensor detected/configured
\item[$\square$]
  Correct operating mode selected
\item[$\square$]
  CSI-2 clock activity present
\item[$\square$]
  Lane activity present
\item[$\square$]
  PHY locks
\item[$\square$]
  Correct virtual channel
\item[$\square$]
  Correct data types
\item[$\square$]
  Correct frame geometry
\item[$\square$]
  First complete frame received
\item[$\square$]
  Frame size correct
\item[$\square$]
  Frame counter increments
\item[$\square$]
  No buffer corruption
\item[$\square$]
  Double buffering works
\item[$\square$]
  Continuous operation stable
\item[$\square$]
  Frame timestamp captured
\end{itemize}

\section{UDP Stream Layer}\label{wiki-hardware-bring-up-and-testing--17-udp-stream-layer}

\subsection{Synthetic stream}\label{wiki-hardware-bring-up-and-testing--171-synthetic-stream}

\begin{itemize}
\tightlist
\item[$\square$]
  UDP socket created
\item[$\square$]
  Synthetic frames generated
\item[$\square$]
  Packet sequence correct
\item[$\square$]
  Frame IDs correct
\item[$\square$]
  Timestamps correct
\item[$\square$]
  Host reconstructs frame
\item[$\square$]
  Lost packet only invalidates affected frame
\item[$\square$]
  Link loss stops stream cleanly
\item[$\square$]
  Reconnect allows new stream
\end{itemize}

\subsection{Real depth stream}\label{wiki-hardware-bring-up-and-testing--172-real-depth-stream}

\begin{itemize}
\tightlist
\item[$\square$]
  VL53L9 frame sent over UDP
\item[$\square$]
  Host reconstructs frame
\item[$\square$]
  Frame rate correct
\item[$\square$]
  End-to-end latency measured
\item[$\square$]
  Packet-loss counters correct
\item[$\square$]
  Frame-drop counters correct
\item[$\square$]
  Acquisition never blocks indefinitely on TX
\end{itemize}

\section{Configuration Storage}\label{wiki-hardware-bring-up-and-testing--18-configuration-storage}

\subsection{Factory information}\label{wiki-hardware-bring-up-and-testing--181-factory-information}

\begin{itemize}
\tightlist
\item[$\square$]
  Serial number writable
\item[$\square$]
  HW revision writable
\item[$\square$]
  MAC writable
\item[$\square$]
  Calibration writable
\item[$\square$]
  Factory CRC validated
\item[$\square$]
  Normal application does not overwrite factory data
\end{itemize}

\subsection{Device configuration}\label{wiki-hardware-bring-up-and-testing--182-device-configuration}

\begin{itemize}
\tightlist
\item[$\square$]
  Config A valid
\item[$\square$]
  Config B valid
\item[$\square$]
  Generation counter works
\item[$\square$]
  CRC works
\item[$\square$]
  Latest valid copy selected
\item[$\square$]
  Power-loss-safe update verified
\item[$\square$]
  Static IP persists
\item[$\square$]
  DHCP flag persists
\item[$\square$]
  Hostname persists
\end{itemize}

\section{Future A/B Update / FoE}\label{wiki-hardware-bring-up-and-testing--19-future-ab-update--foe}

\begin{itemize}
\tightlist
\item[$\square$]
  Application A and B regions independently erasable
\item[$\square$]
  Parameter region excluded from normal firmware update
\item[$\square$]
  FSBL excluded from normal application update
\item[$\square$]
  Full-chip erase never used by normal updater
\item[$\square$]
  Update metadata region defined
\end{itemize}

Later:

\begin{itemize}
\tightlist
\item[$\square$]
  FoE downloads to inactive slot
\item[$\square$]
  Image verification
\item[$\square$]
  Pending slot flag
\item[$\square$]
  Pending image boot
\item[$\square$]
  Successful image confirms itself
\item[$\square$]
  Failed image rolls back
\item[$\square$]
  Power loss during download preserves current firmware
\end{itemize}

\section{Fault / Recovery}\label{wiki-hardware-bring-up-and-testing--20-fault--recovery}

\subsection{Power}\label{wiki-hardware-bring-up-and-testing--power}

\begin{itemize}
\tightlist
\item[$\square$]
  Power cycle during idle
\item[$\square$]
  Power cycle during Ethernet traffic
\item[$\square$]
  Controlled power loss during config write
\item[$\square$]
  Brownout behavior understood
\item[$\square$]
  No unrelated flash corruption
\end{itemize}

\subsection{Reset}\label{wiki-hardware-bring-up-and-testing--reset}

\begin{itemize}
\tightlist
\item[$\square$]
  NRST during normal operation
\item[$\square$]
  Software reset during normal operation
\item[$\square$]
  Reset while SEMPER is in OPI DTR
\item[$\square$]
  Flash recovery works after MCU-only reset
\end{itemize}

\subsection{Communication}\label{wiki-hardware-bring-up-and-testing--communication}

\begin{itemize}
\tightlist
\item[$\square$]
  Ethernet removed during TCP
\item[$\square$]
  Ethernet removed during UDP
\item[$\square$]
  USB removed during control
\item[$\square$]
  Ethernet + USB simultaneously
\item[$\square$]
  Host process crash during TCP
\item[$\square$]
  Device accepts new host afterward
\end{itemize}

\section{Performance}\label{wiki-hardware-bring-up-and-testing--21-performance}

\subsection{Ethernet}\label{wiki-hardware-bring-up-and-testing--ethernet}

TCP throughput: \_\_\_\_\_\_ Mbit/s\\
UDP throughput: \_\_\_\_\_\_ Mbit/s

\begin{itemize}
\tightlist
\item[$\square$]
  Sustained TCP throughput measured
\item[$\square$]
  Sustained UDP throughput measured
\item[$\square$]
  Packet loss measured
\item[$\square$]
  CPU load measured
\item[$\square$]
  Memory usage measured
\end{itemize}

\subsection{Sensor streaming}\label{wiki-hardware-bring-up-and-testing--sensor-streaming}

Frame rate: \_\_\_\_\_\_ Hz\\
Mean latency: \_\_\_\_\_\_ ms\\
Worst-case latency: \_\_\_\_\_\_ ms\\
Dropped frames: \_\_\_\_\_\_ \%

\begin{itemize}
\tightlist
\item[$\square$]
  Maximum sustainable frame rate measured
\item[$\square$]
  Latency measured
\item[$\square$]
  Jitter measured
\item[$\square$]
  Drop rate measured
\end{itemize}

\subsection{Flash}\label{wiki-hardware-bring-up-and-testing--flash}

\begin{itemize}
\tightlist
\item[$\square$]
  Read throughput measured
\item[$\square$]
  Program throughput measured
\item[$\square$]
  Memory-mapped read throughput measured
\item[$\square$]
  OPI DTR stable at target clock
\end{itemize}

\section{Thermal / Full-Load Test}\label{wiki-hardware-bring-up-and-testing--22-thermal--full-load-test}

Run:

\begin{verbatim}
Ethernet active
+ camera streaming
+ IMU active
+ USB active
+ flash traffic
\end{verbatim}

Duration: \_\_\_\_\_\_ min

{\def\LTcaptype{none} % do not increment counter
\begin{longtable}[]{@{}
  >{\raggedright\arraybackslash}p{(\linewidth - 6\tabcolsep) * \real{0.3015}}
  >{\raggedleft\arraybackslash}p{(\linewidth - 6\tabcolsep) * \real{0.2493}}
  >{\raggedleft\arraybackslash}p{(\linewidth - 6\tabcolsep) * \real{0.2586}}
  >{\centering\arraybackslash}p{(\linewidth - 6\tabcolsep) * \real{0.1906}}@{}}
\toprule\noalign{}
\begin{minipage}[b]{\linewidth}\raggedright
\textbf{Component}
\end{minipage} & \begin{minipage}[b]{\linewidth}\raggedleft
\textbf{Temperature}
\end{minipage} & \begin{minipage}[b]{\linewidth}\raggedleft
\textbf{Limit/target}
\end{minipage} & \begin{minipage}[b]{\linewidth}\centering
\textbf{Pass}
\end{minipage} \\
\midrule\noalign{}
\endhead
\bottomrule\noalign{}
\endlastfoot
STM32N657 & \_\_\_\_\_\_ °C & \_\_\_\_\_\_ °C & {[} {]} \\
ADIN1300 & \_\_\_\_\_\_ °C & \_\_\_\_\_\_ °C & {[} {]} \\
SEMPER & \_\_\_\_\_\_ °C & \_\_\_\_\_\_ °C & {[} {]} \\
Main regulator & \_\_\_\_\_\_ °C & \_\_\_\_\_\_ °C & {[} {]} \\
USB/PD power path & \_\_\_\_\_\_ °C & \_\_\_\_\_\_ °C & {[} {]} \\
\end{longtable}
}

\begin{itemize}
\tightlist
\item[$\square$]
  No reset
\item[$\square$]
  No thermal throttling
\item[$\square$]
  Ethernet stable
\item[$\square$]
  Sensor data stable
\end{itemize}

\section{Long-Duration Stability}\label{wiki-hardware-bring-up-and-testing--23-long-duration-stability}

Suggested first target:

\begin{verbatim}
24 h continuous operation
\end{verbatim}

Run with:

\begin{itemize}
\tightlist
\item
  camera streaming
\item
  IMU active
\item
  Ethernet connected
\item
  periodic control requests
\item
  periodic configuration reads
\end{itemize}

Track:

Runtime: \_\_\_\_\_\_\_\_\_\_\_\_\_\_\_\_\_\_\\
Resets: \_\_\_\_\_\_\_\_\_\_\_\_\_\_\_\_\_\_\_\\
Dropped frames: \_\_\_\_\_\_\_\_\_\_\_\\
Network errors: \_\_\_\_\_\_\_\_\_\_\_\\
Sensor errors: \_\_\_\_\_\_\_\_\_\_\_\_\\
Flash errors: \_\_\_\_\_\_\_\_\_\_\_\_\_

\begin{itemize}
\tightlist
\item[$\square$]
  No unexpected reset
\item[$\square$]
  No packet-pool exhaustion
\item[$\square$]
  No thread errors
\item[$\square$]
  No sensor lockups
\item[$\square$]
  No flash errors
\end{itemize}

\section{Final Acceptance}\label{wiki-hardware-bring-up-and-testing--24-final-acceptance}

\begin{itemize}
\tightlist
\item[$\square$]
  Power rails stable
\item[$\square$]
  Debug/recovery works
\item[$\square$]
  Autonomous SEMPER boot works
\item[$\square$]
  Parameter storage works
\item[$\square$]
  Ethernet may be connected after boot
\item[$\square$]
  Ethernet reconnect works
\item[$\square$]
  TCP control works
\item[$\square$]
  USB control works
\item[$\square$]
  IMU works
\item[$\square$]
  VL53L9 acquisition works
\item[$\square$]
  UDP depth stream works
\item[$\square$]
  Configuration survives reset
\item[$\square$]
  Communication interruptions recover cleanly
\item[$\square$]
  Thermal test passes
\item[$\square$]
  Long-duration test passes
\end{itemize}

\section{Issue Log}\label{wiki-hardware-bring-up-and-testing--25-issue-log}

{\def\LTcaptype{none} % do not increment counter
\begin{longtable}[]{@{}
  >{\raggedright\arraybackslash}p{(\linewidth - 10\tabcolsep) * \real{0.1501}}
  >{\raggedright\arraybackslash}p{(\linewidth - 10\tabcolsep) * \real{0.1501}}
  >{\raggedright\arraybackslash}p{(\linewidth - 10\tabcolsep) * \real{0.1853}}
  >{\raggedright\arraybackslash}p{(\linewidth - 10\tabcolsep) * \real{0.1675}}
  >{\raggedright\arraybackslash}p{(\linewidth - 10\tabcolsep) * \real{0.1795}}
  >{\centering\arraybackslash}p{(\linewidth - 10\tabcolsep) * \real{0.1675}}@{}}
\toprule\noalign{}
\begin{minipage}[b]{\linewidth}\raggedright
\textbf{ID}
\end{minipage} & \begin{minipage}[b]{\linewidth}\raggedright
\textbf{Stage}
\end{minipage} & \begin{minipage}[b]{\linewidth}\raggedright
\textbf{Description}
\end{minipage} & \begin{minipage}[b]{\linewidth}\raggedright
\textbf{Severity}
\end{minipage} & \begin{minipage}[b]{\linewidth}\raggedright
\textbf{Resolution}
\end{minipage} & \begin{minipage}[b]{\linewidth}\centering
\textbf{Retested}
\end{minipage} \\
\midrule\noalign{}
\endhead
\bottomrule\noalign{}
\endlastfoot
001 & & & & & {[} {]} \\
002 & & & & & {[} {]} \\
003 & & & & & {[} {]} \\
004 & & & & & {[} {]} \\
\end{longtable}
}

\section{Short-Form Bring-Up Checklist}\label{wiki-hardware-bring-up-and-testing--26-short-form-bring-up-checklist}

\begin{itemize}
\tightlist
\item[$\square$]
  \textbf{1. Visual inspection}
\item[$\square$]
  \textbf{2. Resistance / short checks}
\item[$\square$]
  \textbf{3. Current-limited first power-up}
\item[$\square$]
  \textbf{4. Verify all rails}
\item[$\square$]
  \textbf{5. Verify clocks / reset / BOOT}
\item[$\square$]
  \textbf{6. Establish SWD/SWO}
\item[$\square$]
  \textbf{7. Minimal ThreadX firmware}
\item[$\square$]
  \textbf{8. SEMPER legacy SPI}
\item[$\square$]
  \textbf{9. SEMPER read/write/erase}
\item[$\square$]
  \textbf{10. ExtMemLoader}
\item[$\square$]
  \textbf{11. BootROM → FSBL2}
\item[$\square$]
  \textbf{12. FSBL → Application A}
\item[$\square$]
  \textbf{13. ADIN1300 MDIO}
\item[$\square$]
  \textbf{14. Ethernet link}
\item[$\square$]
  \textbf{15. Ping}
\item[$\square$]
  \textbf{16. TCP echo/control}
\item[$\square$]
  \textbf{17. Ethernet unplug/replug}
\item[$\square$]
  \textbf{18. USB-PD}
\item[$\square$]
  \textbf{19. USB CDC}
\item[$\square$]
  \textbf{20. IMU}
\item[$\square$]
  \textbf{21. VL53L9 / CSI-2}
\item[$\square$]
  \textbf{22. Synthetic UDP stream}
\item[$\square$]
  \textbf{23. Real depth stream}
\item[$\square$]
  \textbf{24. Persistent configuration}
\item[$\square$]
  \textbf{25. Fault/recovery}
\item[$\square$]
  \textbf{26. Thermal/load}
\item[$\square$]
  \textbf{27. 24-hour soak}
\item[$\square$]
  \textbf{28. Full-system acceptance}
\end{itemize}

\appendix

\chapter{Library documentation and source map}\label{wiki-reference-dependencies--library-documentation-and-source-map}

{\small\color{MGSMuted}Source: \nolinkurl{wiki/reference/dependencies.md}\par}\medskip

\section{Firmware contexts}\label{wiki-reference-dependencies--firmware-contexts}

{\def\LTcaptype{none} % do not increment counter
\begin{longtable}[]{@{}
  >{\raggedright\arraybackslash}p{(\linewidth - 2\tabcolsep) * \real{0.4100}}
  >{\raggedright\arraybackslash}p{(\linewidth - 2\tabcolsep) * \real{0.5900}}@{}}
\toprule\noalign{}
\begin{minipage}[b]{\linewidth}\raggedright
\textbf{Location}
\end{minipage} & \begin{minipage}[b]{\linewidth}\raggedright
\textbf{Purpose}
\end{minipage} \\
\midrule\noalign{}
\endhead
\bottomrule\noalign{}
\endlastfoot
{Appli (\nolinkurl{Appli})} & Main ThreadX application, USB, Ethernet, sensors, and control services \\
{FSBL (\nolinkurl{FSBL})} & Bootloader, image validation, SRAM loading, and A/B selection \\
{ExtMemLoader (\nolinkurl{ExtMemLoader})} & CubeProgrammer external flash loader \\
{Drivers (\nolinkurl{Drivers})} & STM32 HAL/LL and CMSIS drivers \\
{Middlewares (\nolinkurl{Middlewares})} & ThreadX, NetX Duo, USBX, USB-PD, and external memory manager \\
\end{longtable}
}

\section{Project implementation}\label{wiki-reference-dependencies--project-implementation}

{\def\LTcaptype{none} % do not increment counter
\begin{longtable}[]{@{}
  >{\raggedright\arraybackslash}p{(\linewidth - 2\tabcolsep) * \real{0.4497}}
  >{\raggedright\arraybackslash}p{(\linewidth - 2\tabcolsep) * \real{0.5503}}@{}}
\toprule\noalign{}
\begin{minipage}[b]{\linewidth}\raggedright
\textbf{Location}
\end{minipage} & \begin{minipage}[b]{\linewidth}\raggedright
\textbf{Purpose}
\end{minipage} \\
\midrule\noalign{}
\endhead
\bottomrule\noalign{}
\endlastfoot
{Firmware update policy (\nolinkurl{common/firmware})} & Shared image validation, flash metadata, and boot policy \\
{SEMPER driver (\nolinkurl{common/s28hs512t})} & S28HS512T SPI/OPI access \\
{Application services (\nolinkurl{Appli/Libs/control})} & Streaming, control, update server, profiling, and parameter reporting \\
{VL53L9 integration (\nolinkurl{Appli/Libs/53L9})} & Sensor application, interface, and firmware library \\
{IMU integration (\nolinkurl{Appli/Libs/IMU})} & LSM6DSV16X acquisition \\
{Ethernet PHY (\nolinkurl{Appli/Libs/PHY})} & ADIN1300 integration \\
{Build and programming scripts (\nolinkurl{scripts})} & Builds, manifests, parameters, SWD programming, and DFU uploader \\
\end{longtable}
}

\section{Bundled library documentation}\label{wiki-reference-dependencies--bundled-library-documentation}

Vendor documents remain with their libraries; the wiki links to them rather than maintaining independent copies or changing submodule contents.

\begin{itemize}
\tightlist
\item
  LSM6DSV16X register driver: {README (\nolinkurl{Appli/Libs/lsm6dsv16x-pid/README.md})}, {release notes (\nolinkurl{Appli/Libs/lsm6dsv16x-pid/Release_Notes.md})}, {contributing (\nolinkurl{Appli/Libs/lsm6dsv16x-pid/CONTRIBUTING.md})}, {security policy (\nolinkurl{Appli/Libs/lsm6dsv16x-pid/SECURITY.md})}, and {code of conduct (\nolinkurl{Appli/Libs/lsm6dsv16x-pid/CODE_OF_CONDUCT.md})}.
\item
  VL53L9 sensor library: {changelog (\nolinkurl{Appli/Libs/53L9/vl53l9/CHANGELOG.md})} and {license (\nolinkurl{Appli/Libs/53L9/vl53l9/LICENSE.txt})}.
\item
  VL53L9 integration licenses: {application (\nolinkurl{Appli/Libs/53L9/app/LICENSE.txt})} and {interface (\nolinkurl{Appli/Libs/53L9/interface/LICENSE.txt})}.
\end{itemize}

The LSM6DSV16X driver is a Git submodule. Initialize it after cloning if its files are missing:

\begin{verbatim}
git submodule update --init --recursive
\end{verbatim}

\section{Desktop tools and tests}\label{wiki-reference-dependencies--desktop-tools-and-tests}

\begin{itemize}
\tightlist
\item
  {Device Monitor (\nolinkurl{test_scripts/Board_Viewer.py})}
\item
  {Device Monitor firmware update worker (\nolinkurl{test_scripts/Firmware_Update.py})}
\item
  {USB stream decoder (\nolinkurl{test_scripts/USB_Debug.py})}
\item
  {Ethernet viewer (\nolinkurl{test_scripts/Ethernet_Debug.py})}
\item
  {Sensor command client (\nolinkurl{test_scripts/Sensor_Control.py})}
\item
  {Firmware uploader (\nolinkurl{scripts/firmware_update.py})}
\item
  {Scene reconstruction (\nolinkurl{test_scripts/Reconstruct_Scene.py})}
\item
  {Test scripts (\nolinkurl{test_scripts})}
\end{itemize}

Run host regressions from the repository root:

\begin{verbatim}
python3 -m unittest discover -s test_scripts -p 'test_*.py'
\end{verbatim}

GUI integration tests require a working display and the dependencies listed in \hyperref[wiki-tools-device-monitor--installation]{Device Monitor}. Hardware acceptance is recorded separately in the \hyperref[wiki-hardware-bring-up-and-testing--magnecko-sense-mgs--bring-up--system-test-protocol]{bring-up plan} and \hyperref[wiki-firmware-firmware-updates--validation]{firmware-update validation guide}.

\chapter{Historical flash-layout proposal}\label{wiki-reference-flash-layout-proposal--stm32n6--s28hs512t-external-flash-memory-map}

{\small\color{MGSMuted}Source: \nolinkurl{wiki/reference/flash-layout-proposal.md}\par}\medskip

\begin{quote}
\textbf{Design reference:} this preserves the original proposed layout. The implemented parameter block and DFU metadata allocation differ from parts of this proposal. Use the \hyperref[wiki-hardware-flash-layout--current-external-flash-layout]{current flash layout} for programming addresses and the \hyperref[wiki-firmware-firmware-updates--application-firmware-updates-and-ab-recovery]{firmware update guide} for A/B behavior.
\end{quote}

\section{Address conventions}\label{wiki-reference-flash-layout-proposal--1-address-conventions}

The S28HS512T has a capacity of:

\begin{verbatim}
64 MiB = 0x04000000 bytes
\end{verbatim}

On STM32N657, the XSPI1 memory-mapped bank starts at:

\begin{verbatim}
XSPI1_BASE = 0x90000000
\end{verbatim}

Therefore:

\begin{verbatim}
Flash offset 0x00000000 -> CPU address 0x90000000
Flash offset 0x00100000 -> CPU address 0x90100000
...
Flash offset 0x03FFFFFF -> CPU address 0x93FFFFFF
\end{verbatim}

The ExtMemLoader should advertise the complete device as:

\begin{verbatim}
#define STM32EXTLOADER_DEVICE_ADDR  0x90000000U
#define STM32EXTLOADER_DEVICE_SIZE  0x04000000U
\end{verbatim}

CubeProgrammer addresses are absolute addresses in this range.\\
The SEMPER driver itself should normally operate on flash-relative offsets.

A safe address conversion is therefore:

\begin{verbatim}
if ((Address < STM32EXTLOADER_DEVICE_ADDR) ||
    (Address >= (STM32EXTLOADER_DEVICE_ADDR + STM32EXTLOADER_DEVICE_SIZE)))
{
    return MEM_FAIL;
}

uint32_t flash_offset = Address - STM32EXTLOADER_DEVICE_ADDR;
\end{verbatim}

Do \textbf{not} use only:

\begin{verbatim}
Address & 0x03FFFFFFU
\end{verbatim}

because invalid addresses would silently wrap into the flash.

\section{Physical erase geometry}\label{wiki-reference-flash-layout-proposal--2-physical-erase-geometry}

The intended device configuration keeps the factory-style \textbf{bottom hybrid sector architecture}.

\begin{verbatim}
Flash offset
0x00000000 ┌─────────────────────────────────────────────┐
           │ 32 × 4 KiB parameter sectors               │
0x00020000 ├─────────────────────────────────────────────┤
           │ 1 × 128 KiB transition region              │
0x00040000 ├─────────────────────────────────────────────┤
           │ 255 × 256 KiB main sectors                 │
           │                                             │
0x04000000 └─────────────────────────────────────────────┘
\end{verbatim}

Geometry:

{\def\LTcaptype{none} % do not increment counter
\begin{longtable}[]{@{}
  >{\raggedright\arraybackslash}p{(\linewidth - 6\tabcolsep) * \real{0.2734}}
  >{\raggedleft\arraybackslash}p{(\linewidth - 6\tabcolsep) * \real{0.3132}}
  >{\raggedleft\arraybackslash}p{(\linewidth - 6\tabcolsep) * \real{0.1772}}
  >{\raggedleft\arraybackslash}p{(\linewidth - 6\tabcolsep) * \real{0.2362}}@{}}
\toprule\noalign{}
\begin{minipage}[b]{\linewidth}\raggedright
\textbf{Region}
\end{minipage} & \begin{minipage}[b]{\linewidth}\raggedleft
\textbf{Offset range}
\end{minipage} & \begin{minipage}[b]{\linewidth}\raggedleft
\textbf{Count}
\end{minipage} & \begin{minipage}[b]{\linewidth}\raggedleft
\textbf{Logical size}
\end{minipage} \\
\midrule\noalign{}
\endhead
\bottomrule\noalign{}
\endlastfoot
Parameter sectors & \texttt{0x00000000} -- \texttt{0x0001FFFF} & 32 & 4 KiB \\
Transition region & \texttt{0x00020000} -- \texttt{0x0003FFFF} & 1 & 128 KiB \\
Main sectors & \texttt{0x00040000} -- \texttt{0x03FFFFFF} & 255 & 256 KiB \\
\end{longtable}
}

Total:

\begin{verbatim}
32 × 4 KiB
+ 1 × 128 KiB
+ 255 × 256 KiB
= 64 MiB
\end{verbatim}

The ExtMemLoader \texttt{StorageInfo} sector description must match this geometry.

\section{Proposed logical flash layout}\label{wiki-reference-flash-layout-proposal--3-proposed-logical-flash-layout}

\begin{verbatim}
Flash offset       CPU address       Purpose
──────────────────────────────────────────────────────────────────

0x00000000         0x90000000
    │
    │  4 KiB       GUARD / INVALID FSBL1
    │
0x00001000         0x90001000
    │
    │  124 KiB     PARAMETER STORAGE
    │               - factory identity
    │               - serial number
    │               - MAC address
    │               - calibration
    │               - IP/network configuration
    │               - boot/update state
    │
0x00020000         0x90020000
    │
    │  128 KiB     RESERVED TRANSITION REGION
    │
0x00040000         0x90040000
    │
    │  512 KiB     FSBL2 SLOT
    │
0x000C0000         0x900C0000
    │
    │  256 KiB     RESERVED
    │
0x00100000         0x90100000
    │
    │  8 MiB       APPLICATION SLOT A
    │
0x00900000         0x90900000
    │
    │  8 MiB       APPLICATION SLOT B
    │
0x01100000         0x91100000
    │
    │  47 MiB      FUTURE / DATA / MODELS / LOGGING
    │
0x04000000         0x94000000 (exclusive)
\end{verbatim}

\section{Guard sector / FSBL1}\label{wiki-reference-flash-layout-proposal--4-guard-sector--fsbl1}

The STM32N6 BootROM searches serial NOR for:

\begin{verbatim}
FSBL1 at flash offset 0x00000000
FSBL2 at flash offset 0x00040000
\end{verbatim}

This design intentionally uses \textbf{FSBL2 as the primary FSBL} so that the hybrid parameter-sector region can be used for device data.

The first 4 KiB sector is therefore reserved:

\begin{verbatim}
0x00000000 – 0x00000FFF
\end{verbatim}

Recommended invariant:

\begin{quote}
This sector must never contain a valid STM32 FSBL image header.
\end{quote}

The simplest policy is to leave it erased (\texttt{0xFF}) and never use it for configuration data.

This ensures that the BootROM rejects FSBL1 and proceeds to the valid FSBL stored at \texttt{0x00040000}.

\section{Parameter sector allocation}\label{wiki-reference-flash-layout-proposal--5-parameter-sector-allocation}

The remaining 31 small sectors provide:

\begin{verbatim}
31 × 4 KiB = 124 KiB
\end{verbatim}

A suggested initial allocation is:

{\def\LTcaptype{none} % do not increment counter
\begin{longtable}[]{@{}
  >{\raggedleft\arraybackslash}p{(\linewidth - 4\tabcolsep) * \real{0.3655}}
  >{\raggedleft\arraybackslash}p{(\linewidth - 4\tabcolsep) * \real{0.2171}}
  >{\raggedright\arraybackslash}p{(\linewidth - 4\tabcolsep) * \real{0.4174}}@{}}
\toprule\noalign{}
\begin{minipage}[b]{\linewidth}\raggedleft
\textbf{Offset}
\end{minipage} & \begin{minipage}[b]{\linewidth}\raggedleft
\textbf{Size}
\end{minipage} & \begin{minipage}[b]{\linewidth}\raggedright
\textbf{Purpose}
\end{minipage} \\
\midrule\noalign{}
\endhead
\bottomrule\noalign{}
\endlastfoot
\texttt{0x00001000} & 4 KiB & Factory information A \\
\texttt{0x00002000} & 4 KiB & Factory information B / backup \\
\texttt{0x00003000} & 4 KiB & Device configuration A \\
\texttt{0x00004000} & 4 KiB & Device configuration B \\
\texttt{0x00005000} & 4 KiB & Firmware/update state A \\
\texttt{0x00006000} & 4 KiB & Firmware/update state B \\
\texttt{0x00007000} -- \texttt{0x0001FFFF} & 100 KiB & Reserved / future parameters \\
\end{longtable}
}

\subsection{Factory information}\label{wiki-reference-flash-layout-proposal--factory-information}

Typical contents:

\begin{itemize}
\tightlist
\item
  board serial number
\item
  hardware revision
\item
  production information
\item
  MAC address
\item
  calibration constants
\item
  manufacturing/test data
\end{itemize}

Factory information should normally be treated as read-only by the application.

\subsection{Device configuration}\label{wiki-reference-flash-layout-proposal--device-configuration}

Typical contents:

\begin{itemize}
\tightlist
\item
  DHCP/static-IP selection
\item
  static IP address
\item
  subnet mask
\item
  gateway
\item
  hostname
\item
  EtherCAT/Ethernet settings
\item
  sensor settings
\item
  user-configurable behavior
\end{itemize}

Use two sectors (A/B) with:

\begin{itemize}
\tightlist
\item
  format/version field
\item
  generation counter
\item
  payload length
\item
  CRC
\end{itemize}

Update the inactive copy first, verify it, then make it current. This provides power-loss-safe configuration updates.

\subsection{Firmware/update state}\label{wiki-reference-flash-layout-proposal--firmwareupdate-state}

Reserve another A/B pair for future firmware-update metadata:

\begin{itemize}
\tightlist
\item
  confirmed application slot
\item
  pending application slot
\item
  update counter
\item
  image version
\item
  image hash
\item
  boot-attempt state
\item
  rollback state
\item
  CRC
\end{itemize}

This can later support FoE-based A/B updates.

\section{FSBL region}\label{wiki-reference-flash-layout-proposal--6-fsbl-region}

The real FSBL is stored at the BootROM FSBL2 location:

\begin{verbatim}
Flash offset : 0x00040000
CPU address  : 0x90040000
\end{verbatim}

Reserve 512 KiB:

\begin{verbatim}
0x00040000 – 0x000BFFFF
\end{verbatim}

This matches the STM32N6 BootROM maximum FSBL-image size and consumes exactly two 256 KiB main sectors.

The following sector is kept reserved:

\begin{verbatim}
0x000C0000 – 0x000FFFFF
\end{verbatim}

This maintains a clean 1 MiB boundary before Application Slot A.

The FSBL should be considered infrastructure and updated only when necessary.

\section{Application slots}\label{wiki-reference-flash-layout-proposal--7-application-slots}

\subsection{Application A}\label{wiki-reference-flash-layout-proposal--application-a}

\begin{verbatim}
Flash offset : 0x00100000
CPU address  : 0x90100000
Size         : 0x00800000 (8 MiB)
End          : 0x008FFFFF
\end{verbatim}

\subsection{Application B}\label{wiki-reference-flash-layout-proposal--application-b}

\begin{verbatim}
Flash offset : 0x00900000
CPU address  : 0x90900000
Size         : 0x00800000 (8 MiB)
End          : 0x010FFFFF
\end{verbatim}

The large slots are intentional. An LRUN application must fit into available SRAM when executed, so the flash slots are considerably larger than is currently required. The additional space allows future image packaging, metadata, AI/model data, and firmware growth without changing the partition table.

For the initial firmware, the FSBL can continue to load Slot A from:

\begin{verbatim}
LRUN source offset = 0x00100000
\end{verbatim}

A future A/B-aware FSBL can choose dynamically between:

\begin{verbatim}
Slot A = 0x00100000
Slot B = 0x00900000
\end{verbatim}

based on the update-state parameter sectors.

\section{Remaining flash}\label{wiki-reference-flash-layout-proposal--8-remaining-flash}

After the two 8 MiB application slots:

\begin{verbatim}
0x01100000 – 0x03FFFFFF
\end{verbatim}

remains available.

Size:

\begin{verbatim}
47 MiB
\end{verbatim}

Possible future uses:

\begin{itemize}
\tightlist
\item
  neural-network models
\item
  lookup/calibration datasets
\item
  event or diagnostic logging
\item
  crash dumps
\item
  web-interface assets
\item
  sensor recordings
\item
  additional firmware images
\item
  factory test data
\item
  reserved expansion space
\end{itemize}

Do not assign this region permanently until there is a concrete need.

\section{ExtMemLoader configuration}\label{wiki-reference-flash-layout-proposal--9-extmemloader-configuration}

The ExtMemLoader does \textbf{not} enforce the semantic partition table above. It only exposes the physical external memory to STM32CubeProgrammer.

For the custom board using XSPI1, configure:

\begin{verbatim}
External loader start address : 0x90000000
External loader end address   : 0x93FFFFFF
Device size                   : 0x04000000
Page size                     : 0x00000100
Erase value                   : 0xFF
\end{verbatim}

In CubeMX these correspond to the ExtMemLoader parameters, and the generated configuration should contain:

\begin{verbatim}
#define STM32EXTLOADER_DEVICE_ADDR                0x90000000U
#define STM32EXTLOADER_DEVICE_SIZE                0x04000000U

#define STM32EXTLOADER_DEVICE_PAGE_SIZE           0x00000100U
#define STM32EXTLOADER_DEVICE_INITIAL_CONTENT     0xFFU

#define STM32EXTLOADER_DEVICE_4K_SECTOR_NUMBERS   0x20U
#define STM32EXTLOADER_DEVICE_4K_SECTOR_SIZE      0x1000U

#define STM32EXTLOADER_DEVICE_128K_SECTOR_NUMBERS 0x1U
#define STM32EXTLOADER_DEVICE_128K_SECTOR_SIZE    0x20000U

#define STM32EXTLOADER_DEVICE_256K_SECTOR_NUMBERS 0xFFU
#define STM32EXTLOADER_DEVICE_256K_SECTOR_SIZE    0x40000U
\end{verbatim}

The \texttt{StorageInfo} structure in the ExtMemLoader is what STM32CubeProgrammer reads to discover:

\begin{itemize}
\tightlist
\item
  memory start address
\item
  total size
\item
  page size
\item
  erase value
\item
  erase-region geometry
\end{itemize}

The current STM32N6570-DK example advertises \texttt{0x70000000} because it uses XSPI2.\\
For this custom board, SEMPER is on XSPI1, therefore the advertised base must be \texttt{0x90000000}.

\section{ExtMemLoader address translation}\label{wiki-reference-flash-layout-proposal--10-extmemloader-address-translation}

STM32CubeProgrammer calls the loader using absolute external-memory addresses such as:

\begin{verbatim}
0x90040000
0x90100000
0x90900000
\end{verbatim}

The low-level SEMPER driver should receive flash-relative offsets:

\begin{verbatim}
0x00040000
0x00100000
0x00900000
\end{verbatim}

Therefore the memory wrapper should translate:

\begin{verbatim}
static bool ExtMemLoader_AddressToOffset(uint32_t address,
                                         uint32_t size,
                                         uint32_t *offset)
{
    const uint32_t base = STM32EXTLOADER_DEVICE_ADDR;
    const uint32_t flash_size = STM32EXTLOADER_DEVICE_SIZE;

    if (offset == NULL)
    {
        return false;
    }

    if (address < base)
    {
        return false;
    }

    uint32_t local = address - base;

    if (local >= flash_size)
    {
        return false;
    }

    if (size > (flash_size - local))
    {
        return false;
    }

    *offset = local;
    return true;
}
\end{verbatim}

This is preferred over masking the address.

\section{Important erase/update rules}\label{wiki-reference-flash-layout-proposal--11-important-eraseupdate-rules}

\begin{enumerate}
\def\labelenumi{\arabic{enumi}.}
\tightlist
\item
  Never mass-erase the flash during a normal firmware update.
\item
  Never erase the guard sector at \texttt{0x00000000} as part of an application update.
\item
  Do not erase factory-information sectors during normal operation.
\item
  Update device configuration using A/B sectors.
\item
  For future FoE updates, erase/program only the inactive application slot.
\item
  Verify an application image before marking it pending.
\item
  Keep the currently confirmed application until the new application has successfully booted and confirmed itself.
\item
  Treat the FSBL as separately updateable infrastructure.
\item
  Any FSBL update procedure needs its own power-loss/recovery strategy because this design intentionally does not use FSBL1 as a redundant FSBL copy.
\end{enumerate}

\section{Summary}\label{wiki-reference-flash-layout-proposal--12-summary}

\begin{verbatim}
0x00000000  Guard / invalid FSBL1             4 KiB
0x00001000  Parameter storage                 124 KiB
0x00020000  Reserved transition region        128 KiB
0x00040000  FSBL2                             512 KiB
0x000C0000  Reserved                          256 KiB
0x00100000  Application A                     8 MiB
0x00900000  Application B                     8 MiB
0x01100000  Future/data/models/logs           47 MiB
0x04000000  End of flash
\end{verbatim}

Memory-mapped through XSPI1:

\begin{verbatim}
0x90000000 – 0x93FFFFFF
\end{verbatim}

\end{document}
